% Sistemas Geodésicos de Referência — Parte 2: os referenciais celestes (ICRS e ICRF) e a ponte
% até o chão. "A Terra gira em relação a quê?" Continua sgr_terrestres.tex, cujo último slide
% promete a Parte 2 ("Os referenciais celestes: os que não giram com a Terra"). Quatro aulas de
% 50 minutos, uma por parte. Tudo começa em linguagem de ensino médio, em passos pequenos, com
% apêndice técnico no fim. Mesmo formato, mesma paleta e mesmas convenções de quase_geoide.tex.
%
% Bilíngue (AGENTS.md): o mesmo arquivo gera as duas versões.
%   Português:  latexmk -pdf sgr_celestes.tex      -> sgr_celestes.pdf
%   English:    latexmk -pdf sgr_celestes_en.tex   -> sgr_celestes_en.pdf
%   Limpar auxiliares:  latexmk -c sgr_celestes.tex sgr_celestes_en.tex
% Todo texto visível passa por \tr{português}{english}. Números decimais passam por \num{}, que
% troca vírgula por ponto na versão em inglês. Números derivados e coordenadas das figuras:
% sgr_celestes_numeros.py (rodar sem argumentos, ou com --figuras).
%
% Acompanham a aula os simuladores astrometria/ (estrelas, AR e Dec, tempo de Greenwich),
% GNSS_simul/ (órbitas: ECI -> ECEF com o ângulo sideral) e sistemas_coordenadas/.
%
% Fontes principais (valores e datas conferidos em 18/10/2026; ver os slides "Fontes e leituras"):
%   IERS Conventions (2010), IERS Technical Note 36; resoluções da IAU (1997 B2; 2000 B1.3 e
%   B1.6-B1.9; 2006 B1; 2009 B3; 2018 B2; 2021 B3); Ma et al. 1998, AJ 116, 516 (ICRF); Fey et al.
%   2015, AJ 150, 58 (ICRF2); Charlot et al. 2020, A&A 644, A159 (ICRF3); Gaia Collaboration,
%   Klioner et al. 2022, A&A 667, A148 (Gaia-CRF3); ESA 1997, SP-1200 (Hipparcos); Capitaine &
%   Wallace 2006, A&A 450, 855; Meeus 1998 (Astronomical Algorithms); Stephenson et al. 2016, Proc.
%   R. Soc. A 472, 20160404; Agnew 2024, Nature 628, 333; Ashby 2003, Living Rev. Relativ. 6, 1;
%   Takamoto et al. 2020, Nat. Photon. 14, 411; Vallado et al. 2006, AIAA 2006-6753 (SGP4, TEME);
%   Navipedia (tempo no GNSS); BIPM (Brochura do SI, 9a ed.; CGPM 2022 Res. 4; CGPM 2026 Res. 5,
%   cujo texto oficial NAO foi conferido: dados da declaracao do presidente do CCTF e da imprensa);
%   IVS Biennial Report 2017+2018 (estacao de Fortaleza); Observatorio Nacional (Hora Legal,
%   Decreto 4.264/2002); IBGE (RBMC, 153 estacoes ativas em dez/2025).

\newif\ifEN
\ifdefined\aulaEN \ENtrue \fi

\documentclass[aspectratio=169,11pt]{beamer}

\usetheme[progressbar=frametitle, sectionpage=progressbar, numbering=fraction]{metropolis}
\metroset{block=fill}

\usepackage[T1]{fontenc}
\usepackage{lmodern}
\ifEN
  \usepackage[brazilian,english]{babel}
\else
  \usepackage[english,brazilian]{babel}
\fi
\usepackage{amsmath, amssymb, mathtools}
\usepackage{booktabs}
\usepackage{siunitx}
\usepackage[most]{tcolorbox}
\usepackage{xspace}
\usepackage{tikz}
\usetikzlibrary{arrows.meta, calc, positioning, decorations.pathreplacing,
                decorations.pathmorphing, patterns, babel}

% ------------------------------------------------------------------ idioma / language
\newcommand{\tr}[2]{\ifEN#2\else#1\fi}
\ifEN
  \sisetup{output-decimal-marker={.}, group-separator={\,}, group-minimum-digits=4,
           group-digits=integer}
\else
  \sisetup{output-decimal-marker={,}, group-separator={\,}, group-minimum-digits=4,
           group-digits=integer}
\fi

% ------------------------------------------------------------------ paleta (a mesma das outras aulas)
\definecolor{mStone}{HTML}{1C1917}
\definecolor{mStoneMid}{HTML}{78716C}
\definecolor{mStoneLight}{HTML}{F5F5F4}
\definecolor{mTeal}{HTML}{0D9488}
\definecolor{mRose}{HTML}{E11D48}
\definecolor{mAmber}{HTML}{D97706}
\definecolor{mIndigo}{HTML}{4F46E5}
\definecolor{mSky}{HTML}{0284C7}

% As ideias que se repetem em toda figura desta aula (as mesmas das aulas anteriores)
\colorlet{corSist}{mIndigo}    % o SISTEMA — a regra escrita: ICRS, modelos e convenções da IAU
\colorlet{corFrame}{mTeal}     % o FRAME / o MEDIDO — fontes do ICRF, VLBI, EOP, o chão (ITRF)
\colorlet{corTempo}{mAmber}    % o TEMPO — época, giro da Terra, relógios, a eclíptica (o ano)
\colorlet{corErro}{mRose}      % a DIFERENÇA — erro, correção, alerta
\colorlet{corCeu}{mSky}        % o CÉU — esfera celeste, equador celeste, estrelas

\setbeamercolor{normal text}{fg=mStone, bg=white}
\setbeamercolor{palette primary}{fg=white, bg=mStone}
\setbeamercolor{frametitle}{fg=white, bg=mStone}
\setbeamercolor{alerted text}{fg=mRose}
\setbeamercolor{example text}{fg=mTeal}
\setbeamercolor{progress bar}{fg=mTeal, bg=mTeal!20}
\setbeamercolor{progress bar in head/foot}{fg=mTeal, bg=mTeal!20}
\setbeamercolor{progress bar in section page}{fg=mTeal, bg=mTeal!20}
\setbeamercolor{title separator}{fg=mTeal}
\setbeamercolor{block title}{fg=white, bg=mTeal}
\setbeamercolor{block body}{fg=mStone, bg=mTeal!8}

% ------------------------------------------------------------------ caixas
\newtcolorbox{ideia}[1][]{enhanced, colback=mTeal!7, colframe=mTeal, boxrule=0.6pt, arc=2pt,
  left=6pt, right=6pt, top=4pt, bottom=4pt, fonttitle=\bfseries\small, coltitle=white, #1}
\newtcolorbox{cuidado}[1][]{enhanced, colback=mAmber!8, colframe=mAmber, boxrule=0.6pt, arc=2pt,
  left=6pt, right=6pt, top=4pt, bottom=4pt, fonttitle=\bfseries\small, coltitle=white, #1}
\newtcolorbox{armadilha}[1][]{enhanced, colback=mRose!6, colframe=mRose, boxrule=0.6pt, arc=2pt,
  left=6pt, right=6pt, top=4pt, bottom=4pt, fonttitle=\bfseries\small, coltitle=white, #1}
\newtcolorbox{regra}[1][]{enhanced, colback=mIndigo!7, colframe=mIndigo, boxrule=0.6pt, arc=2pt,
  left=6pt, right=6pt, top=4pt, bottom=4pt, fonttitle=\bfseries\small, coltitle=white, #1}

% ------------------------------------------------------------------ estilos TikZ
\tikzset{
  seta/.style={-{Stealth[length=2.4mm]}, thick},
  seta2/.style={{Stealth[length=2mm]}-{Stealth[length=2mm]}},
  rotulo/.style={font=\scriptsize, text=mStone},
  caixa/.style={draw, rounded corners=3pt, thick, align=center, font=\scriptsize},
  pessoa/.pic={
    \fill (0,0.40) circle (0.075);
    \draw[line width=0.9pt] (0,0.33) -- (0,0.13);
    \draw[line width=0.9pt] (0,0.13) -- (-0.08,0) (0,0.13) -- (0.08,0);
    \draw[line width=0.9pt] (-0.11,0.26) -- (0.11,0.26);
  },
  antena/.pic={
    \draw[thick] (0,-0.02) -- (0,0.14);
    \draw[thick, fill=white] (-0.2,0.34) .. controls (-0.1,0.10) and (0.1,0.10) .. (0.2,0.34) -- cycle;
    \draw[thick] (0,0.2) -- (0,0.42);
    \fill (0,0.43) circle (0.02);
  },
  satelite/.pic={
    \draw[thick, fill=white] (-0.07,-0.05) rectangle (0.07,0.05);
    \draw[thick, fill=corFrame!30] (-0.30,-0.035) rectangle (-0.09,0.035);
    \draw[thick, fill=corFrame!30] (0.09,-0.035) rectangle (0.30,0.035);
  },
}

\newcommand{\bnum}[1]{\textbf{\textcolor{mTeal}{#1}}}
\newcommand{\sist}[1]{\textcolor{corSist}{\textbf{#1}}}
\newcommand{\fram}[1]{\textcolor{corFrame}{\textbf{#1}}}
\newcommand{\tmpo}[1]{\textcolor{corTempo}{\textbf{#1}}}
\newcommand{\asec}{\ensuremath{{}^{\prime\prime}}}
\newcommand{\amin}{\ensuremath{{}^{\prime}}}
\newcommand{\gr}{\ensuremath{{}^{\circ}}}
\newcommand{\mas}{\ensuremath{\,\mathrm{mas}}\xspace}
\newcommand{\uas}{\ensuremath{\,\mu\mathrm{as}}\xspace}
\newcommand{\hms}[3]{\ensuremath{#1^{\mathrm{h}}\,#2^{\mathrm{m}}\,}\num{#3}\textsuperscript{s}}
\newcommand{\dms}[3]{\ensuremath{#1\gr\,#2\amin\,}\num{#3}\ensuremath{\asec}}

% ---- esfera celeste em projeção ortográfica (AR0 = 40°, elevação 22°, raio 1) ----
% Geometria gerada por sgr_celestes_numeros.py (funções esfera, partes): equador, eclíptica,
% arcos de AR e Dec da estrela de exemplo (AR 75°, Dec +30°) e pontos notáveis.
\newcommand{\esferaContorno}{\draw[mStoneMid!70, thick, fill=corCeu!5] (0,0) circle (1);}
\newcommand{\esferaEquador}{%
  \draw[corCeu!70, dashed] plot[smooth] coordinates {(0.996,-0.033) (1.000,0.000) (0.996,0.033) (0.985,0.065) (0.966,0.097) (0.940,0.128) (0.906,0.158) (0.866,0.187) (0.819,0.215) (0.766,0.241) (0.707,0.265) (0.643,0.287) (0.574,0.307) (0.500,0.324) (0.423,0.340) (0.342,0.352) (0.259,0.362) (0.174,0.369) (0.087,0.373) (0.000,0.375) (-0.087,0.373) (-0.174,0.369) (-0.259,0.362) (-0.342,0.352) (-0.423,0.340) (-0.500,0.324) (-0.574,0.307) (-0.643,0.287) (-0.707,0.265) (-0.766,0.241) (-0.819,0.215) (-0.866,0.187) (-0.906,0.158) (-0.940,0.128) (-0.966,0.097) (-0.985,0.065) (-0.996,0.033) (-1.000,0.000)};
  \draw[corCeu, very thick] plot[smooth] coordinates {(-0.643,-0.287) (-0.574,-0.307) (-0.500,-0.324) (-0.423,-0.340) (-0.342,-0.352) (-0.259,-0.362) (-0.174,-0.369) (-0.087,-0.373) (0.000,-0.375) (0.087,-0.373) (0.174,-0.369) (0.259,-0.362) (0.342,-0.352) (0.423,-0.340) (0.500,-0.324) (0.574,-0.307) (0.643,-0.287) (0.707,-0.265) (0.766,-0.241) (0.819,-0.215) (0.866,-0.187) (0.906,-0.158) (0.940,-0.128) (0.966,-0.097) (0.985,-0.065) (0.996,-0.033)};
  \draw[corCeu, very thick] plot[smooth] coordinates {(-1.000,0.000) (-0.996,-0.033) (-0.985,-0.065) (-0.966,-0.097) (-0.940,-0.128) (-0.906,-0.158) (-0.866,-0.187) (-0.819,-0.215) (-0.766,-0.241) (-0.707,-0.265) (-0.643,-0.287)};
}
\newcommand{\esferaEcliptica}{%
  \draw[corTempo!70, dashed] plot[smooth] coordinates {(0.952,0.298) (0.951,0.307) (0.944,0.315) (0.930,0.320) (0.908,0.322) (0.880,0.323) (0.844,0.320) (0.803,0.315) (0.755,0.308) (0.702,0.299) (0.643,0.287) (0.579,0.273) (0.511,0.257) (0.439,0.239) (0.364,0.219) (0.286,0.198) (0.205,0.175) (0.123,0.150) (0.041,0.125) (-0.042,0.098) (-0.125,0.071) (-0.207,0.043) (-0.287,0.015) (-0.365,-0.013) (-0.441,-0.041) (-0.513,-0.069) (-0.581,-0.096) (-0.644,-0.122) (-0.703,-0.148) (-0.756,-0.172) (-0.804,-0.195) (-0.845,-0.217) (-0.880,-0.237) (-0.909,-0.255) (-0.930,-0.272) (-0.944,-0.286) (-0.952,-0.298)};
  \draw[corTempo, very thick] plot[smooth] coordinates {(-0.643,-0.287) (-0.579,-0.273) (-0.511,-0.257) (-0.439,-0.239) (-0.364,-0.219) (-0.286,-0.198) (-0.205,-0.175) (-0.123,-0.150) (-0.041,-0.125) (0.042,-0.098) (0.125,-0.071) (0.207,-0.043) (0.287,-0.015) (0.365,0.013) (0.441,0.041) (0.513,0.069) (0.581,0.096) (0.644,0.122) (0.703,0.148) (0.756,0.172) (0.804,0.195) (0.845,0.217) (0.880,0.237) (0.909,0.255) (0.930,0.272) (0.944,0.286) (0.952,0.298)};
  \draw[corTempo, very thick] plot[smooth] coordinates {(-0.952,-0.298) (-0.951,-0.307) (-0.944,-0.315) (-0.930,-0.320) (-0.908,-0.322) (-0.880,-0.323) (-0.844,-0.320) (-0.803,-0.315) (-0.755,-0.308) (-0.702,-0.299) (-0.643,-0.287)};
}
\newcommand{\esferaEixo}{\draw[corSist, dashed] (0,0) -- (0.000,-0.927); \draw[corSist, very thick] (0,0) -- (0.000,0.927) -- (0.000,1.094);}
\newcommand{\esferaArcoAR}[1]{\draw[#1] plot[smooth] coordinates {(-0.643,-0.287) (-0.574,-0.307) (-0.500,-0.324) (-0.423,-0.340) (-0.342,-0.352) (-0.259,-0.362) (-0.174,-0.369) (-0.087,-0.373) (0.000,-0.375) (0.087,-0.373) (0.174,-0.369) (0.259,-0.362) (0.342,-0.352) (0.423,-0.340) (0.500,-0.324) (0.574,-0.307)};}
\newcommand{\esferaArcoDec}[1]{\draw[#1] plot[smooth] coordinates {(0.574,-0.307) (0.573,-0.258) (0.570,-0.208) (0.567,-0.158) (0.561,-0.107) (0.554,-0.056) (0.546,-0.005) (0.535,0.046) (0.524,0.097) (0.511,0.148) (0.497,0.198)};}
\newcommand{\esferaMeridiano}[1]{\draw[#1] plot[smooth] coordinates {(0.497,0.198) (0.464,0.297) (0.426,0.392) (0.384,0.484) (0.337,0.570) (0.287,0.650) (0.233,0.722) (0.177,0.787) (0.119,0.843) (0.060,0.890) (0.000,0.927)};}
\newcommand{\ptGama}{(-0.643,-0.287)}
\newcommand{\ptPoloN}{(0.000,0.927)}
\newcommand{\ptEstrela}{(0.497,0.198)}
\newcommand{\ptPe}{(0.574,-0.307)}
\newcommand{\ptSolJun}{(0.703,0.148)}
\newcommand{\ptSolDez}{(-0.703,-0.148)}
\newcommand{\ptEqSet}{(0.643,0.287)}
\newcommand{\ptSol}{(0.207,-0.043)}
\newcommand{\ptEclA}{(-0.483,-0.250)}
\newcommand{\ptEclB}{(-0.140,-0.155)}
\newcommand{\ptEqA}{(0.866,-0.187)}

% ---- trilha do satélite (lon, lat) em graus: 2 voltas por dia sideral, passo de 10 min ----
% Gerada por sgr_celestes_numeros.py (função trilha_gps, inclinação 55°).
\newcommand{\trilhaUm}[1]{\draw[#1] plot[smooth] coordinates {(0.0,0.0) (0.4,4.1) (0.8,8.2) (1.2,12.3) (1.8,16.3) (2.5,20.3) (3.3,24.2) (4.4,28.1) (5.7,31.9) (7.4,35.5) (9.4,39.0) (11.9,42.2) (14.9,45.3) (18.5,48.0) (22.8,50.4) (27.7,52.4) (33.2,53.8) (39.1,54.7) (45.3,55.0) (51.5,54.7) (57.4,53.7) (62.8,52.2) (67.7,50.2) (71.9,47.8) (75.4,45.0) (78.4,41.9) (80.8,38.6) (82.8,35.1) (84.4,31.5) (85.7,27.7) (86.8,23.9) (87.6,19.9) (88.3,15.9) (88.8,11.9) (89.3,7.8) (89.7,3.7) (90.0,-0.4) (90.4,-4.5) (90.8,-8.6) (91.3,-12.7) (91.9,-16.7) (92.6,-20.7) (93.4,-24.6) (94.5,-28.5) (95.9,-32.2) (97.6,-35.8) (99.6,-39.3) (102.2,-42.6) (105.2,-45.6) (108.9,-48.3) (113.2,-50.6) (118.2,-52.5) (123.7,-53.9) (129.7,-54.8) (135.9,-55.0) (142.1,-54.6) (148.0,-53.6) (153.3,-52.0) (158.1,-50.0) (162.2,-47.5) (165.7,-44.7) (168.6,-41.6) (171.0,-38.3) (173.0,-34.8) (174.6,-31.1) (175.8,-27.3) (176.8,-23.5) (177.7,-19.5) (178.3,-15.5) (178.9,-11.5) (179.3,-7.4) (179.7,-3.3)};}
\newcommand{\trilhaDois}[1]{\draw[#1] plot[smooth] coordinates {(-179.9,0.8) (-179.6,4.9) (-179.1,9.0) (-178.7,13.1) (-178.1,17.1) (-177.4,21.1) (-176.5,25.0) (-175.4,28.8) (-174.0,32.6) (-172.3,36.2) (-170.1,39.6) (-167.5,42.9) (-164.4,45.8) (-160.7,48.5) (-156.3,50.8) (-151.3,52.7) (-145.7,54.0) (-139.7,54.8) (-133.5,55.0) (-127.3,54.5) (-121.5,53.5) (-116.2,51.8) (-111.5,49.8) (-107.4,47.3) (-104.0,44.4) (-101.1,41.3) (-98.8,38.0) (-96.9,34.4) (-95.3,30.8) (-94.1,27.0) (-93.1,23.1) (-92.3,19.1) (-91.6,15.1) (-91.1,11.1) (-90.7,7.0) (-90.3,2.9) (-89.9,-1.2) (-89.5,-5.3) (-89.1,-9.4) (-88.6,-13.5) (-88.0,-17.5) (-87.3,-21.5) (-86.4,-25.4) (-85.2,-29.2) (-83.8,-32.9) (-82.1,-36.5) (-79.9,-39.9) (-77.3,-43.2) (-74.1,-46.1) (-70.3,-48.8) (-65.8,-51.0) (-60.8,-52.9) (-55.1,-54.2) (-49.1,-54.9) (-42.9,-55.0) (-36.7,-54.4) (-30.9,-53.3) (-25.7,-51.7) (-21.0,-49.5) (-17.0,-47.0) (-13.7,-44.1) (-10.9,-41.0) (-8.6,-37.6) (-6.7,-34.1) (-5.2,-30.4) (-4.0,-26.6) (-3.0,-22.7) (-2.2,-18.7) (-1.6,-14.7) (-1.0,-10.7) (-0.6,-6.6) (-0.2,-2.5) (0.1,1.6)};}

% ---- 303 pontos ilustrativos (distribuição uniforme, projeção de Hammer): sgr_celestes_numeros.py ----
\newcommand{\mapaFontes}[1]{\draw[#1] plot[only marks, mark=*, mark size=0.9pt] coordinates {(-0.24,1.50) (-0.14,1.40) (0.38,1.39) (-0.55,1.40) (0.06,1.31) (0.74,1.39) (-0.48,1.29) (0.43,1.27) (-0.95,1.38) (-0.18,1.21) (0.85,1.28) (-0.83,1.25) (0.26,1.17) (1.20,1.34) (-0.50,1.16) (0.76,1.17) (-1.19,1.24) (-0.02,1.10) (1.22,1.21) (-0.87,1.12) (0.53,1.08) (-1.51,1.27) (-0.38,1.04) (1.08,1.10) (-1.25,1.11) (0.21,1.00) (1.56,1.17) (-0.79,1.01) (0.83,1.00) (-1.62,1.13) (-0.18,0.95) (1.41,1.05) (-1.20,1.00) (0.48,0.93) (1.88,1.14) (-0.61,0.91) (1.15,0.94) (-1.61,1.00) (0.07,0.87) (1.74,1.01) (-1.06,0.89) (0.79,0.86) (-1.98,1.03) (-0.38,0.82) (1.47,0.89) (-1.50,0.89) (0.36,0.80) (2.05,0.98) (-0.86,0.80) (1.11,0.80) (-1.92,0.90) (-0.11,0.75) (1.80,0.85) (-1.33,0.78) (0.68,0.73) (-2.30,0.94) (-0.60,0.71) (1.44,0.76) (-1.79,0.79) (0.20,0.68) (2.11,0.82) (-1.10,0.69) (1.01,0.68) (-2.22,0.81) (-0.31,0.64) (1.77,0.71) (-1.60,0.69) (0.53,0.62) (2.42,0.79) (-0.83,0.61) (1.35,0.63) (-2.07,0.69) (0.01,0.57) (2.10,0.68) (-1.36,0.60) (0.87,0.56) (-2.50,0.71) (-0.53,0.54) (1.69,0.58) (-1.86,0.59) (0.35,0.51) (2.42,0.64) (-1.07,0.52) (1.22,0.51) (-2.33,0.60) (-0.20,0.47) (2.03,0.54) (-1.61,0.50) (0.70,0.46) (2.72,0.61) (-0.76,0.44) (1.57,0.47) (-2.11,0.49) (0.14,0.41) (2.35,0.50) (-1.32,0.42) (1.05,0.41) (-2.58,0.50) (-0.43,0.38) (1.91,0.42) (-1.85,0.41) (0.50,0.36) (2.67,0.46) (-1.00,0.35) (1.41,0.35) (-2.36,0.40) (-0.07,0.32) (2.24,0.37) (-1.56,0.33) (0.86,0.30) (-2.82,0.40) (-0.66,0.29) (1.76,0.30) (-2.09,0.31) (0.29,0.26) (2.56,0.32) (-1.23,0.26) (1.22,0.25) (-2.58,0.30) (-0.30,0.23) (2.10,0.25) (-1.79,0.23) (0.65,0.21) (2.86,0.27) (-0.89,0.19) (1.58,0.20) (-2.31,0.21) (0.06,0.17) (2.42,0.20) (-1.46,0.16) (1.02,0.15) (-2.79,0.19) (-0.53,0.13) (1.92,0.14) (-2.01,0.14) (0.43,0.11) (2.73,0.14) (-1.12,0.10) (1.38,0.10) (-2.52,0.10) (-0.16,0.08) (2.25,0.08) (-1.69,0.07) (0.80,0.06) (-2.98,0.07) (-0.76,0.04) (1.73,0.04) (-2.22,0.03) (0.21,0.02) (2.56,0.02) (-1.34,0.01) (1.16,0.00) (-2.71,-0.01) (-0.39,-0.01) (2.06,-0.02) (-1.90,-0.03) (0.58,-0.04) (2.85,-0.06) (-0.99,-0.05) (1.51,-0.06) (-2.41,-0.08) (-0.02,-0.07) (2.37,-0.09) (-1.56,-0.09) (0.94,-0.09) (-2.87,-0.13) (-0.62,-0.11) (1.85,-0.12) (-2.09,-0.14) (0.35,-0.13) (2.66,-0.17) (-1.20,-0.15) (1.29,-0.15) (-2.57,-0.20) (-0.25,-0.16) (2.16,-0.20) (-1.75,-0.19) (0.71,-0.19) (2.91,-0.26) (-0.84,-0.20) (1.62,-0.22) (-2.26,-0.25) (0.12,-0.22) (2.44,-0.28) (-1.40,-0.25) (1.06,-0.25) (-2.71,-0.32) (-0.47,-0.26) (1.93,-0.29) (-1.92,-0.30) (0.48,-0.28) (2.69,-0.37) (-1.04,-0.30) (1.39,-0.32) (-2.40,-0.37) (-0.11,-0.31) (2.21,-0.38) (-1.58,-0.35) (0.82,-0.34) (-2.80,-0.47) (-0.68,-0.35) (1.70,-0.39) (-2.07,-0.42) (0.25,-0.37) (2.46,-0.47) (-1.23,-0.40) (1.15,-0.41) (-2.49,-0.51) (-0.32,-0.41) (1.97,-0.47) (-1.73,-0.47) (0.59,-0.43) (2.65,-0.58) (-0.87,-0.45) (1.45,-0.48) (-2.18,-0.55) (0.03,-0.47) (2.21,-0.57) (-1.38,-0.51) (0.91,-0.50) (-2.54,-0.65) (-0.52,-0.51) (1.72,-0.57) (-1.85,-0.59) (0.36,-0.53) (2.39,-0.68) (-1.04,-0.56) (1.21,-0.58) (-2.24,-0.69) (-0.18,-0.57) (1.95,-0.67) (-1.51,-0.63) (0.68,-0.60) (2.52,-0.81) (-0.69,-0.62) (1.46,-0.66) (-1.91,-0.72) (0.15,-0.63) (2.12,-0.78) (-1.17,-0.68) (0.96,-0.68) (-2.23,-0.84) (-0.36,-0.67) (1.68,-0.76) (-1.58,-0.76) (0.45,-0.70) (2.22,-0.90) (-0.83,-0.73) (1.20,-0.76) (-1.92,-0.87) (-0.05,-0.74) (1.83,-0.87) (-1.25,-0.81) (0.72,-0.78) (-2.15,-1.00) (-0.51,-0.79) (1.39,-0.86) (-1.60,-0.90) (0.24,-0.81) (1.91,-1.00) (-0.93,-0.86) (0.94,-0.87) (-1.84,-1.02) (-0.21,-0.86) (1.52,-0.97) (-1.27,-0.95) (0.49,-0.90) (1.90,-1.13) (-0.62,-0.92) (1.10,-0.97) (-1.52,-1.05) (0.06,-0.94) (1.56,-1.10) (-0.95,-1.00) (0.68,-0.99) (-1.65,-1.19) (-0.33,-0.99) (1.19,-1.09) (-1.19,-1.10) (0.28,-1.03) (1.49,-1.23) (-0.64,-1.07) (0.80,-1.10) (-1.31,-1.22) (-0.08,-1.08) (1.16,-1.21) (-0.86,-1.17) (0.42,-1.13) (-1.28,-1.35) (-0.36,-1.16) (0.80,-1.22) (-0.94,-1.28) (0.10,-1.20) (0.97,-1.35) (-0.52,-1.26) (0.44,-1.27) (-0.83,-1.40) (-0.12,-1.29) (0.58,-1.37) (-0.48,-1.38) (0.14,-1.36) (0.42,-1.48) (-0.10,-1.44)};}


\title{\tr{A Terra gira em relação a quê?}{The Earth spins relative to what?}}
\subtitle{\tr{Sistemas Geodésicos de Referência --- Parte 2: os referenciais celestes (ICRS e
ICRF) e a ponte até o chão --- quatro aulas}{Geodetic Reference Systems --- Part 2: the
celestial reference systems (ICRS and ICRF) and the bridge to the ground --- four classes}}
\author{}
\institute{\tr{Simuladores de Astrometria e de Sistemas de Coordenadas}{Astrometry and
Coordinate Systems Simulators}}
\date{}

\begin{document}

\maketitle

% ==================================================================
\begin{frame}{\tr{O que vamos ver}{What we will see}}
  \setbeamertemplate{section in toc}[sections numbered]
  \footnotesize
  \begin{columns}[T]
    \begin{column}{0.48\textwidth}
      \tableofcontents[sections={1-9}]
    \end{column}
    \begin{column}{0.48\textwidth}
      \tableofcontents[sections={10-18}]
    \end{column}
  \end{columns}
\end{frame}

\begin{frame}{\tr{Esta é a Parte 2}{This is Part 2}}
  \begin{columns}[T]
    \begin{column}{0.50\textwidth}
      \begin{regra}[title={\tr{Parte 1 (a aula anterior)}{Part 1 (the previous class)}}]
        \small\tr{Os referenciais \textbf{terrestres}: ITRS e ITRF, SIRGAS2000, o WGS 84. Os que
        \emph{giram junto com a Terra}.}{The \textbf{terrestrial} reference systems: ITRS and
        ITRF, SIRGAS2000, WGS 84. The ones that \emph{turn with the Earth}.}
      \end{regra}
    \end{column}
    \begin{column}{0.46\textwidth}
      \begin{ideia}[title={\tr{Parte 2 (esta aula)}{Part 2 (this class)}}]
        \small\tr{Os referenciais \textbf{celestes}: ICRS e ICRF. Os que \bnum{não giram com a
        Terra} --- e a \bnum{ponte} que liga o céu ao chão.}{The \textbf{celestial} reference
        systems: ICRS and ICRF. The ones that \bnum{do not turn with the Earth} --- and the
        \bnum{bridge} that links the sky to the ground.}
      \end{ideia}
    \end{column}
  \end{columns}
  \bigskip
  \centering
  \tr{A Parte 2 também \bnum{começa do zero}: se você já sabe o que é ascensão reta, ela passa
  rápido --- mas as ideias dela voltam todas na aula 3.}{Part 2 also \bnum{starts from
  scratch}: if you already know what right ascension is, it goes quickly --- but its ideas all
  come back in class 3.}
\end{frame}

\begin{frame}{\tr{Quatro aulas}{Four classes}}
  \begin{columns}[T]
    \begin{column}{0.47\textwidth}
      \begin{regra}[title={\tr{Aula 1 · Um céu que não gira}{Class 1 · A sky that does not
        spin}}]
        \small\tr{Por que um referencial que não gira com a Terra; coordenadas de uma estrela; o
        que, no céu, \emph{se mexe}.}{Why a reference system that does not turn with the Earth;
        coordinates of a star; what, in the sky, \emph{moves}.}
      \end{regra}
      \vspace{2pt}
      \begin{regra}[title={\tr{Aula 3 · A ponte céu--chão}{Class 3 · The sky--ground
        bridge}}]
        \small\tr{Como a Terra está virada no espaço: rotação, precessão, nutação, polo --- e os
        EOP.}{How the Earth is turned in space: rotation, precession, nutation, pole --- and the
        EOP.}
      \end{regra}
    \end{column}
    \begin{column}{0.47\textwidth}
      \begin{regra}[title={\tr{Aula 2 · Os faróis do Universo}{Class 2 · The beacons of the
        Universe}}]
        \small\tr{Quasares, VLBI, o ICRS e o ICRF; depois o Hipparcos e o Gaia.}{Quasars, VLBI,
        the ICRS and the ICRF; then Hipparcos and Gaia.}
      \end{regra}
      \vspace{2pt}
      \begin{regra}[title={\tr{Aula 4 · O tempo e a prática}{Class 4 · Time and practice}}]
        \small\tr{UT1, UTC, TAI, o segundo intercalar e o tempo do GNSS; da órbita ao receptor; o
        Brasil.}{UT1, UTC, TAI, the leap second and GNSS time; from the orbit to the receiver;
        Brazil.}
      \end{regra}
    \end{column}
  \end{columns}
  \vspace{2pt}
  \centering\footnotesize
  \tr{Cada aula tem 50 minutos, começa relembrando a anterior e termina com um resumo.
  Vamos \bnum{devagar}: uma ideia por slide.}{Each class is 50 minutes long, starts by
  recalling the previous one and ends with a summary. We will go \bnum{slowly}: one idea per
  slide.}
\end{frame}

% ==================================================================
%  AULA 1 · UM CÉU QUE NÃO GIRA
% ==================================================================
\begin{frame}[standout]
  \tr{Parte 1}{Part 1}\\[6pt]
  {\normalsize \tr{Um céu que não gira}{A sky that does not spin}}\\[4pt]
  {\scriptsize \tr{aula 1 de 4}{class 1 of 4}}
\end{frame}

\section{\tr{Onde paramos}{Where we left off}}

\begin{frame}{\tr{Onde paramos: a Parte 1}{Where we left off: Part 1}}
  \begin{columns}[T]
    \begin{column}{0.54\textwidth}
      \tr{Na Parte 1, sobre os referenciais \bnum{terrestres}:}{In Part 1, on the
      \bnum{terrestrial} reference systems:}
      \begin{enumerate}
        \setlength{\itemsep}{4pt}
        \item \tr{uma coordenada precisa dizer em qual \bnum{frame} e em qual
              \bnum{época};}{a coordinate must say in which \bnum{frame} and at which
              \bnum{epoch};}
        \item \tr{o \sist{sistema} é a regra escrita; o \fram{frame}, as estações de
              verdade;}{the \sist{system} is the written rule; the \fram{frame}, the real
              stations;}
        \item \tr{o ITRS/ITRF fica \bnum{preso à Terra}: gira junto com ela;}{the ITRS/ITRF is
              \bnum{fixed to the Earth}: it turns with it;}
        \item \tr{a orientação dele vem de um acordo de 1984 --- e quem a amarra são os
              \bnum{quasares}, vistos pelo VLBI.}{its orientation comes from a 1984 agreement
              --- and what ties it down are \bnum{quasars}, seen by VLBI.}
      \end{enumerate}
    \end{column}
    \begin{column}{0.42\textwidth}
      \begin{cuidado}[title={\tr{Ficou faltando uma peça}{One piece is missing}}]
        \tr{Tudo isso serve a quem está \alert{parado na Terra}. E quem não está --- um
        satélite, a Lua, uma estrela, um quasar?}{All of that serves those who are \alert{at
        rest on the Earth}. What about those who are not --- a satellite, the Moon, a star, a
        quasar?}
        \medskip

        \tr{Esses precisam de um referencial que \bnum{não gira com a Terra}.}{They need a
        reference system that \bnum{does not turn with the Earth}.}
      \end{cuidado}
    \end{column}
  \end{columns}
\end{frame}

\begin{frame}{\tr{As perguntas de hoje}{Today's questions}}
  \begin{columns}[T]
    \begin{column}{0.58\textwidth}
      \begin{enumerate}
        \setlength{\itemsep}{8pt}
        \item \tr{Por que precisamos de um referencial que não gira com a Terra?}{Why do we need
              a reference system that does not turn with the Earth?}
        \item \tr{Como se dão coordenadas a uma estrela?}{How does a star get coordinates?}
        \item \tr{O que quer dizer ``J2000.0'' no catálogo?}{What does ``J2000.0'' mean in the
              catalog?}
        \item \tr{O céu parece parado --- mas o que, de fato, \emph{se mexe}?}{The sky looks
              still --- but what, in fact, \emph{moves}?}
      \end{enumerate}
    \end{column}
    \begin{column}{0.38\textwidth}
      \begin{ideia}[title={\tr{Combinado}{Deal}}]
        \tr{No fim da aula, você responde às quatro --- com as suas palavras.}{By the end of the
        class you will answer all four --- in your own words.}
        \medskip

        \tr{Elas voltam no último slide.}{They come back on the last slide.}
      \end{ideia}
    \end{column}
  \end{columns}
\end{frame}

% ==================================================================
\section{\tr{Duas surpresas}{Two surprises}}

\begin{frame}{\tr{O mesmo satélite, dois desenhos}{The same satellite, two drawings}}
  \small
  \tr{Um satélite GPS dá uma volta em torno da Terra em cerca de \bnum{12 horas}. O mesmo
  satélite, durante um dia, em dois desenhos:}{A GPS satellite goes once around the Earth in
  about \bnum{12 hours}. The same satellite, over one day, in two drawings:}
  \begin{columns}[T]
    \begin{column}{0.34\textwidth}
      \centering
      \begin{tikzpicture}[scale=0.56]
        \fill[corFrame!15] (0,0) circle (0.55);
        \draw[corFrame, thick] (0,0) circle (0.55);
        \draw[mStoneMid, -{Stealth[length=1.6mm]}] (0.12,0.33) arc (70:150:0.36);
        \draw[corSist, thick] (0,0) circle (2.15);
        \foreach \ho/\ag in {0/0, 3/90, 6/180, 9/270} {
          \fill[corTempo] ({2.15*cos(\ag)},{2.15*sin(\ag)}) circle (0.1);
          \node[rotulo, fill=white, inner sep=1pt] at ({2.68*cos(\ag)},{2.68*sin(\ag)}) {\ho\,h};
        }
        \draw[corSist, seta] (2.15,0.15) arc (4:40:2.15);
        \node[rotulo, text=corFrame] at (0,0) {\tr{Terra}{Earth}};
      \end{tikzpicture}
      \par\smallskip
      {\scriptsize\textcolor{corSist}{\textbf{\tr{visto do espaço: um círculo}{seen from space: a circle}}}}
    \end{column}
    \begin{column}{0.62\textwidth}
      \centering
      \begin{tikzpicture}[x=0.0158cm, y=0.0158cm]
        \draw[mStoneMid!60] (-180,-90) rectangle (180,90);
        \foreach \xx in {-120,-60,0,60,120} { \draw[mStoneMid!25] (\xx,-90) -- (\xx,90); }
        \foreach \yy in {-60,-30,0,30,60} { \draw[mStoneMid!25] (-180,\yy) -- (180,\yy); }
        \foreach \xx in {-180,-90,0,90,180} { \node[rotulo, below] at (\xx,-90) {\xx\gr}; }
        \foreach \yy in {-60,0,60} { \node[rotulo, left] at (-180,\yy) {\yy\gr}; }
        \trilhaDois{corFrame!70, thick, dashed}
        \trilhaUm{corFrame, very thick}
        \foreach \xx/\yy/\ho/\pos in {0/0/0/below right, 45.3/55.0/3/above, 90/-0.4/6/below right,
                                      135.9/-55.0/9/below} {
          \fill[corTempo] (\xx,\yy) circle (4.5);
          \node[rotulo, \pos, inner sep=2pt] at (\xx,\yy) {\ho\,h};
        }
        \node[rotulo, text=corFrame, align=left, right] at (-172,80) {\tr{1ª volta (0 a 12 h)}{1st orbit (0 to 12 h)}};
        \node[rotulo, text=corFrame!70, align=left, right] at (-172,66) {\tr{2ª volta (12 a 24 h)}{2nd orbit (12 to 24 h)}};
      \end{tikzpicture}
      \par
      {\scriptsize\textcolor{corFrame}{\textbf{\tr{visto de um mapa preso à Terra: um laço retorcido}{seen from a map fixed to the Earth: a twisted loop}}}}
    \end{column}
  \end{columns}
  \smallskip
  \begin{ideia}
    \scriptsize\tr{O satélite é o mesmo. Mudou o \bnum{referencial}: à direita, o mapa
    \alert{gira} embaixo dele (uma volta a cada 23 h 56 min).}{The satellite is the same. The
    \bnum{reference system} changed: on the right, the map \alert{turns} underneath it (one
    turn every 23 h 56 min).}
  \end{ideia}
\end{frame}

\begin{frame}{\tr{O norte do céu muda de endereço}{The north of the sky changes address}}
  \begin{columns}[T]
    \begin{column}{0.52\textwidth}
      \small
      \tr{Hoje o polo norte celeste aponta quase para \bnum{Polaris}, a Estrela Polar: fica a
      menos de $1\gr$ dela.}{Today the north celestial pole points almost at \bnum{Polaris}, the
      North Star: it lies less than $1\gr$ away.}
      \medskip

      \tr{Mas o eixo da Terra não aponta sempre para o mesmo lugar. Há uns \num{4800} anos, no
      tempo das pirâmides, a ``estrela do norte'' era \bnum{Thuban}. Daqui a uns \num{11700}
      anos será \bnum{Vega}.}{But the Earth's axis does not always point to the same place.
      About \num{4800} years ago, in the time of the pyramids, the ``north star'' was
      \bnum{Thuban}. In about \num{11700} years it will be \bnum{Vega}.}
      \medskip

      \tr{No Hemisfério Sul não há uma estrela brilhante no polo: é o Cruzeiro do Sul que serve
      de guia.}{In the Southern Hemisphere there is no bright star at the pole: the Southern
      Cross serves as the guide.}
    \end{column}
    \begin{column}{0.44\textwidth}
      \centering
      \begin{tikzpicture}[scale=0.98]
        \fill[mStone] (0,0) circle (0.035);
        \node[rotulo, text=mStoneMid, align=center] at (0,-0.4) {\tr{polo norte da}{north pole of the}\\\tr{eclíptica}{ecliptic}};
        \draw[corSist, thick] (0,0) circle (1.9);
        \draw[corSist, seta] (150:1.9) arc (150:105:1.9);
        \draw[corSist, seta] (75:1.9) arc (75:12:1.9);
        \draw[corSist, seta] (-6:1.9) arc (-6:-60:1.9);
        \fill[corCeu] (88.57:1.938) circle (0.065);
        \draw[corErro, thick] (90:1.9) circle (0.15);
        \node[rotulo, above] at (90:2.12) {Polaris, \tr{hoje}{today}};
        \fill[corCeu] (157.46:1.917) circle (0.065);
        \node[rotulo, left, align=right] at (157.46:2.0) {Thuban\\{\scriptsize\tr{$\approx$ 2800 a.C.}{$\approx$ 2800 BC}}};
        \fill[corCeu] (285.32:2.293) circle (0.065);
        \node[rotulo, below right, align=left] at (285.32:2.35) {Vega\\{\scriptsize\tr{$\approx$ ano \num{13700}}{$\approx$ year \num{13700}}}};
        \node[rotulo, text=corSist, align=center] at (-0.1,0.95) {\tr{uma volta:}{one lap:}\\$\approx$ \num{25800} \tr{anos}{years}};
      \end{tikzpicture}
      \par\smallskip
      {\scriptsize\textcolor{mStoneMid}{\tr{O polo celeste anda num círculo de $23{,}4\gr$ de
      raio.}{The celestial pole walks along a circle of radius $23.4\gr$.}}}
    \end{column}
  \end{columns}
\end{frame}

\begin{frame}{\tr{A tese da Parte 2}{The thesis of Part 2}}
  \begin{columns}[T]
    \begin{column}{0.54\textwidth}
      \tr{Duas surpresas, uma lição: o que está ``parado'' \bnum{depende de quem olha}.}{Two
      surprises, one lesson: what is ``still'' \bnum{depends on who is looking}.}
      \medskip

      \tr{Quem precisa de um referencial que \alert{não gira com a Terra}?}{Who needs a reference
      system that \alert{does not turn with the Earth}?}
      \small
      \begin{itemize}
        \setlength{\itemsep}{1pt}
        \item \tr{\textbf{satélites} (GNSS, laser, DORIS): a órbita obedece à gravidade no
              espaço, não ao mapa;}{\textbf{satellites} (GNSS, laser, DORIS): the orbit obeys
              gravity in space, not the map;}
        \item \tr{\textbf{sondas} e missões à Lua e a Marte;}{\textbf{probes} and missions to
              the Moon and Mars;}
        \item \tr{a \textbf{astronomia}: estrelas, planetas, galáxias;}{\textbf{astronomy}:
              stars, planets, galaxies;}
        \item \tr{a \textbf{própria Terra}: para medir o quanto ela gira, é preciso algo que
              não gira com ela.}{the \textbf{Earth itself}: to measure how much it turns, you
              need something that does not turn with it.}
      \end{itemize}
    \end{column}
    \begin{column}{0.42\textwidth}
      \begin{regra}[title={\tr{A pergunta que organiza as quatro aulas}{The question that
        organizes the four classes}}]
        \centering\large \tr{Girar \textbf{em relação a quê}?}{Spinning \textbf{relative to
        what}?}
      \end{regra}
      \medskip
      \begin{ideia}[title={\tr{Aonde vamos chegar}{Where we are headed}}]
        \small\tr{O sistema \sist{ICRS} e o frame \fram{ICRF}: um céu preso a \bnum{quasares},
        tão distantes que parecem imóveis.}{The \sist{ICRS} system and the \fram{ICRF} frame: a
        sky tied to \bnum{quasars}, so distant that they look motionless.}
      \end{ideia}
    \end{column}
  \end{columns}
\end{frame}

% ==================================================================
\section{\tr{Por que um céu que não gira?}{Why a sky that does not spin?}}

\begin{frame}{\tr{O carrossel, de novo}{The carousel, again}}
  \begin{columns}[T]
    \begin{column}{0.44\textwidth}
      \small
      \tr{Você está num carrossel e joga uma bola para um amigo do outro lado.}{You are on a
      carousel and throw a ball to a friend on the other side.}
      \medskip

      \tr{\bnum{Visto de fora}, a bola vai em linha reta --- mas o amigo girou e ``saiu do
      caminho''.}{\bnum{Seen from outside}, the ball goes in a straight line --- but your friend
      turned and ``got out of the way''.}
      \medskip

      \tr{\bnum{Visto do carrossel}, a bola faz uma curva, como se alguém a empurrasse.}{\bnum{Seen
      from the carousel}, the ball makes a curve, as if someone pushed it.}
      \medskip

      {\small\tr{Lembra do ``empurrãozinho para fora'' da aula de gravimetria? Ele só aparece
      porque a gente está num referencial que gira.}{Remember the ``little push outward'' from
      the gravimetry class? It only appears because we are in a reference system that
      turns.}}
    \end{column}
    \begin{column}{0.54\textwidth}
      \centering
      \begin{tikzpicture}[scale=0.86]
        \begin{scope}
          \fill[mStoneLight] (0,0) circle (1.8);
          \draw[mStoneMid, thick] (0,0) circle (1.8);
          \draw[mStoneMid, -{Stealth[length=1.8mm]}] (1.45,-0.9) arc (-32:22:1.72);
          \fill[mStone] (-1.2,0) circle (0.07);
          \node[rotulo, left] at (-1.26,0) {A};
          \draw[mStoneMid] (1.2,0) circle (0.07);
          \fill[corTempo] ({1.2*cos(50)},{1.2*sin(50)}) circle (0.08);
          \node[rotulo, above right] at ({1.2*cos(50)},{1.2*sin(50)}) {B};
          \draw[corSist, seta, line width=1.2pt] (-1.15,0) -- (1.12,0);
          \node[rotulo, text=corSist, align=center] at (0,-2.45) {\tr{visto de fora:}{seen from outside:}\\\tr{reta}{straight}};
        \end{scope}
        \begin{scope}[xshift=4.3cm]
          \fill[mStoneLight] (0,0) circle (1.8);
          \draw[mStoneMid, thick] (0,0) circle (1.8);
          \draw[mStoneMid, -{Stealth[length=1.8mm]}] (1.45,-0.9) arc (-32:22:1.72);
          \fill[mStone] (-1.2,0) circle (0.07);
          \node[rotulo, left] at (-1.26,0) {A};
          \fill[corTempo] (1.2,0) circle (0.08);
          \node[rotulo, above right] at (1.2,0) {B};
          \draw[corErro, -{Stealth[length=2.2mm]}, line width=1.2pt, domain=0:1, samples=40, variable=\t]
            plot ({(-1.2+2.4*\t)*cos(50*\t)}, {-(-1.2+2.4*\t)*sin(50*\t)});
          \node[rotulo, text=corErro, align=center] at (0,-2.45) {\tr{visto do carrossel:}{seen from the carousel:}\\\tr{curva}{curved}};
        \end{scope}
      \end{tikzpicture}
    \end{column}
  \end{columns}
\end{frame}

\begin{frame}{\tr{A Terra deixa pistas, mas não diz ``em relação a quê''}{The Earth leaves clues, but does not say ``relative to what''}}
  \begin{columns}[T]
    \begin{column}{0.54\textwidth}
      \small
      \tr{Num carrossel fechado você ainda \emph{sente} que está girando. Com a Terra também:
      em 1851, em Paris, Foucault pendurou um pêndulo de \bnum{67 metros} no Panteão --- e o
      plano de oscilação foi \bnum{girando} sozinho.}{On a closed carousel you can still
      \emph{feel} that you are turning. The Earth is similar: in 1851, in Paris, Foucault hung
      a \bnum{67-meter} pendulum in the Panthéon --- and its plane of swing \bnum{turned} by
      itself.}
      \medskip

      \tr{Em \textbf{Curitiba}: uma volta a cada \bnum{2,3 dias} (uns $155\gr$ por dia), no
      sentido anti-horário visto de cima.}{In \textbf{Curitiba}: one turn every \bnum{2.3 days}
      (about $155\gr$ per day), counterclockwise seen from above.}
      \medskip

      \begin{cuidado}[title={\tr{Mas falta uma resposta}{But an answer is missing}}]
        \footnotesize\tr{O pêndulo prova que a Terra gira, mas não diz \alert{em relação a
        quê}. Newton imaginou um ``espaço absoluto''; \bnum{na prática, usamos os objetos mais
        distantes que conseguimos ver}.}{The pendulum proves that the Earth turns, but does not
        say \alert{relative to what}. Newton imagined an ``absolute space''; \bnum{in practice,
        we use the most distant objects we can see}.}
      \end{cuidado}
    \end{column}
    \begin{column}{0.42\textwidth}
      \centering
      \begin{tikzpicture}[scale=0.95]
        \draw[mStoneMid!40] (0,0) circle (1.95);
        \foreach \ag/\cor/\rot in {0/corSist/{\tr{agora}{now}}, 38.8/corTempo/{$+6$\,h}, 77.5/corErro/{$+12$\,h}} {
          \draw[\cor, very thick] ({-1.7*cos(\ag)},{-1.7*sin(\ag)}) -- ({1.7*cos(\ag)},{1.7*sin(\ag)});
          \node[rotulo, text=\cor] at ({2.3*cos(\ag)},{2.3*sin(\ag)}) {\rot};
        }
        \fill (0,0) circle (0.05);
        \draw[mStoneMid, seta] (1.35,0) arc (0:70:1.35);
        \node[rotulo, text=mStoneMid, align=center] at (0,-2.55) {\tr{visto de cima, em Curitiba}{seen from above, in Curitiba}};
      \end{tikzpicture}
    \end{column}
  \end{columns}
\end{frame}

% ==================================================================
\section{\tr{O céu como uma esfera}{The sky as a sphere}}

\begin{frame}{\tr{A esfera celeste}{The celestial sphere}}
  \begin{columns}[T]
    \begin{column}{0.50\textwidth}
      \small
      \tr{As estrelas estão a distâncias muito diferentes. Mas, para dizer \emph{onde} uma
      estrela está no céu, só interessa a \bnum{direção}.}{Stars are at very different
      distances. But to say \emph{where} a star is in the sky, only the \bnum{direction}
      matters.}
      \medskip

      \tr{Então imaginamos uma esfera gigante em volta da Terra, como o domo de um planetário,
      e ``pintamos'' cada astro no ponto onde o raio de visão o cruza: é a \bnum{esfera
      celeste}.}{So we imagine a giant sphere around the Earth, like the dome of a planetarium,
      and ``paint'' each object at the point where the line of sight crosses it: the
      \bnum{celestial sphere}.}
      \medskip

      \begin{ideia}[title={\tr{Direção, não posição}{Direction, not position}}]
        \footnotesize\tr{Na Parte 1, uma coordenada era uma posição, em metros. Aqui, é só uma
        direção: \bnum{dois ângulos}.}{In Part 1, a coordinate was a position, in meters. Here,
        it is only a direction: \bnum{two angles}.}
      \end{ideia}
    \end{column}
    \begin{column}{0.46\textwidth}
      \centering
      \begin{tikzpicture}[scale=2.05]
        \esferaContorno
        \esferaEquador
        \fill[corFrame!40] (0,0) circle (0.045);
        \draw[corFrame] (0,0) circle (0.045);
        \node[rotulo, text=corFrame, below right] at (0.03,-0.03) {\tr{Terra}{Earth}};
        \foreach \dx/\dy/\fa/\fb in {0.55/0.45/1.30/1.55, -0.52/0.40/1.30/1.55} {
          \draw[mStoneMid, dashed] (0,0) -- ({\fb*\dx},{\fb*\dy});
          \fill[corCeu] (\dx,\dy) circle (0.03);
          \fill[corTempo] ({\fa*\dx},{\fa*\dy}) circle (0.035);
          \fill[corTempo] ({\fb*\dx},{\fb*\dy}) circle (0.035);
        }
        \node[rotulo, text=corCeu, below right, align=left] at (0.55,0.45) {\tr{o ponto}{the point}\\\tr{na esfera}{on the sphere}};
        \node[rotulo, text=corTempo, align=center] at (0,-1.22) {\tr{estrelas a distâncias diferentes,}{stars at different distances,}\\\tr{na mesma direção}{in the same direction}};
      \end{tikzpicture}
    \end{column}
  \end{columns}
\end{frame}

\begin{frame}{\tr{Latitude e longitude do céu}{Latitude and longitude of the sky}}
  \begin{columns}[T]
    \begin{column}{0.50\textwidth}
      \small
      \tr{Na Terra, damos a latitude e a longitude. No céu, é a mesma ideia:}{On the Earth, we
      give latitude and longitude. In the sky, it is the same idea:}
      \smallskip
      \begin{itemize}
        \setlength{\itemsep}{2pt}
        \item \tr{o \bnum{equador celeste} é o equador da Terra ``esticado'' até a esfera; os
              \bnum{polos celestes} são o eixo da Terra esticado;}{the \bnum{celestial equator}
              is the Earth's equator ``stretched'' out to the sphere; the \bnum{celestial
              poles} are the Earth's axis stretched out;}
        \item \textcolor{corErro}{\textbf{\tr{declinação $\delta$}{declination $\delta$}}}:
              \tr{a ``latitude'' --- de $-90\gr$ a $+90\gr$;}{the ``latitude'' --- from
              $-90\gr$ to $+90\gr$;}
        \item \textcolor{corTempo}{\textbf{\tr{ascensão reta $\alpha$}{right ascension
              $\alpha$}}}: \tr{a ``longitude'' --- medida para leste, a partir de um ponto
              de partida $\gamma$.}{the ``longitude'' --- measured eastward, from a starting
              point $\gamma$.}
      \end{itemize}
      \smallskip
      \tr{Em vez de graus, a ascensão reta costuma vir em \bnum{horas}.}{Instead of degrees,
      right ascension usually comes in \bnum{hours}.}
    \end{column}
    \begin{column}{0.46\textwidth}
      \centering
      \begin{tikzpicture}[scale=2.05]
        \esferaContorno
        \esferaEquador
        \esferaEixo
        \esferaMeridiano{mStoneMid, dashed}
        \esferaArcoAR{corTempo, line width=2.2pt}
        \esferaArcoDec{corErro, line width=2.2pt}
        \fill[corFrame] \ptGama circle (0.03);
        \node[rotulo, text=corFrame, below left] at \ptGama {$\gamma$};
        \fill[corCeu] \ptEstrela circle (0.04);
        \node[rotulo, above right] at \ptEstrela {\tr{estrela}{star}};
        \node[rotulo, text=corSist, above left] at (0,1.0) {\tr{polo norte celeste}{north celestial pole}};
        \node[rotulo, text=corTempo, below] at (0.05,-0.44) {$\alpha$};
        \node[rotulo, text=corErro, right] at (0.585,-0.05) {$\delta$};
        \node[rotulo, text=corCeu, left] at (-0.78,0.05) {\tr{equador}{equator}};
      \end{tikzpicture}
    \end{column}
  \end{columns}
\end{frame}

\begin{frame}{\tr{Horas viram graus}{Hours become degrees}}
  \begin{columns}[T]
    \begin{column}{0.46\textwidth}
      \tr{O céu dá uma volta em 24 horas. Então:}{The sky makes one turn in 24 hours. So:}
      \medskip
      \centering
      \renewcommand{\arraystretch}{1.35}
      \begin{tabular}{rcl}
        \toprule
        24 h & $=$ & $360\gr$ \\
        1 h & $=$ & $15\gr$ \\
        1 min & $=$ & $15\amin$ \\
        1 s & $=$ & $15\asec$ \\
        \bottomrule
      \end{tabular}
      \par\medskip
      \raggedright
      {\small\tr{(minuto e segundo \emph{de tempo} viram minuto e segundo \emph{de arco}, vezes
      15.)}{(minutes and seconds \emph{of time} become minutes and seconds \emph{of arc},
      times 15.)}}
    \end{column}
    \begin{column}{0.52\textwidth}
      \centering
      \small
      \renewcommand{\arraystretch}{1.4}
      \begin{tabular}{lcc}
        \toprule
        & \textcolor{corTempo}{$\alpha$} & \textcolor{corErro}{$\delta$} \\
        \midrule
        \textbf{Sirius} & 6 h 45 min \num{8.9} s & $-16\gr42\amin58\asec$ \\
        & $=\num{101.287}\gr$ & $=\num{-16.716}\gr$ \\
        \midrule
        \textbf{Acrux} & 12 h 26 min \num{35.9} s & $-63\gr05\amin57\asec$ \\
        & $=\num{186.650}\gr$ & $=\num{-63.099}\gr$ \\
        \bottomrule
      \end{tabular}
      \par\medskip
      \raggedright
      {\footnotesize\tr{Sirius é a estrela mais brilhante do céu noturno; Acrux é a estrela do
      ``pé'' do \bnum{Cruzeiro do Sul}. Valores para o equinócio e a época J2000.0.}{Sirius is the
      brightest star in the night sky; Acrux is the ``foot'' star of the \bnum{Southern
      Cross}. Values for the equinox and epoch J2000.0.}}
    \end{column}
  \end{columns}
\end{frame}

\begin{frame}{\tr{O ponto de partida: eclíptica e ponto vernal}{The starting point: ecliptic and vernal point}}
  \begin{columns}[T]
    \begin{column}{0.52\textwidth}
      \small
      \tr{Para a longitude da Terra, o mundo combinou o meridiano de Greenwich. E para a
      ascensão reta?}{For longitude on the Earth, the world agreed on the Greenwich meridian.
      And for right ascension?}
      \medskip

      \tr{Ao longo do ano, o Sol anda entre as estrelas por uma linha: a \tmpo{eclíptica}. Ela
      é inclinada \bnum{$23{,}4\gr$} em relação ao equador celeste --- a mesma inclinação do
      eixo da Terra, a causa das estações.}{Over the year, the Sun walks among the stars along
      a line: the \tmpo{ecliptic}. It is tilted by \bnum{$23.4\gr$} relative to the celestial
      equator --- the same tilt as the Earth's axis, the cause of the seasons.}
      \medskip

      \tr{As duas linhas se cruzam em dois pontos, os \bnum{equinócios}. O de março é o
      \textbf{ponto vernal} $\gamma$: onde o Sol cruza o equador, de sul para norte. É o
      zero da ascensão reta.}{The two lines cross at two points, the \bnum{equinoxes}. The
      March one is the \textbf{vernal point} $\gamma$: where the Sun crosses the equator,
      from south to north. It is the zero of right ascension.}
    \end{column}
    \begin{column}{0.44\textwidth}
      \centering
      \begin{tikzpicture}[scale=1.95]
        \esferaContorno
        \esferaEquador
        \esferaEcliptica
        \draw[corTempo, seta] \ptEclA -- \ptEclB;
        \fill[corTempo] \ptSol circle (0.045);
        \node[rotulo, text=corTempo, above left] at \ptSol {\tr{Sol}{Sun}};
        \fill[corFrame] \ptGama circle (0.035);
        \node[rotulo, text=corFrame, below left, align=right] at (-0.68,-0.74) {$\gamma$ \tr{(março)}{(March)}};
        \draw[mStoneMid!60, thin] \ptGama -- (-0.68,-0.74);
        \fill[corSist] \ptSolJun circle (0.03);
        \node[rotulo, text=mStoneMid, above right, align=left] at (0.80,0.60) {\tr{solstício}{solstice}\\\tr{de junho}{in June}};
        \draw[mStoneMid!60, thin] \ptSolJun -- (0.80,0.60);
        \node[rotulo, text=corCeu, below] at (0.1,-0.40) {\tr{equador celeste}{celestial equator}};
        \node[rotulo, text=corTempo, above] at (0.1,0.43) {\tr{eclíptica}{ecliptic}};
      \end{tikzpicture}
      \par
      {\scriptsize\textcolor{mStoneMid}{\tr{O ângulo entre as duas linhas é $23{,}4\gr$.}{The
      angle between the two lines is $23.4\gr$.}}}
    \end{column}
  \end{columns}
\end{frame}

\begin{frame}{\tr{Sua vez}{Your turn}}
  \begin{columns}[T]
    \begin{column}{0.50\textwidth}
      \begin{regra}[title={\tr{Três continhas}{Three little sums}}]
        \small
        \begin{enumerate}
          \setlength{\itemsep}{4pt}
          \item[(a)] \tr{Vega tem ascensão reta de $18$ h $36$ min $57$ s. Quantos graus?}{Vega has
                right ascension $18$ h $36$ min $57$ s. How many degrees?}
          \item[(b)] \tr{Uma estrela tem $\alpha = 150\gr$ e $\delta = -30\gr$. Qual é a ascensão
                reta em horas? Ela está no céu norte ou sul?}{A star has $\alpha = 150\gr$ and
                $\delta = -30\gr$. What is the right ascension in hours? Is it in the northern
                or southern sky?}
          \item[(c)] \tr{Quantos graus a Terra gira em \textbf{4 minutos}?}{How many degrees
                does the Earth turn in \textbf{4 minutes}?}
        \end{enumerate}
      \end{regra}
    \end{column}
    \begin{column}{0.46\textwidth}
      \onslide<2->{
      \begin{ideia}[title={\tr{Respostas}{Answers}}]
        \small\raggedright
        (a) $\num{18.6158}\ \text{h} \times 15 = \bnum{\num{279.24}\gr}$.
        \medskip

        (b) $150 / 15 = \bnum{10}$ h, \tr{céu \bnum{sul} (declinação negativa).}{\bnum{southern}
        sky (negative declination).}
        \medskip

        (c) $4 \times 15\amin = 60\amin = \bnum{1\gr}$. \tr{Guarde: \alert{4 minutos de tempo
        $=$ 1 grau}.}{Remember: \alert{4 minutes of time $=$ 1 degree}.}
      \end{ideia}
      }
    \end{column}
  \end{columns}
\end{frame}

% ==================================================================
\section{\tr{O céu se mexe}{The sky moves}}

\begin{frame}{\tr{O que quer dizer ``J2000.0''}{What ``J2000.0'' means}}
  \begin{columns}[T]
    \begin{column}{0.52\textwidth}
      \small
      \tr{Todo catálogo de estrelas traz, ao lado das coordenadas, uma etiqueta: \bnum{equinócio
      e época J2000.0}.}{Every star catalog carries, next to the coordinates, a label:
      \bnum{equinox and epoch J2000.0}.}
      \medskip

      \tr{J2000.0 é \bnum{1.º de janeiro de 2000, ao meio-dia}.}{J2000.0 is \bnum{1 January
      2000, at noon}.}
      \medskip

      \tr{É a mesma lição da Parte 1: uma coordenada precisa dizer em qual \bnum{frame} e em
      qual \bnum{época}. No céu, a etiqueta diz: ``os pontos de partida --- o equador e o
      equinócio --- são os que existiam naquele instante''.}{It is the same lesson as in Part 1:
      a coordinate must say in which \bnum{frame} and at which \bnum{epoch}. In the sky, the
      label says: ``the starting points --- the equator and the equinox --- are the ones that
      existed at that instant''.}
    \end{column}
    \begin{column}{0.44\textwidth}
      \begin{armadilha}[title={\tr{Por que a etiqueta é necessária}{Why the label is needed}}]
        \footnotesize\tr{Porque os \alert{dois pontos de partida se mexem}, e as
        \alert{próprias estrelas também}.}{Because the \alert{two starting points move}, and the
        \alert{stars themselves too}.}
      \end{armadilha}
      \smallskip
      \begin{ideia}[title={\tr{Três coisas se mexem}{Three things move}}]
        \footnotesize
        \begin{enumerate}
          \setlength{\itemsep}{2pt}
          \item \tr{o equinócio e o equador: \bnum{precessão};}{the equinox and the equator:
                \bnum{precession};}
          \item \tr{o equador de novo, ondulando: \bnum{nutação};}{the equator again,
                rippling: \bnum{nutation};}
          \item \tr{as estrelas: \bnum{movimento próprio}.}{the stars: \bnum{proper
                motion}.}
        \end{enumerate}
      \end{ideia}
    \end{column}
  \end{columns}
\end{frame}

\begin{frame}{\tr{Precessão: o pião da Terra}{Precession: the Earth's spinning top}}
  \begin{columns}[T]
    \begin{column}{0.50\textwidth}
      \tr{Faça um pião girar inclinado. Além de girar sobre o próprio eixo, o eixo dele
      desenha, devagar, um \bnum{cone}.}{Spin a top while it is tilted. Besides spinning around
      its own axis, its axis slowly draws a \bnum{cone}.}
      \medskip

      \tr{A Terra faz o mesmo: o eixo gira em torno do eixo da eclíptica, mantendo a inclinação
      de $23{,}4\gr$. Uma volta completa leva $\approx$ \bnum{\num{25800} anos}.}{The Earth does
      the same: its axis turns around the axis of the ecliptic, keeping the tilt of $23.4\gr$.
      One full turn takes $\approx$ \bnum{\num{25800} years}.}
      \medskip

      \begin{ideia}[title={\tr{O nome}{The name}}]
        \small\tr{Esse movimento lento do eixo é a \bnum{precessão}. Foi o que moveu o polo
        celeste de Thuban a Polaris, e o levará a Vega.}{This slow motion of the axis is
        \bnum{precession}. It is what moved the celestial pole from Thuban to Polaris, and will
        take it to Vega.}
      \end{ideia}
    \end{column}
    \begin{column}{0.46\textwidth}
      \centering
      \begin{tikzpicture}[scale=1.22]
        \draw[mStoneMid, dashed] (0,-1.3) -- (0,3.1);
        \node[rotulo, text=mStoneMid, above] at (0,3.1) {\tr{eixo da eclíptica (fixo)}{ecliptic axis (fixed)}};
        \begin{scope}[rotate=-23.4]
          \fill[corFrame!15] (0,0) ellipse (0.9 and 0.82);
          \draw[corFrame, thick] (0,0) ellipse (0.9 and 0.82);
          \draw[corFrame!70] (-0.9,0) arc (180:360:0.9 and 0.2);
          \draw[corFrame!70, dashed] (-0.9,0) arc (180:0:0.9 and 0.2);
          \draw[mStone, very thick] (0,-1.25) -- (0,1.9);
          \draw[mStoneMid, -{Stealth[length=1.6mm]}] (-0.4,1.4) arc (200:520:0.4 and 0.1);
        \end{scope}
        \draw[corSist, dashed] (0,1.744) ellipse (0.754 and 0.17);
        \draw[corSist, dashed] (0,0) -- (-0.754,1.744) (0,0) -- (0.754,1.744);
        \draw[corSist, seta] (0.66,1.60) arc (-25:-155:0.70 and 0.17);
        \node[rotulo, text=corSist, align=left, left] at (-0.75,2.15) {\tr{o eixo desenha}{the axis draws}\\\tr{um cone}{a cone}};
        \node[rotulo, text=mStone, above right] at (0.72,1.86) {\tr{eixo da Terra}{Earth's axis}};
        \draw[mStone] (0,0.85) arc (90:66.6:0.85);
        \node[rotulo, above right] at (0.10,0.80) {$\num{23.4}\gr$};
      \end{tikzpicture}
    \end{column}
  \end{columns}
\end{frame}

\begin{frame}{\tr{Quem empurra o pião?}{What pushes the top?}}
  \begin{columns}[T]
    \begin{column}{0.50\textwidth}
      \tr{A Terra não é uma esfera perfeita: o raio no Equador é \bnum{21 km maior} que nos
      polos --- um ``barrigão''.}{The Earth is not a perfect sphere: its radius at the Equator
      is \bnum{21 km larger} than at the poles --- a ``belly''.}
      \medskip

      \tr{O Sol e a Lua puxam mais forte o lado do barrigão que está \emph{mais perto}, e menos
      o lado mais longe. A diferença é um torque que tenta \alert{endireitar o eixo}.}{The Sun
      and the Moon pull harder on the side of the belly that is \emph{nearer}, and less on the
      far side. The difference is a torque that tries to \alert{straighten the axis}.}
      \medskip

      \tr{Mas a Terra gira, como um giroscópio. Em vez de tombar, o eixo \bnum{dá a volta}: é a
      precessão.}{But the Earth spins, like a gyroscope. Instead of toppling, the axis
      \bnum{goes around}: that is precession.}
    \end{column}
    \begin{column}{0.46\textwidth}
      \centering
      \begin{tikzpicture}[scale=0.98]
        \fill[corTempo] (-3.0,0) circle (0.3);
        \node[rotulo, below] at (-3.0,-0.38) {\tr{Sol e Lua}{Sun and Moon}};
        \draw[mStoneMid, dashed] (-2.6,0) -- (2.6,0);
        \node[rotulo, text=mStoneMid, above left, align=right] at (2.6,0.02) {\tr{plano da}{plane of}\\\tr{eclíptica}{ecliptic}};
        \begin{scope}[rotate=-23.4]
          \fill[corFrame!15] (0,0) ellipse (0.78 and 0.66);
          \draw[corFrame, thick] (0,0) ellipse (0.78 and 0.66);
          \draw[mStone, thick] (0,-1.1) -- (0,1.2);
          \draw[corFrame, thick] (-1.1,0) -- (1.1,0);
        \end{scope}
        \draw[corErro, -{Stealth[length=2.4mm]}, line width=1.8pt] (-0.95,0.40) -- (-2.45,0.12);
        \draw[corErro, -{Stealth[length=1.8mm]}, line width=0.9pt] (0.95,-0.40) -- (0.35,-0.52);
        \node[rotulo, text=corErro, align=center] at (-1.55,0.95) {\tr{puxão maior}{stronger pull}};
        \node[rotulo, text=corErro, align=center] at (1.45,-0.95) {\tr{puxão menor}{weaker pull}};
        \draw[corSist, -{Stealth[length=2mm]}, thick] (-0.5,-1.0) arc (-150:-60:0.55);
        \node[rotulo, text=corSist, below] at (0,-1.25) {\tr{torque}{torque}};
      \end{tikzpicture}
    \end{column}
  \end{columns}
\end{frame}

\begin{frame}{\tr{O ponto vernal escorrega}{The vernal point slips}}
  \begin{columns}[T]
    \begin{column}{0.48\textwidth}
      \tr{Se o equador celeste muda de lugar, o ponto onde ele cruza a eclíptica --- o ponto
      vernal --- \bnum{escorrega para o oeste}, ao longo da eclíptica:}{If the celestial
      equator changes place, the point where it crosses the ecliptic --- the vernal point ---
      \bnum{slips westward}, along the ecliptic:}
      \[ \boxed{\; \num{50.29}\asec \text{ \tr{por ano}{per year}} \;} \]
      \tr{Isso dá $1{,}4\gr$ por século, ou $1\gr$ a cada \bnum{72 anos}.}{That is $1.4\gr$ per
      century, or $1\gr$ every \bnum{72 years}.}
      \medskip

      \tr{E como a ascensão reta é medida \emph{a partir do ponto vernal}, a ascensão reta de
      toda estrela ``fixa'' \alert{vai mudando}.}{And since right ascension is measured
      \emph{from the vernal point}, the right ascension of every ``fixed'' star
      \alert{keeps changing}.}
    \end{column}
    \begin{column}{0.48\textwidth}
      \centering
      \begin{tikzpicture}[x=0.82cm, y=0.95cm]
        \draw[corTempo, very thick] (-0.9,0) -- (5.5,0);
        \node[rotulo, text=corTempo, above left] at (5.5,0.02) {\tr{eclíptica}{ecliptic}};
        \draw[corCeu, very thick, domain=-0.9:5.5, samples=80, variable=\t] plot ({\t},{0.8*sin(60*\t)});
        \draw[corSist, thick, dashed, domain=-0.9:5.5, samples=80, variable=\t] plot ({\t},{0.8*sin(60*(\t+0.45))});
        \node[rotulo, text=corCeu, above] at (1.5,0.83) {\tr{equador em 2000}{equator in 2000}};
        \node[rotulo, text=corSist, below] at (4.4,-0.84) {\tr{equador hoje}{equator today}};
        \fill[corFrame] (0,0) circle (0.06);
        \node[rotulo, text=corFrame, above right] at (0.02,0.03) {$\gamma$};
        \fill[corSist] (-0.45,0) circle (0.06);
        \node[rotulo, text=corSist, below left] at (-0.45,-0.03) {$\gamma'$};
        \fill[mStone] (2.5,0.5) circle (0.07);
        \node[rotulo, above left] at (2.45,0.55) {\tr{estrela fixa}{fixed star}};
        \draw[corFrame, seta2] (0,-1.25) -- (2.5,-1.25);
        \draw[corSist, seta2] (-0.45,-1.75) -- (2.5,-1.75);
        \node[rotulo, text=corFrame, above] at (1.25,-1.25) {$\alpha_{2000}$};
        \node[rotulo, text=corSist, below] at (1.0,-1.75) {$\alpha_{\text{\tr{hoje}{today}}}$ \tr{(maior)}{(larger)}};
        \draw[mStoneMid!60, dotted] (2.5,0.5) -- (2.5,-1.75);
      \end{tikzpicture}
      \par
      {\scriptsize\textcolor{mStoneMid}{\tr{(deslocamento exagerado)}{(shift exaggerated)}}}
    \end{column}
  \end{columns}
\end{frame}

\begin{frame}{\tr{Pare e pense}{Stop and think}}
  \begin{columns}[T]
    \begin{column}{0.50\textwidth}
      \tr{O ponto vernal escorrega para o \bnum{oeste}. A ascensão reta se mede para o
      \bnum{leste}, a partir dele.}{The vernal point slips \bnum{westward}. Right ascension is
      measured \bnum{eastward}, from it.}
      \medskip
      \begin{regra}[title={\tr{A pergunta}{The question}}]
        \tr{A ascensão reta de uma estrela fixa, hoje, é \textbf{maior} ou \textbf{menor} do que
        em 2000? E por quanto, para Sirius?}{Is the right ascension of a fixed star, today,
        \textbf{larger} or \textbf{smaller} than in 2000? And by how much, for Sirius?}
      \end{regra}
    \end{column}
    \begin{column}{0.46\textwidth}
      \onslide<2->{
      \begin{ideia}[title={\tr{Resposta}{Answer}}]
        \small\raggedright
        \tr{\bnum{Maior}. O ponto de partida andou para trás, então a distância medida a partir
        dele, para a frente, \emph{cresce}.}{\bnum{Larger}. The starting point moved backward,
        so the distance measured from it, forward, \emph{grows}.}
        \medskip

        \tr{Sirius: de $101{,}287\gr$ para $101{,}587\gr$ --- uns \bnum{+$1\,078\asec$}
        ($+0{,}30\gr$) em 26,8 anos.}{Sirius: from $101.287\gr$ to $101.587\gr$ --- about
        \bnum{+$1\,078\asec$} ($+0.30\gr$) in 26.8 years.}
      \end{ideia}
      }
    \end{column}
  \end{columns}
\end{frame}

\begin{frame}{\tr{Hiparco e Spica (cerca de 129 a.C.)}{Hipparchus and Spica (about 129 BC)}}
  \begin{columns}[T]
    \begin{column}{0.54\textwidth}
      \tr{O astrônomo grego \textbf{Hiparco} comparou a posição de \bnum{Spica} com observações
      feitas uns 150 anos antes, por Timocharis.}{The Greek astronomer \textbf{Hipparchus}
      compared the position of \bnum{Spica} with observations made about 150 years earlier by
      Timocharis.}
      \medskip

      \tr{Spica estava uns \bnum{$2\gr$} mais longe do ponto vernal. As estrelas não tinham se
      mexido tanto: quem tinha se mexido era o \alert{ponto de partida}.}{Spica was about
      \bnum{$2\gr$} farther from the vernal point. The stars had not moved that much: what had
      moved was the \alert{starting point}.}
      \medskip

      \tr{Hiparco estimou uns $1\gr$ por século (o valor de hoje é $1{,}4\gr$). Foi a primeira
      vez que alguém percebeu que o céu tem uma \bnum{época}.}{Hipparchus estimated about
      $1\gr$ per century (today's value is $1.4\gr$). It was the first time anyone realized
      that the sky has an \bnum{epoch}.}
    \end{column}
    \begin{column}{0.42\textwidth}
      \begin{ideia}[title={\tr{A conta}{The sum}}]
        \small\tr{$144$ anos $\times$ $50{,}29\asec$ por ano $= \num{7242}\asec \approx
        \bnum{2{,}0\gr}$.}{$144$ years $\times$ $50.29\asec$ per year $= \num{7242}\asec \approx
        \bnum{2.0\gr}$.}
        \medskip

        \tr{Conferem: o deslocamento medido pelos gregos é o que a precessão de hoje prevê.}{It
        checks out: the shift measured by the Greeks is what today's precession predicts.}
      \end{ideia}
    \end{column}
  \end{columns}
\end{frame}

\begin{frame}{\tr{Exemplo: Sirius em J2000.0 e hoje}{Example: Sirius in J2000.0 and today}}
  \begin{columns}[T]
    \begin{column}{0.56\textwidth}
      \centering
      \small
      \renewcommand{\arraystretch}{1.4}
      \begin{tabular}{lrr}
        \toprule
        & \textcolor{corTempo}{$\alpha$} & \textcolor{corErro}{$\delta$} \\
        \midrule
        \textbf{J2000.0} & $\num{101.2872}\gr$ & $\num{-16.7161}\gr$ \\
        \tr{\textbf{hoje} (18/10/2026)}{\textbf{today} (2026-10-18)} & $\num{101.5865}\gr$ &
          $\num{-16.7457}\gr$ \\
        \midrule
        \tr{diferença}{difference} & $+\num{1078}\asec$ & $-\num{106}\asec$ \\
        \bottomrule
      \end{tabular}
      \par\medskip
      {\scriptsize\textcolor{mStoneMid}{\tr{Só precessão (equinócio e equador médios da data).
      O movimento próprio de Sirius acrescenta $\approx -33\asec$ em $\delta$.}{Precession only
      (mean equinox and equator of date). Sirius's proper motion adds $\approx -33\asec$ in
      $\delta$.}}}
    \end{column}
    \begin{column}{0.40\textwidth}
      \begin{ideia}[title={\tr{O tamanho disso}{How big that is}}]
        \small\tr{Total: $\num{0.288}\gr = 17{,}3\amin$.}{Total: $\num{0.288}\gr = 17.3\amin$.}
        \medskip

        \tr{É \bnum{60\,\% do diâmetro da Lua cheia}.}{That is \bnum{60\,\% of the diameter of
        the full Moon}.}
        \medskip

        \tr{Em 27 anos, o ``mesmo'' ponto do catálogo deixou de ser o ponto onde a estrela
        está.}{In 27 years, the ``same'' catalog point stopped being where the star is.}
      \end{ideia}
    \end{column}
  \end{columns}
\end{frame}

\begin{frame}{\tr{Nutação: o pião balança}{Nutation: the top wobbles}}
  \begin{columns}[T]
    \begin{column}{0.46\textwidth}
      \tr{A Lua complica. A órbita dela é inclinada uns $5\gr$ em relação à eclíptica, e o plano
      dessa órbita também gira: uma volta em \bnum{18,6 anos}.}{The Moon complicates things.
      Its orbit is tilted about $5\gr$ relative to the ecliptic, and the plane of that orbit
      also turns: one turn in \bnum{18.6 years}.}
      \medskip

      \tr{Por isso o puxão da Lua no barrigão varia, e o eixo da Terra ``balança'' enquanto
      percorre o círculo: a \bnum{nutação}.}{So the Moon's pull on the belly varies, and the
      Earth's axis ``wobbles'' while it goes around its circle: \bnum{nutation}.}
      \medskip

      \tr{Amplitudes: até $17{,}2\asec$ em longitude e $9{,}2\asec$ na inclinação.}{Amplitudes:
      up to $17.2\asec$ in longitude and $9.2\asec$ in the tilt.}
    \end{column}
    \begin{column}{0.50\textwidth}
      \centering
      \begin{tikzpicture}[x=0.0070cm, y=0.050cm]
        \draw[corTempo, dashed, thick] (-10,0) -- (780,0);
        \draw[corSist, very thick, domain=0:37.2, samples=160, variable=\t]
          plot ({20*\t - 6.84*sin(\t*360/18.6)}, {9.2*cos(\t*360/18.6)});
        \node[rotulo, text=corTempo, above, align=center] at (390,12) {\tr{só precessão (reta)}{precession only (straight)}};
        \node[rotulo, text=corSist, below, align=center] at (390,-12) {\tr{precessão + nutação (ondulado)}{precession + nutation (wavy)}};
        \draw[mStoneMid, seta2] (780,-9.2) -- (780,9.2);
        \node[rotulo, text=mStoneMid, right] at (785,0) {$\pm \num{9.2}\asec$};
        \draw[mStoneMid, seta2] (0,-22) -- (744,-22);
        \node[rotulo, below] at (372,-22) {\tr{37 anos (2 ciclos)}{37 years (2 cycles)}};
      \end{tikzpicture}
      \par
      {\scriptsize\textcolor{mStoneMid}{\tr{Trajetória do polo no céu. Escala vertical ampliada
      $\approx$ 7 vezes.}{Path of the pole in the sky. Vertical scale enlarged $\approx$ 7
      times.}}}
    \end{column}
  \end{columns}
\end{frame}

\begin{frame}{\tr{As estrelas também andam}{Stars move too}}
  \begin{columns}[T]
    \begin{column}{0.52\textwidth}
      \tr{As estrelas viajam pela Galáxia a dezenas de quilômetros por segundo. Para as
      mais próximas, dá para medir o deslocamento aparente: o \bnum{movimento próprio}.}{Stars
      travel through the Galaxy at tens of kilometers per second. For the closest ones, the
      apparent shift can be measured: the \bnum{proper motion}.}
      \medskip
      \begin{itemize}
        \setlength{\itemsep}{3pt}
        \item \tr{\textbf{Estrela de Barnard}: $10{,}4\asec$ por ano --- o recorde. Em
              \bnum{174 anos} anda o diâmetro da Lua.}{\textbf{Barnard's Star}: $10.4\asec$ per
              year --- the record. In \bnum{174 years} it moves the diameter of the Moon.}
        \item \tr{\textbf{Sirius}: $1{,}34\asec$ por ano.}{\textbf{Sirius}: $1.34\asec$ per
              year.}
        \item \tr{A maioria das estrelas: bem menos de $0{,}1\asec$ por ano.}{Most stars: far
              less than $0.1\asec$ per year.}
      \end{itemize}
    \end{column}
    \begin{column}{0.44\textwidth}
      \centering
      \begin{tikzpicture}[scale=0.95]
        \fill[mStoneLight] (0,0) circle (1.4);
        \draw[mStoneMid] (0,0) circle (1.4);
        \node[rotulo, align=center] at (0,0) {\tr{Lua cheia}{full Moon}\\$\num{0.5}\gr$};
        \draw[corErro, seta, very thick] (-1.4,-2.0) -- (1.4,-2.0);
        \node[rotulo, text=corErro, align=center, below] at (0,-2.05) {\tr{Barnard em 174 anos}{Barnard in 174 years}};
        \draw[corSist, seta, very thick] (-1.4,-3.0) -- (-1.4+0.36,-3.0);
        \node[rotulo, text=corSist, right] at (-0.95,-3.0) {\tr{Sirius em 174 anos}{Sirius in 174 years}};
      \end{tikzpicture}
    \end{column}
  \end{columns}
\end{frame}

\begin{frame}{\tr{O que se mexe, em uma tabela}{What moves, in one table}}
  \centering
  \small
  \renewcommand{\arraystretch}{1.5}
  \begin{tabular}{p{4.3cm}p{3.1cm}p{5.0cm}}
    \toprule
    \tr{o que se mexe}{what moves} & \tr{quanto}{how much} & \tr{o que muda nas coordenadas}{what changes in the coordinates} \\
    \midrule
    \tr{o equinócio (precessão)}{the equinox (precession)} & $\num{50.29}\asec$/\tr{ano}{yr} &
      \tr{a ascensão reta de todas as estrelas}{the right ascension of all stars} \\
    \tr{o polo (precessão)}{the pole (precession)} & $20\asec$/\tr{ano}{yr} &
      \tr{a declinação de todas as estrelas}{the declination of all stars} \\
    \tr{o equador (nutação)}{the equator (nutation)} & $\pm 17\asec$ \tr{em}{in} \num{18.6} \tr{anos}{yr} &
      \tr{uma ondulação nas duas}{a ripple in both} \\
    \tr{as estrelas (mov.~próprio)}{the stars (proper motion)} & \tr{até}{up to} $10\asec$/\tr{ano}{yr} &
      \tr{cada estrela, a seu modo}{each star, in its own way} \\
    \bottomrule
  \end{tabular}
  \bigskip

  \tr{Nada do que usamos até aqui como ``parado'' está \alert{de fato parado}. E agora?}{Nothing we
  have used so far as ``still'' is \alert{really still}. Now what?}
\end{frame}

\begin{frame}{\tr{Experimente: o Simulador de Astrometria}{Try it: the Astrometry Simulator}}
  \begin{columns}[T]
    \begin{column}{0.50\textwidth}
      \small
      \tr{Na pasta \texttt{astrometria} você é um agrimensor sem GNSS: aponta o teodolito para
      três estrelas, anota a \emph{distância zenital} e o \emph{ângulo horizontal}, e o
      programa ajusta, por \bnum{mínimos quadrados}, a latitude, a longitude e a constante do
      azimute.}{In the \texttt{astrometria} folder you are a surveyor without GNSS: you point
      the theodolite at three stars, record the \emph{zenith distance} and the
      \emph{horizontal angle}, and the program adjusts, by \bnum{least squares}, the latitude,
      the longitude and the azimuth constant.}
      \medskip
      \begin{itemize}
        \setlength{\itemsep}{2pt}
        \item \tr{catálogo de $\approx$ \num{2800} estrelas (AR e Dec em graus, J2000.0);}{catalog of $\approx$ \num{2800} stars (RA and Dec in degrees, J2000.0);}
        \item \tr{o ``Tempo Médio de Greenwich'' vira a hora sideral; $\text{hora sideral local} = \text{Greenwich} + \text{longitude}$.}{``Greenwich Mean Time'' becomes sidereal time; $\text{local sidereal time} = \text{Greenwich} + \text{longitude}$.}
      \end{itemize}
    \end{column}
    \begin{column}{0.46\textwidth}
      \begin{cuidado}[title={\tr{O que o simulador deixa de fora}{What the simulator leaves out}}]
        \footnotesize
        \begin{itemize}
          \setlength{\itemsep}{2pt}
          \item \tr{\textbf{precessão}: as coordenadas são as de J2000.0 --- em 2026, até $0{,}37\gr$ de diferença ($\approx 40$ km);}{\textbf{precession}: coordinates are those of J2000.0 --- in 2026, up to $0.37\gr$ off ($\approx 40$ km);}
          \item \tr{\textbf{nutação}: até $17\asec$ ($\approx 0{,}5$ km);}{\textbf{nutation}: up to $17\asec$ ($\approx 0.5$ km);}
          \item \tr{\textbf{UT1$-$UTC}: a hora é tratada como UT1; o erro chega a $0{,}9$ s $= 13{,}5\asec$ ($\approx 0{,}4$ km no Equador).}{\textbf{UT1$-$UTC}: the time is treated as UT1; the error reaches $0.9$ s $= 13.5\asec$ ($\approx 0.4$ km at the Equator).}
        \end{itemize}
      \end{cuidado}
      \medskip
      {\small\tr{São simplificações boas para ensinar. Organizar o que fica de fora é o assunto
      desta Parte 2.}{These are good simplifications for teaching. Organizing what is left out
      is the subject of this Part 2.}}
    \end{column}
  \end{columns}
\end{frame}

\begin{frame}{\tr{Para levar para casa --- aula 1}{Take-home --- class 1}}
  \footnotesize
  \begin{enumerate}
    \setlength{\itemsep}{2pt}
    \item \tr{Quem não gira com a Terra --- satélites, sondas, estrelas --- precisa de um
          referencial que \bnum{também não gira}. ``Girar'' é sempre \emph{em relação a
          algo}.}{Whatever does not turn with the Earth --- satellites, probes, stars --- needs a
          reference system that \bnum{does not turn either}. ``Turning'' is always
          \emph{relative to something}.}
    \item \tr{O céu é uma \bnum{esfera} de direções. As coordenadas são dois ângulos:
          \textcolor{corTempo}{\textbf{ascensão reta $\alpha$}} (de $\gamma$, para leste; $24$ h
          $= 360\gr$) e \textcolor{corErro}{\textbf{declinação $\delta$}}.}{The sky is a
          \bnum{sphere} of directions. Coordinates are two angles: \textcolor{corTempo}{\textbf{right
          ascension $\alpha$}} (from $\gamma$, eastward; $24$ h $= 360\gr$) and
          \textcolor{corErro}{\textbf{declination $\delta$}}.}
    \item \tr{$\gamma$, o ponto vernal, é onde a eclíptica cruza o equador celeste. O catálogo
          diz ``\bnum{J2000.0}'': \emph{frame e época}, como na Parte 1.}{$\gamma$, the vernal
          point, is where the ecliptic crosses the celestial equator. The catalog says
          ``\bnum{J2000.0}'': \emph{frame and epoch}, as in Part 1.}
    \item \tr{\bnum{Precessão}: o eixo da Terra dá uma volta em $\approx \num{25800}$ anos; o
          ponto vernal escorrega $50{,}29\asec$/ano. \bnum{Nutação}: $\pm 17\asec$ em 18,6
          anos. \bnum{Movimento próprio}: as estrelas andam.}{\bnum{Precession}: the Earth's
          axis makes one turn in $\approx \num{25800}$ years; the vernal point slips
          $50.29\asec$/yr. \bnum{Nutation}: $\pm 17\asec$ in 18.6 years. \bnum{Proper motion}:
          stars move.}
    \item \tr{Então ``fixo'' não existe nem no equador nem nas estrelas. Precisamos de algo
          \bnum{ainda mais distante e imóvel}.}{So ``fixed'' exists neither in the equator nor in
          the stars. We need something \bnum{even more distant and motionless}.}
  \end{enumerate}
  \begin{ideia}
    \scriptsize \tr{\textbf{Na próxima aula:} os faróis do Universo --- quasares, VLBI e o
    \sist{ICRS} / \fram{ICRF}. \textbf{Experimente:} o Simulador de Astrometria, ajustando a sua
    posição com três estrelas.}{\textbf{Next class:} the beacons of the Universe --- quasars,
    VLBI and the \sist{ICRS} / \fram{ICRF}. \textbf{Try it:} the Astrometry Simulator, adjusting
    your position with three stars.}
  \end{ideia}
\end{frame}


% ==================================================================
%  AULA 2 · OS FARÓIS DO UNIVERSO
% ==================================================================
\begin{frame}[standout]
  \tr{Parte 2}{Part 2}\\[6pt]
  {\normalsize \tr{Os faróis do Universo}{The beacons of the Universe}}\\[4pt]
  {\scriptsize \tr{aula 2 de 4}{class 2 of 4}}
\end{frame}

\section{\tr{Parado de verdade}{Truly still}}

\begin{frame}{\tr{Onde paramos: a aula 1}{Where we left off: class 1}}
  \begin{columns}[T]
    \begin{column}{0.54\textwidth}
      \tr{Na aula 1:}{In class 1:}
      \begin{enumerate}
        \setlength{\itemsep}{4pt}
        \item \tr{o céu é uma \bnum{esfera de direções}, com duas coordenadas: ascensão reta
              $\alpha$ e declinação $\delta$;}{the sky is a \bnum{sphere of directions}, with two
              coordinates: right ascension $\alpha$ and declination $\delta$;}
        \item \tr{elas dependem do equador e do equinócio --- e eles \bnum{se mexem}:
              precessão ($50{,}29\asec$/ano) e nutação ($\pm 17\asec$);}{they depend on the
              equator and the equinox --- and those \bnum{move}: precession
              ($50.29\asec$/yr) and nutation ($\pm 17\asec$);}
        \item \tr{as próprias estrelas andam (até $10\asec$/ano).}{the stars themselves move (up
              to $10\asec$/yr).}
      \end{enumerate}
    \end{column}
    \begin{column}{0.42\textwidth}
      \begin{cuidado}[title={\tr{A pergunta de hoje}{Today's question}}]
        \tr{O que fica parado \alert{de verdade}?}{What stays still \alert{for real}?}
        \medskip

        \tr{Resposta provisória: algo tão longe que, para nós, não se mexe.}{Provisional
        answer: something so far away that, to us, it does not move.}
      \end{cuidado}
    \end{column}
  \end{columns}
\end{frame}

\begin{frame}{\tr{As perguntas de hoje}{Today's questions}}
  \begin{columns}[T]
    \begin{column}{0.58\textwidth}
      \begin{enumerate}
        \setlength{\itemsep}{8pt}
        \item \tr{O que fica parado de verdade?}{What stays still for real?}
        \item \tr{Como medir a direção de algo tão distante, com precisão de milionésimos de
              segundo de arco?}{How do we measure the direction of something so distant, to
              within millionths of an arcsecond?}
        \item \tr{O que são o \sist{ICRS} e o \fram{ICRF}: o sistema e o frame do céu?}{What are
              the \sist{ICRS} and the \fram{ICRF}: the system and the frame of the sky?}
        \item \tr{Onde entram o Hipparcos e o Gaia?}{Where do Hipparcos and Gaia come in?}
      \end{enumerate}
    \end{column}
    \begin{column}{0.38\textwidth}
      \begin{ideia}[title={\tr{Combinado}{Deal}}]
        \tr{No fim da aula, você responde às quatro --- com as suas palavras.}{By the end of the
        class you will answer all four --- in your own words.}
      \end{ideia}
    \end{column}
  \end{columns}
\end{frame}

\begin{frame}{\tr{Quanto mais longe, menos se mexe}{The farther away, the less it moves}}
  \begin{columns}[T]
    \begin{column}{0.50\textwidth}
      \small
      \tr{Numa viagem de carro, as árvores da beira da estrada passam voando. A montanha, lá
      longe, anda devagar. E a Lua parece acompanhar você, quase parada.}{On a car trip, the
      trees by the roadside fly past. The mountain far away moves slowly. And the Moon seems
      to follow you, almost motionless.}
      \medskip

      \tr{Para o mesmo movimento seu, o deslocamento aparente é \bnum{inversamente
      proporcional à distância}: o dobro da distância, metade do deslocamento.}{For the same
      motion of yours, the apparent shift is \bnum{inversely proportional to the distance}:
      twice the distance, half the shift.}
      \medskip

      \begin{ideia}[title={\tr{A ideia}{The idea}}]
        \footnotesize\tr{Quer algo que não se mexe? Procure o que está \bnum{mais longe} que
        você conseguir enxergar.}{Want something that does not move? Look for the thing that is
        \bnum{farthest away} you can see.}
      \end{ideia}
    \end{column}
    \begin{column}{0.46\textwidth}
      \centering
      \begin{tikzpicture}[scale=0.66]
        \draw[mStoneMid, thick] (0,-1.3) -- (0,1.3);
        \fill[corTempo] (0,-0.9) circle (0.08);
        \fill[corTempo] (0,0.9) circle (0.08);
        \node[rotulo, left] at (-0.1,-0.9) {A};
        \node[rotulo, left] at (-0.1,0.9) {B};
        \draw[corTempo, seta] (-0.6,-0.9) -- (-0.6,0.9);
        \node[rotulo, text=corTempo, rotate=90] at (-1.0,0) {\tr{você anda}{you move}};
        \foreach \xx/\nome in {1.7/{\tr{árvore}{tree}}, 3.9/{\tr{montanha}{mountain}}, 7.4/{\tr{Lua}{Moon}}} {
          \pgfmathsetmacro{\phi}{atan(0.9/\xx)}
          \draw[mStoneMid!70] (0,-0.9) -- (\xx,0) -- (0,0.9);
          \draw[corErro, thick] ({\xx-0.55*cos(\phi)},{0.55*sin(\phi)}) arc ({180-\phi}:{180+\phi}:0.55);
          \node[rotulo, below] at (\xx,-0.12) {\nome};
        }
        \fill[corFrame] (1.7,0) circle (0.1);
        \fill[mStoneMid] (3.9,0) circle (0.1);
        \fill[corCeu] (7.4,0) circle (0.1);
        \node[rotulo, text=corErro, above] at (1.45,0.72) {\tr{grande}{large}};
        \node[rotulo, text=corErro, above] at (3.65,0.38) {\tr{médio}{medium}};
        \node[rotulo, text=corErro, above] at (7.15,0.2) {\tr{minúsculo}{tiny}};
      \end{tikzpicture}
      \par\smallskip
      {\scriptsize\textcolor{mStoneMid}{\tr{O ângulo que o objeto ``vê'' (em vermelho) diminui com a
      distância.}{The angle the object ``sees'' (in red) shrinks with distance.}}}
    \end{column}
  \end{columns}
\end{frame}

\begin{frame}{\tr{De quanto é o deslocamento aparente?}{How big is the apparent shift?}}
  \begin{columns}[T]
    \begin{column}{0.38\textwidth}
      \small
      \tr{Algo que se move de lado, com velocidade $v$ a uma distância $d$, parece andar no céu a}{Something moving sideways, at speed $v$ and distance $d$, seems to move across the sky at}
      \[ \boxed{\; \frac{v}{d} \;} \]
      \tr{radianos por segundo.}{radians per second.}
      \bigskip

      \tr{Barnard anda \bnum{150 milhões de vezes mais} que um quasar.}{Barnard moves \bnum{150 million times more} than a quasar.}
    \end{column}
    \begin{column}{0.58\textwidth}
      \centering
      \begin{tikzpicture}[x=0.6cm, y=0.74cm]
        \foreach \yy/\len/\cor/\lab in {
            3.9/9.017/corErro/{\tr{estrela de Barnard (6 anos-luz): 10\,400\,000}{Barnard's Star (6 light-years): 10,400,000}},
            2.8/8.127/corTempo/{\tr{Sirius (8,6 anos-luz): 1\,340\,000}{Sirius (8.6 light-years): 1,340,000}},
            1.7/5.0/corCeu/{\tr{estrelas do outro lado da Galáxia: $\sim$ 1\,000}{stars across the Galaxy: $\sim$ 1,000}},
            0.6/0.845/corFrame/{\tr{quasar a 10 bilhões de anos-luz: 0,07}{quasar 10 billion light-years away: 0.07}}} {
          \fill[\cor!80] (0,\yy) rectangle (\len,{\yy+0.4});
          \node[rotulo, anchor=south west] at (0,{\yy+0.4}) {\lab};
        }
        \draw[mStoneMid] (0,0.3) -- (10,0.3);
        \foreach \xx/\lab in {0/{$10^{-2}$}, 2/{$10^{0}$}, 4/{$10^{2}$}, 6/{$10^{4}$}, 8/{$10^{6}$}, 10/{$10^{8}$}} {
          \draw[mStoneMid] (\xx,0.22) -- (\xx,0.38);
          \node[rotulo, below] at (\xx,0.2) {\lab};
        }
        \node[rotulo, text=mStoneMid] at (5,-0.6) {\tr{movimento aparente ($\mu$as por ano, escala logarítmica)}{apparent motion ($\mu$as per year, logarithmic scale)}};
      \end{tikzpicture}
    \end{column}
  \end{columns}
\end{frame}

\begin{frame}{\tr{O que é um quasar?}{What is a quasar?}}
  \begin{columns}[T]
    \begin{column}{0.58\textwidth}
      \small
      \tr{No centro de algumas galáxias há um \bnum{buraco negro gigante} (milhões a bilhões de
      vezes a massa do Sol) engolindo gás. O gás esquenta ao cair e brilha mais que a galáxia
      inteira; jatos de partículas saem pelos polos.}{At the center of some galaxies sits a
      \bnum{giant black hole} (millions to billions of solar masses) swallowing gas. The gas heats
      up as it falls and outshines the whole galaxy; jets of particles leave through the poles.}
      \smallskip

      \tr{De longe é só um ponto de luz --- ``quase uma estrela'': \emph{quasi-stellar radio
      source}, ou \bnum{quasar}. O primeiro com o espectro decifrado foi o \textbf{3C 273}
      (M.~Schmidt, 1963): sua luz saiu há uns \bnum{2 bilhões de anos}.}{From afar it is just a
      point of light --- ``almost a star'': \emph{quasi-stellar radio source}, or
      \bnum{quasar}. The first with a deciphered spectrum was \textbf{3C 273} (M.~Schmidt,
      1963): its light left about \bnum{2 billion years ago}.}
      \smallskip

      \begin{ideia}[title={\tr{Por que rádio?}{Why radio?}}]
        \footnotesize\tr{Muitos quasares são fortes emissores de \bnum{rádio}; as ondas de rádio
        atravessam as nuvens, valem de dia e de noite --- e, como veremos, dá para juntar antenas
        distantes.}{Many quasars are strong \bnum{radio} emitters; radio waves cross clouds, work
        day and night --- and, as we will see, distant antennas can be combined.}
      \end{ideia}
    \end{column}
    \begin{column}{0.38\textwidth}
      \centering
      \begin{tikzpicture}[scale=1.0]
        \fill[corTempo!35] (0,0) ellipse (1.1 and 0.26);
        \draw[corTempo, thick] (0,0) ellipse (1.1 and 0.26);
        \fill[mStone] (0,0) circle (0.11);
        \fill[corErro!50] (0,0.1) -- (-0.15,1.4) -- (0.15,1.4) -- cycle;
        \fill[corErro!50] (0,-0.1) -- (-0.15,-1.4) -- (0.15,-1.4) -- cycle;
        \fill[corFrame!40] (0,1.6) ellipse (0.45 and 0.25);
        \fill[corFrame!40] (0,-1.6) ellipse (0.45 and 0.25);
        \draw[corFrame] (0,1.6) ellipse (0.45 and 0.25);
        \draw[corFrame] (0,-1.6) ellipse (0.45 and 0.25);
        \node[rotulo, text=corTempo, right] at (1.15,0.1) {\tr{disco de gás}{gas disc}};
        \node[rotulo, text=mStone, below] at (0,-0.12) {\tr{buraco negro}{black hole}};
        \node[rotulo, text=corFrame, right] at (0.5,1.7) {\tr{jato e lóbulo}{jet and lobe}};
        \node[rotulo, text=corFrame, right] at (0.5,-1.7) {\tr{de rádio}{of radio}};
      \end{tikzpicture}
      \par
      {\scriptsize\textcolor{mStoneMid}{\tr{esquema; fora de escala}{sketch; not to scale}}}
    \end{column}
  \end{columns}
\end{frame}

% ==================================================================
\section{\tr{VLBI: ouvir com duas orelhas}{VLBI: listening with two ears}}

\begin{frame}{\tr{Quão fino precisa ser o olhar?}{How sharp must the gaze be?}}
  \begin{columns}[T]
    \begin{column}{0.54\textwidth}
      \small
      \tr{Uma antena só enxerga detalhes maiores que, aproximadamente, o comprimento da onda
      dividido pelo diâmetro:}{A single antenna only sees details larger than roughly the
      wavelength divided by the diameter:}

      \centerline{$\displaystyle \theta \approx \num{1.22}\,\frac{\lambda}{D}$}
      \begin{itemize}
        \setlength{\itemsep}{0pt}
        \item \tr{olho humano: $\approx 60\asec$;}{human eye: $\approx 60\asec$;}
        \item \tr{prato de \textbf{100 m}, ondas de 3,6 cm: $\approx 90\asec$ --- igual ao olho;}{\textbf{100 m} dish, 3.6 cm waves: $\approx 90\asec$ --- same as the eye;}
        \item \tr{duas antenas a \textbf{10\,000 km} uma da outra: $\lambda/B \approx \bnum{0{,}7\mas}$.}{two antennas \textbf{10\,000 km} apart: $\lambda/B \approx \bnum{0.7\mas}$.}
      \end{itemize}
      \smallskip
      \tr{Isso é \bnum{120\,000 vezes mais fino}. E ainda queremos chegar a $30\uas$.}{That is
      \bnum{120\,000 times sharper}. And we still want to reach $30\uas$.}
    \end{column}
    \begin{column}{0.42\textwidth}
      \begin{regra}[title={\tr{A ideia do VLBI}{The idea of VLBI}}]
        \small\tr{\emph{Very Long Baseline Interferometry}: em vez de uma antena gigante,
        \bnum{duas antenas muito distantes} --- um telescópio do tamanho da Terra, feito de
        pedaços.}{\emph{Very Long Baseline Interferometry}: instead of one giant antenna,
        \bnum{two very distant antennas} --- a telescope the size of the Earth, made of
        pieces.}
      \end{regra}
      \smallskip
      \begin{cuidado}[title={\tr{Curiosidade}{Fun fact}}]
        \footnotesize\tr{A mesma técnica, com antenas em vários continentes, deu em 2019 a
        primeira \bnum{imagem de um buraco negro}: o da galáxia M87.}{The same technique, with
        antennas on several continents, gave in 2019 the first \bnum{image of a black hole}: the
        one in galaxy M87.}
      \end{cuidado}
    \end{column}
  \end{columns}
\end{frame}

\begin{frame}{\tr{Duas orelhas e um trovão}{Two ears and a thunderclap}}
  \begin{columns}[T]
    \begin{column}{0.48\textwidth}
      \small
      \tr{Você sabe de que lado vem um trovão: o som chega a uma orelha um pouquinho
      \emph{antes} da outra. O cérebro mede esse \bnum{atraso} e dá a direção.}{You know which
      side a thunderclap comes from: the sound reaches one ear a little \emph{before} the
      other. The brain measures that \bnum{delay} and gives the direction.}
      \medskip

      \tr{Orelhas a $d = 0{,}20$ m, som a $343$ m/s: o atraso máximo é}{Ears $d = 0.20$ m apart,
      sound at $343$ m/s: the maximum delay is}
      \[ \frac{d}{v} = \frac{\num{0.20}}{343} \approx \bnum{\num{0.6}}\ \text{ms} \]
      \tr{E o atraso depende do ângulo $\theta$:}{And the delay depends on the angle $\theta$:}
      \[ \boxed{\; \text{\tr{atraso}{delay}} = \frac{d\,\cos\theta}{v} \;} \]
    \end{column}
    \begin{column}{0.48\textwidth}
      \centering
      \begin{tikzpicture}[scale=1.15]
        \draw[mStoneMid, thick, fill=mStoneLight] (0,0) circle (0.75);
        \draw[mStoneMid, thick, fill=mStoneLight] (-0.85,0) ellipse (0.09 and 0.2);
        \draw[mStoneMid, thick, fill=mStoneLight] (0.85,0) ellipse (0.09 and 0.2);
        \draw[mStoneMid, thick] (-0.12,0.74) -- (0,0.95) -- (0.12,0.74);
        \node[rotulo, below] at (0,-0.85) {\tr{vista de cima}{seen from above}};
        % frentes de onda, direção do som: vem de (cos35, sin35)
        \foreach \px/\py in {0.85/0, -0.85/0} {
          \draw[corCeu, thick] ({\px+1.5*(-sin(35))},{\py+1.5*cos(35)}) -- ({\px-1.5*(-sin(35))},{\py-1.5*cos(35)});
        }
        \draw[corErro, very thick] (-0.85,0) -- ({-0.85+1.7*cos(35)*cos(35)},{1.7*cos(35)*sin(35)});
        \node[rotulo, text=corErro, above left] at (-0.2,0.45) {$d\cos\theta$};
        \draw[mStoneMid, dashed] (-0.85,0) -- (0.85,0);
        \node[rotulo, below] at (0,0) {$d$};
        \draw[mStoneMid] (-0.4,0) arc (0:35:0.45);
        \node[rotulo, text=mStone] at (-0.12,0.13) {$\theta$};
        \draw[corCeu, seta] ({2.3*cos(35)},{2.3*sin(35)}) -- ({1.55*cos(35)},{1.55*sin(35)});
        \node[rotulo, text=corCeu, right] at ({2.3*cos(35)},{2.3*sin(35)}) {\tr{o som}{the sound}};
      \end{tikzpicture}
    \end{column}
  \end{columns}
\end{frame}

\begin{frame}{\tr{O mesmo, com ondas de rádio}{The same, with radio waves}}
  \begin{columns}[T]
    \begin{column}{0.50\textwidth}
      \small
      \tr{Troque as orelhas por \bnum{duas antenas} distantes e o som pelas ondas de rádio de um
      quasar. A onda chega a uma antena antes da outra; medindo o atraso $\tau$, descobre-se
      $\theta$:}{Replace the ears by \bnum{two distant antennas} and the sound by radio waves from
      a quasar. The wave reaches one antenna before the other; measure the delay $\tau$ and
      $\theta$ follows:}
      \[ \boxed{\; \tau = \frac{B\,\cos\theta}{c} \;} \]
      \begin{itemize}
        \setlength{\itemsep}{2pt}
        \item \tr{$B$: a \bnum{linha de base}, a distância entre as antenas;}{$B$: the \bnum{baseline}, the distance between the antennas;}
        \item \tr{$\theta$: o ângulo entre a linha de base e a direção do quasar;}{$\theta$: the angle between the baseline and the quasar direction;}
        \item \tr{$c$: a velocidade da luz.}{$c$: the speed of light.}
      \end{itemize}
    \end{column}
    \begin{column}{0.48\textwidth}
      \centering
      \begin{tikzpicture}[scale=0.95]
        \draw[mStoneMid!70, thick, fill=corFrame!8] (0,0) circle (1.7);
        \coordinate (A1) at (150:1.7);
        \coordinate (A2) at (40:1.7);
        \pic[corFrame, scale=1.5, transform shape] at (A1) {antena};
        \pic[corFrame, scale=1.5, transform shape] at (A2) {antena};
        \draw[mStoneMid, dashed] (A1) -- (A2);
        \node[rotulo, below] at ($(A1)!0.5!(A2)$) {$B$};
        % direção ao quasar u = (cos60, sin60); frentes de onda perpendiculares a u
        \draw[corCeu, thin, -{Stealth[length=1.8mm]}] ($(A1)+(1.3,2.25)$) -- (A1);
        \draw[corCeu, thin, -{Stealth[length=1.8mm]}] ($(A2)+(1.3,2.25)$) -- (A2);
        \draw[corCeu, thick] ($(A1)+(0.866,-0.5)$) -- ($(A1)+(-0.866,0.5)$);
        \draw[corCeu, thick] ($(A2)+(0.866,-0.5)$) -- ($(A2)+(-2.0,1.155)$);
        \node[rotulo, text=corCeu, above] at ($(A1)!0.5!(A2)+(1.3,2.25)$) {\tr{quasar (muito longe)}{quasar (very far)}};
        % trajeto extra: projeção da linha de base sobre u
        \pgfmathsetmacro{\sx}{(1.7*cos(40)-1.7*cos(150))*0.5+(1.7*sin(40)-1.7*sin(150))*0.866}
        \draw[corErro, line width=1.8pt] (A1) -- ($(A1)+(\sx*0.5,\sx*0.866)$);
        \node[rotulo, text=corErro, right] at ($(A1)+(\sx*0.25+0.05,\sx*0.433+0.12)$) {$B\cos\theta$};
        % ângulo theta: entre a linha de base e u
        \draw[mStoneMid] ($(A1)+(5:0.5)$) arc (5:60:0.5);
        \node[rotulo, text=mStone] at ($(A1)+(28:0.82)$) {$\theta$};
      \end{tikzpicture}
    \end{column}
  \end{columns}
\end{frame}

\begin{frame}{\tr{Do atraso para o ângulo}{From the delay to the angle}}
  \begin{columns}[T]
    \begin{column}{0.52\textwidth}
      \small
      \tr{Invertendo a fórmula:}{Inverting the formula:}
      \[ \cos\theta = \frac{c\,\tau}{B}, \qquad \delta\theta \approx \frac{c\,\delta\tau}{B\,\sin\theta} \]
      \tr{Com uma linha de base de $B = \num{10000}$ km:}{With a baseline of $B = \num{10000}$ km:}
      \medskip
      \centering
      \renewcommand{\arraystretch}{1.35}
      \begin{tabular}{lr}
        \toprule
        \tr{atraso máximo $B/c$}{maximum delay $B/c$} & $\num{33.4}$ ms \\
        \tr{precisão no atraso $\delta\tau$}{precision of the delay $\delta\tau$} & $10$ ps \\
        \tr{(luz em 10 ps anda)}{(light in 10 ps travels)} & $3$ mm \\
        \tr{precisão no ângulo $\delta\theta$}{precision of the angle $\delta\theta$} & \bnum{$\approx 60\uas$} \\
        \bottomrule
      \end{tabular}
    \end{column}
    \begin{column}{0.44\textwidth}
      \begin{armadilha}[title={\tr{Quão pequeno é 10 ps?}{How small is 10 ps?}}]
        \small\tr{$10$ ps $= 0{,}000\,000\,000\,01$ s: um centésimo de bilionésimo de segundo.
        Para medir o atraso tão bem, é preciso um relógio \alert{extraordinário} em cada
        antena.}{$10$ ps $= 0.000\,000\,000\,01$ s: one hundredth of a billionth of a second. To
        measure the delay that well, you need an \alert{extraordinary} clock at each antenna.}
      \end{armadilha}
    \end{column}
  \end{columns}
\end{frame}

\begin{frame}{\tr{Como gravar sem fio}{How to record without wires}}
  \begin{columns}[T]
    \begin{column}{0.46\textwidth}
      \small
      \tr{As antenas estão a milhares de quilômetros: não dá para ligar um fio entre elas.}{The
      antennas are thousands of kilometers apart: you cannot run a wire between them.}
      \medskip
      \begin{enumerate}
        \setlength{\itemsep}{3pt}
        \item \tr{Cada antena \bnum{grava} o ``chiado'' de rádio do quasar, junto com a hora de
              um \bnum{relógio atômico} (um maser de hidrogênio).}{Each antenna \bnum{records} the
              radio ``hiss'' of the quasar, together with the time from an \bnum{atomic clock} (a
              hydrogen maser).}
        \item \tr{As gravações vão para um computador, o \bnum{correlator}.}{The recordings go to
              a computer, the \bnum{correlator}.}
        \item \tr{Ele \bnum{desliza} uma gravação sobre a outra até as duas ficarem iguais. O
              deslizamento é o atraso $\tau$.}{It \bnum{slides} one recording over the other
              until the two match. The slide is the delay $\tau$.}
      \end{enumerate}
    \end{column}
    \begin{column}{0.52\textwidth}
      \centering
      \begin{tikzpicture}[x=1.0cm, y=1.0cm]
        \draw[corFrame, thick, domain=0:6, samples=240, variable=\t]
          plot ({\t},{1.9+0.30*sin(deg(7*\t))+0.22*sin(deg(17*\t+1))+0.16*sin(deg(29*\t+2))});
        \draw[corSist, thick, domain=0:6, samples=240, variable=\t]
          plot ({\t},{0.1+0.30*sin(deg(7*(\t-0.5)))+0.22*sin(deg(17*(\t-0.5)+1))+0.16*sin(deg(29*(\t-0.5)+2))});
        \node[rotulo, text=corFrame, left] at (0,1.9) {\tr{antena 1}{antenna 1}};
        \node[rotulo, text=corSist, left] at (0,0.1) {\tr{antena 2}{antenna 2}};
        \draw[mStoneMid, dashed] (2.3,0.1) -- (2.3,2.3);
        \draw[mStoneMid, dashed] (2.8,-0.35) -- (2.8,1.5);
        \draw[corErro, seta2] (2.3,1.05) -- (2.8,1.05);
        \node[rotulo, text=corErro, above] at (2.55,1.08) {$\tau$};
        \node[rotulo, text=mStoneMid, below] at (3,-0.6) {\tr{o mesmo sinal, gravado em dois lugares}{the same signal, recorded in two places}};
      \end{tikzpicture}
    \end{column}
  \end{columns}
\end{frame}

\begin{frame}{\tr{Sua vez}{Your turn}}
  \begin{columns}[T]
    \begin{column}{0.50\textwidth}
      \begin{regra}[title={\tr{Uma linha de base de 6\,000 km}{A 6\,000 km baseline}}]
        \small
        \begin{enumerate}
          \setlength{\itemsep}{4pt}
          \item[(a)] \tr{Um quasar está a $\theta = 60\gr$ da linha de base. Qual é o atraso
                $\tau$?}{A quasar is at $\theta = 60\gr$ from the baseline. What is the delay
                $\tau$?}
          \item[(b)] \tr{Para saber $\theta$ com $1\mas$ de precisão, com que precisão temos de
                medir $\tau$?}{To know $\theta$ to $1\mas$, how precisely must we measure
                $\tau$?}
          \item[(c)] \tr{Quanto a luz anda nesse tempo?}{How far does light travel in that
                time?}
        \end{enumerate}
      \end{regra}
    \end{column}
    \begin{column}{0.46\textwidth}
      \onslide<2->{
      \begin{ideia}[title={\tr{Respostas}{Answers}}]
        \small\raggedright
        (a) $\tau = \dfrac{6\times10^{6}\,\cos 60\gr}{\num{299792458}} \approx \bnum{\num{10.0}}$ ms.
        \medskip

        (b) $\delta\tau = \dfrac{B\sin\theta\,\delta\theta}{c} \approx \bnum{84}$ ps
        ($1\mas = \num{4.85e-9}$ rad).
        \medskip

        (c) $c \times 84$ ps $\approx \bnum{\num{2.5}}$ cm.
      \end{ideia}
      }
    \end{column}
  \end{columns}
\end{frame}

\begin{frame}{\tr{O atraso conhece os dois referenciais}{The delay knows both reference systems}}
  \begin{columns}[T]
    \begin{column}{0.46\textwidth}
      \small
      \tr{A conta do atraso, escrita com vetores:}{The delay, written with vectors:}
      \[ \tau = -\,\frac{\vec{b}\cdot\hat{s}}{c} \]
      \begin{itemize}
        \setlength{\itemsep}{2pt}
        \item \tr{$\vec b$: a linha de base --- conhecida no \fram{ITRF}, preso à Terra;}{$\vec b$: the baseline --- known in the \fram{ITRF}, fixed to the Earth;}
        \item \tr{$\hat s$: a direção do quasar --- no \fram{ICRF}, preso ao céu.}{$\hat s$: the direction of the quasar --- in the \fram{ICRF}, fixed to the sky.}
      \end{itemize}
      \smallskip
      \tr{Para multiplicar, os dois precisam estar \bnum{no mesmo referencial}: é preciso saber
      \emph{como a Terra está virada}. Isso é a aula 3.}{To multiply, both must be \bnum{in the
      same reference system}: you need to know \emph{how the Earth is turned}. That is class
      3.}
    \end{column}
    \begin{column}{0.50\textwidth}
      \centering
      \begin{tikzpicture}[scale=0.95, caixa/.append style={font=\footnotesize}]
        \node[caixa, draw=corFrame, fill=corFrame!10, text width=2.4cm] (B) at (0,2.4) {$\vec b$\\\tr{no chão (ITRF)}{on the ground (ITRF)}};
        \node[caixa, draw=corTempo, fill=corTempo!10, text width=2.2cm] (R) at (3.4,2.4) {\tr{como a Terra está virada}{how the Earth is turned}};
        \node[caixa, draw=corSist, fill=corSist!8, text width=2.4cm] (S) at (3.4,0.4) {$\hat s$\\\tr{no céu (ICRF)}{in the sky (ICRF)}};
        \node[caixa, draw=corErro, fill=corErro!6, text width=2.4cm] (T) at (0,0.4) {$\tau = -\vec b\cdot\hat s / c$\\\tr{o que se mede}{what is measured}};
        \draw[seta, mStoneMid] (B) -- (R);
        \draw[seta, mStoneMid] (R) -- (S);
        \draw[seta, mStoneMid] (S) -- (T);
        \node[rotulo, text=mStoneMid, below] at (1.7,1.6) {\tr{$\vec b$ agora no céu}{$\vec b$ now in the sky}};
      \end{tikzpicture}
      \bigskip

      \tr{Muitos atrasos, poucas incógnitas: \bnum{mínimos quadrados}, de novo. As incógnitas: antenas, quasares, orientação da Terra, relógios, atmosfera.}{Many delays, few unknowns: \bnum{least squares}, again. The unknowns: antennas, quasars, Earth orientation, clocks, atmosphere.}
    \end{column}
  \end{columns}
\end{frame}

\begin{frame}{\tr{O IVS e as antenas pelo mundo}{The IVS and the antennas around the world}}
  \begin{columns}[T]
    \begin{column}{0.44\textwidth}
      \small
      \tr{O \bnum{IVS} (\emph{International VLBI Service for Geodesy and Astrometry}) combina
      os observatórios de dezenas de países.}{The \bnum{IVS} (\emph{International VLBI Service
      for Geodesy and Astrometry}) combines the observatories of dozens of countries.}
      \smallskip
      \begin{itemize}
        \setlength{\itemsep}{3pt}
        \item \tr{sessões de \textbf{24 h}, algumas por semana: a orientação completa da Terra;}{\textbf{24 h} sessions, a few per week: the full orientation of the Earth;}
        \item \tr{sessões de \textbf{1 h} (``\emph{Intensives}''), quase todo dia: a rotação (UT1$-$UTC), com resultado rápido;}{\textbf{1 h} sessions (``\emph{Intensives}''), almost every day: the rotation (UT1$-$UTC), with fast results;}
        \item \tr{a nova geração, o \bnum{VGOS}: antenas de 12 m e meta de \bnum{1 mm}.}{the new generation, \bnum{VGOS}: 12 m antennas and a goal of \bnum{1 mm}.}
      \end{itemize}
    \end{column}
    \begin{column}{0.52\textwidth}
      \centering
      \begin{tikzpicture}[x=0.0150cm, y=0.0150cm]
        \draw[mStoneMid!60] (-180,-90) rectangle (180,90);
        \foreach \xx in {-120,-60,0,60,120} { \draw[mStoneMid!20] (\xx,-90) -- (\xx,90); }
        \foreach \yy in {-60,-30,0,30,60} { \draw[mStoneMid!20] (-180,\yy) -- (180,\yy); }
        \foreach \xx/\yy in {-71.5/42.6, 12.9/49.1, 27.7/-25.9, -159.7/22.1, 147.4/-42.8, 102.2/51.8,
                             121.2/31.1, -3.1/40.5, 11.9/78.9, 174.7/-36.4, -25.1/37.0, 11.9/57.4, 11.6/44.5} {
          \fill[corFrame] (\xx,\yy) circle (3.4);
        }
        \fill[corErro] (-38.4,-3.9) circle (5.2);
        \node[rotulo, text=corErro, below left] at (-38.4,-5) {Fortaleza};
        \node[rotulo, below right] at (-71.5,42.0) {Westford};
        \node[rotulo, above] at (12.9,52) {Wettzell};
        \node[rotulo, below] at (27.7,-28) {Hartebeesthoek};
        \node[rotulo, below] at (-159.7,19) {Kokee};
        \node[rotulo, below left] at (147.4,-44) {Hobart};
        \node[rotulo, text=mStoneMid] at (0,-108) {\tr{algumas estações da rede (posições aproximadas)}{some stations of the network (approximate positions)}};
      \end{tikzpicture}
      \par\smallskip
      \begin{cuidado}[title={\tr{No Brasil}{In Brazil}}]
        \scriptsize\tr{\textbf{Fortaleza}, Rádio Observatório Espacial do Nordeste (Eusébio, CE): antena de 14,2 m em VLBI desde \bnum{1993}; em novembro de 2025 foi inaugurada uma nova antena com terminal \bnum{VGOS}. O Hemisfério Sul tem poucas antenas: cada uma conta muito.}{\textbf{Fortaleza}, Northeast Space Radio Observatory (Eusébio, CE): 14.2 m antenna in VLBI since \bnum{1993}; in November 2025 a new antenna with a \bnum{VGOS} terminal was inaugurated. The Southern Hemisphere has few antennas: each one counts a lot.}
      \end{cuidado}
    \end{column}
  \end{columns}
\end{frame}

% ==================================================================
\section{\tr{Do ICRS ao ICRF}{From the ICRS to the ICRF}}

\begin{frame}{\tr{Sistema e frame, de novo}{System and frame, again}}
  \centering
  \small
  \renewcommand{\arraystretch}{1.5}
  \begin{tabular}{l>{\raggedright\arraybackslash}p{4.4cm}>{\raggedright\arraybackslash}p{4.4cm}}
    \toprule
    & \textcolor{corFrame}{\textbf{\tr{terrestre (Parte 1)}{terrestrial (Part 1)}}} &
      \textcolor{corSist}{\textbf{\tr{celeste (Parte 2)}{celestial (Part 2)}}} \\
    \midrule
    \sist{\tr{sistema}{system}} & ITRS & \textbf{ICRS} \\
    \fram{frame} & ITRF2020 & \textbf{ICRF3} \\
    \tr{o que a lista traz}{what the list holds} & \tr{estações: coordenada \emph{e} velocidade}{stations: coordinate \emph{and} velocity} & \tr{quasares: ascensão reta e declinação}{quasars: right ascension and declination} \\
    \tr{gira com a Terra?}{turns with the Earth?} & \tr{sim}{yes} & \bnum{\tr{não}{no}} \\
    \tr{quem adota}{who adopts it} & IUGG / IERS & \tr{IAU (com o IERS)}{IAU (with the IERS)} \\
    \bottomrule
  \end{tabular}
  \bigskip

  \tr{A mesma dupla: a \sist{regra escrita} e as \fram{coisas de verdade}.}{The same pair: the
  \sist{written rule} and the \fram{real things}.}
\end{frame}

\begin{frame}{\tr{A regra do ICRS, em quatro linhas}{The rule of the ICRS, in four lines}}
  \centering
  \small
  \renewcommand{\arraystretch}{1.2}
  \begin{tabular}{l>{\raggedright\arraybackslash}p{9.0cm}}
    \toprule
    \textbf{\tr{origem}{origin}} & \tr{o \bnum{baricentro} (centro de massa) do Sistema Solar}{the \bnum{barycenter} (center of mass) of the Solar System} \\
    \textbf{\tr{orientação}{orientation}} & \tr{eixos \bnum{presos a objetos extragalácticos distantes}, escolhidos para coincidir com o equador e o equinócio médios de J2000.0 (a $\approx 18\mas$), para não ``quebrar'' os catálogos antigos}{axes \bnum{tied to distant extragalactic objects}, chosen to coincide with the mean equator and equinox of J2000.0 (to $\approx 18\mas$), so as not to ``break'' the old catalogs} \\
    \textbf{\tr{escala}{scale}} & \tr{só direções; o espaço-tempo em volta usa o metro e o segundo do SI}{directions only; the surrounding space-time uses the SI meter and second} \\
    \textbf{\tr{tempo}{time}} & \tr{os eixos \bnum{não giram} em relação ao Universo distante}{the axes \bnum{do not rotate} relative to the distant Universe} \\
    \bottomrule
  \end{tabular}
  \medskip

  \begin{cuidado}[title={\tr{Repare no que não tem aqui}{Notice what is missing here}}]
    \small\tr{Nenhuma fonte, nenhuma coordenada: \alert{nada que dê para medir}. É o sistema --- e
    por isso ele não basta.}{No source, no coordinate: \alert{nothing you can measure}. This is
    the system --- and that is why it is not enough.}
  \end{cuidado}
\end{frame}

\begin{frame}{\tr{Por que o baricentro, e não a Terra?}{Why the barycenter, and not the Earth?}}
  \begin{columns}[T]
    \begin{column}{0.52\textwidth}
      \small
      \tr{O Sol não fica parado: ele e os planetas giram em torno do \bnum{centro de massa} do
      conjunto, o baricentro.}{The Sun does not stay put: it and the planets turn around the
      \bnum{center of mass} of the whole, the barycenter.}
      \medskip

      \tr{Só Júpiter já puxa o baricentro para \emph{fora} da superfície do Sol: a $1{,}07$ raio
      solar do centro.}{Jupiter alone already pulls the barycenter \emph{outside} the Sun's
      surface: $1.07$ solar radii from its center.}
      \medskip

      \tr{E a Terra \alert{acelera} ao dar a volta no Sol: um referencial no centro dela não é
      ``parado''. O baricentro é o melhor ponto de repouso que o Sistema Solar oferece.}{And the
      Earth \alert{accelerates} as it goes around the Sun: a reference system at its center is not
      ``still''. The barycenter is the best resting point the Solar System offers.}
    \end{column}
    \begin{column}{0.44\textwidth}
      \centering
      \begin{tikzpicture}[scale=1.3]
        \fill[corTempo!40] (0,0) circle (1.0);
        \draw[corTempo, thick] (0,0) circle (1.0);
        \draw[mStoneMid, dashed] (1.07,0) circle (1.07);
        \fill[mStone] (1.07,0) circle (0.05);
        \node[rotulo, text=mStone, right] at (1.12,0.1) {\tr{baricentro}{barycenter}};
        \fill[mStone] (0,0) circle (0.03);
        \node[rotulo, text=corTempo!60!black] at (-0.4,0.45) {\tr{Sol}{Sun}};
        \draw[corSist, seta] (0.22,0.98) arc (80:35:1.07 and 1.07);
        \node[rotulo, text=mStoneMid, align=center] at (1.1,-1.75) {\tr{o centro do Sol dá a volta}{the Sun's center circles}\\\tr{em torno do baricentro}{around the barycenter}};
      \end{tikzpicture}
    \end{column}
  \end{columns}
\end{frame}

\begin{frame}{\tr{O frame: o ICRF}{The frame: the ICRF}}
  \begin{columns}[T]
    \begin{column}{0.36\textwidth}
      \small
      \tr{O \fram{ICRF} é uma \bnum{lista de posições}: de cada quasar, o nome, a ascensão reta e a
      declinação --- com a incerteza.}{The \fram{ICRF} is a \bnum{list of positions}: for each
      quasar, the name, the right ascension and the declination --- with the uncertainty.}
      \medskip

      \tr{Como a lista é medida \emph{com o VLBI}, a precisão é de milésimos de segundo de arco e
      melhor.}{Since the list is measured \emph{with VLBI}, the precision is thousandths of an
      arcsecond and better.}
      \medskip

      \tr{Mesma forma do ITRF: nomes e números. Só que os números são \bnum{direções}, e não
      envelhecem com a precessão.}{Same form as the ITRF: names and numbers. Only the numbers are
      \bnum{directions}, and they do not age with precession.}
    \end{column}
    \begin{column}{0.62\textwidth}
      \centering
      \footnotesize
      \renewcommand{\arraystretch}{1.35}
      \setlength{\tabcolsep}{4pt}
      \begin{tabular}{llll}
        \toprule
        \tr{fonte}{source} & $\alpha$ & $\delta$ & \tr{o que é}{what it is} \\
        \midrule
        3C 273 & \hms{12}{29}{6.70} & \dms{+2}{3}{8.6} & \tr{quasar (1963)}{quasar (1963)} \\
        3C 279 & \hms{12}{56}{11.17} & \dms{-5}{47}{21.5} & \tr{quasar variável}{variable quasar} \\
        3C 84 & \hms{3}{19}{48.16} & \dms{+41}{30}{42.1} & \tr{galáxia NGC 1275}{galaxy NGC 1275} \\
        3C 345 & \hms{16}{42}{58.81} & \dms{+39}{48}{37.0} & \tr{quasar com jato}{quasar with a jet} \\
        M87 & \hms{12}{30}{49.42} & \dms{+12}{23}{28.0} & \tr{galáxia gigante}{giant galaxy} \\
        \bottomrule
      \end{tabular}
      \par\smallskip
      {\scriptsize\textcolor{mStoneMid}{\tr{Fontes famosas; valores de catálogo arredondados (J2000.0). ``3C'' = Terceiro Catálogo de Cambridge, de fontes de rádio.}{Famous sources; rounded catalog values (J2000.0). ``3C'' = Third Cambridge Catalogue of radio sources.}}}
    \end{column}
  \end{columns}
\end{frame}

\begin{frame}{\tr{Três gerações do ICRF}{Three generations of the ICRF}}
  \begin{columns}[T]
    \begin{column}{0.50\textwidth}
      \centering
      \begin{tikzpicture}[x=0.00095cm, y=0.8cm]
        \foreach \nome/\tot/\nd/\yy in {ICRF1/608/212/2.4, ICRF2/3414/295/1.2, ICRF3/4536/303/0} {
          \fill[corFrame!35] (0,\yy) rectangle (\tot,{\yy+0.7});
          \fill[corSist] (0,\yy) rectangle (\nd,{\yy+0.7});
          \node[rotulo, left] at (-60,{\yy+0.35}) {\nome};
          \node[rotulo, right] at ({\tot+60},{\yy+0.35}) {\tot\,/\,\textcolor{corSist}{\nd}};
        }
        \draw[mStoneMid] (0,-0.4) -- (0,3.4);
      \end{tikzpicture}
      \par
      {\scriptsize\textcolor{corFrame!80!black}{\tr{fontes no total}{sources in total}} $/$ \textcolor{corSist}{\tr{definidoras}{defining}}}
    \end{column}
    \begin{column}{0.46\textwidth}
      \footnotesize
      \renewcommand{\arraystretch}{1.3}
      \setlength{\tabcolsep}{3pt}
      \begin{tabular}{@{}lll@{}}
        \toprule
        & \tr{adotado (IAU)}{adopted (IAU)} & \tr{em vigor}{in force} \\
        \midrule
        ICRF1 & 1997 \tr{(Kyoto)}{(Kyoto)} & 1/1/1998 \\
        ICRF2 & 2009 \tr{(Rio de Janeiro)}{(Rio de Janeiro)} & 1/1/2010 \\
        ICRF3 & 2018 \tr{(Viena)}{(Vienna)} & 1/1/2019 \\
        \bottomrule
      \end{tabular}
      \medskip

      \tr{O ICRF2 foi adotado na Assembleia Geral da IAU \bnum{no Rio de Janeiro}. O ICRF3 traz três
      frequências: 8,4 GHz (4\,536 fontes), 24 GHz (824) e 32 GHz (678).}{ICRF2 was adopted at the IAU General Assembly \bnum{in Rio de Janeiro}. ICRF3 has three frequencies: 8.4 GHz (4,536 sources), 24 GHz (824) and 32 GHz (678).}
    \end{column}
  \end{columns}
\end{frame}

\begin{frame}{\tr{As fontes definidoras}{The defining sources}}
  \begin{columns}[T]
    \begin{column}{0.46\textwidth}
      \small
      \tr{Nem toda fonte da lista manda na orientação dos eixos. Escolhem-se as \bnum{definidoras}
      (no ICRF3, \bnum{303} entre 4\,536): posição estável, estrutura compacta e \bnum{bem
      espalhadas} pelo céu.}{Not every source on the list controls the orientation of the axes.
      The \bnum{defining} ones are chosen (in the ICRF3, \bnum{303} out of 4,536): stable
      position, compact structure and \bnum{well spread} over the sky.}
      \medskip

      \tr{É como escolher os melhores marcos de uma rede: os outros ficam na lista, mas não
      amarram os eixos.}{It is like choosing the best monuments of a network: the others stay on
      the list, but do not tie the axes.}
      \medskip

      \begin{cuidado}[title={\tr{Por que só os estáveis?}{Why only the stable ones?}}]
        \footnotesize\tr{Um quasar tem \emph{jato}: o ponto de rádio mais brilhante pode andar um pouco quando o jato muda. Isso mexeria nos eixos.}{A quasar has a \emph{jet}: the brightest radio spot can wander a little when the jet changes. That would shake the axes.}
      \end{cuidado}
    \end{column}
    \begin{column}{0.52\textwidth}
      \centering
      \begin{tikzpicture}[scale=1.0]
        \draw[mStoneMid!70, fill=corCeu!5] (0,0) ellipse (3 and 1.5);
        \draw[mStoneMid!30] (-3,0) -- (3,0);
        \draw[mStoneMid!30] (0,-1.5) -- (0,1.5);
        \mapaFontes{corSist}
        \node[rotulo, text=mStoneMid] at (0,-1.8) {\tr{303 fontes definidoras (distribuição ilustrativa)}{303 defining sources (illustrative distribution)}};
      \end{tikzpicture}
    \end{column}
  \end{columns}
\end{frame}

\begin{frame}{\tr{A escada dos ângulos}{The ladder of angles}}
  \begin{columns}[T]
    \begin{column}{0.50\textwidth}
      \centering
      \begin{tikzpicture}[x=1.0cm, y=0.56cm]
        \draw[mStoneMid] (0,0) -- (0,9.3);
        \foreach \yy/\lab in {0.0/{$10\uas$}, 1.0/{$100\uas$}, 2.0/{$1\mas$}, 3.0/{$10\mas$}, 4.0/{$100\mas$}, 5.0/{$1\asec$}, 6.0/{$10\asec$}, 7.0/{$100\asec$}, 8.0/{$1\,000\asec$}} {
          \draw[mStoneMid] (-0.08,\yy) -- (0.08,\yy);
          \node[rotulo, left] at (-0.12,\yy) {\lab};
        }
        \fill[corCeu] (0,8.26) circle (0.07); \node[rotulo, right, text=corCeu] at (0.15,8.26) {\tr{Lua cheia ($0{,}5\gr$)}{full Moon ($0.5\gr$)}};
        \fill[corTempo] (0,6.78) circle (0.07); \node[rotulo, right, text=corTempo] at (0.15,6.78) {\tr{olho nu $\approx 60\asec$}{naked eye $\approx 60\asec$}};
        \fill[corSist] (0,4.0) circle (0.07); \node[rotulo, right, text=corSist] at (0.15,4.0) {\tr{FK5 (1988): décimos de $\asec$}{FK5 (1988): tenths of $\asec$}};
        \fill[corFrame] (0,2.0) circle (0.07); \node[rotulo, right, text=corFrame] at (0.15,2.0) {Hipparcos (1997): $\approx 1\mas$};
        \fill[corFrame] (0,0.6) circle (0.07);
        \draw[corFrame] (0.07,0.6) -- (0.35,1.05); \node[rotulo, right, text=corFrame] at (0.35,1.05) {ICRF2 (2010): $40\uas$};
        \fill[corErro] (0,0.48) circle (0.07);
        \draw[corErro] (0.07,0.48) -- (0.35,-0.05); \node[rotulo, right, text=corErro] at (0.35,-0.05) {ICRF3 (2019): $30\uas$};
      \end{tikzpicture}
    \end{column}
    \begin{column}{0.46\textwidth}
      \begin{ideia}[title={\tr{O que isso vale na Terra}{What that is worth on the ground}}]
        \small
        \renewcommand{\arraystretch}{1.15}
        \begin{tabular}{@{}rl@{}}
          $1\asec$ & $= \num{30.9}$ m \\
          $1\mas$ & $= \num{3.1}$ cm \\
          $30\uas$ & $= \bnum{\num{0.9}}$ mm \\
        \end{tabular}
      \end{ideia}
      \smallskip
      \footnotesize
      \tr{Se o céu de referência errasse $30\uas$ na orientação, \alert{todos os pontos da Terra}
      ficariam deslocados cerca de \bnum{1 mm}. Por isso a meta de 1 mm para o ITRF exige um céu
      com esta qualidade.}{If the reference sky erred by $30\uas$ in orientation, \alert{every
      point on the Earth} would be displaced by about \bnum{1 mm}. That is why the 1 mm goal for
      the ITRF demands a sky of this quality.}
      \smallskip

      \tr{Para ter uma ideia: $30\uas$ é o tamanho de uns 6 cm \emph{na Lua}, visto da
      Terra.}{To get a feel: $30\uas$ is the size of about 6 cm \emph{on the Moon}, seen from the
      Earth.}
    \end{column}
  \end{columns}
\end{frame}

\begin{frame}{\tr{Sua vez}{Your turn}}
  \begin{columns}[T]
    \begin{column}{0.52\textwidth}
      \begin{regra}[title={\tr{Três continhas}{Three little sums}}]
        \small
        \begin{enumerate}
          \setlength{\itemsep}{4pt}
          \item[(a)] \tr{Quantos centímetros, na superfície da Terra, vale $1\mas$? (use
                $1\asec = 30{,}9$ m)}{How many centimeters, on the Earth's surface, is $1\mas$
                worth? (use $1\asec = 30.9$ m)}
          \item[(b)] \tr{Um erro de $0{,}3\mas$ na orientação do céu desloca os pontos da Terra
                de quanto?}{An error of $0.3\mas$ in the orientation of the sky displaces points
                on the Earth by how much?}
          \item[(c)] \tr{Um quasar anda de lado a $1\,000$ km/s, a $10$ bilhões de anos-luz.
                Quantos $\mu$as por ano? (Dica: $v/d$, em radianos por ano.)}{A quasar moves
                sideways at $1\,000$ km/s, $10$ billion light-years away. How many $\mu$as per
                year? (Hint: $v/d$, in radians per year.)}
        \end{enumerate}
      \end{regra}
    \end{column}
    \begin{column}{0.44\textwidth}
      \onslide<2->{
      \begin{ideia}[title={\tr{Respostas}{Answers}}]
        \small\raggedright
        (a) $\num{0.001}\asec \times \num{30.9}$ m $\approx \bnum{\num{3.1}}$ cm.
        \medskip

        (b) $\num{0.3} \times \num{3.1}$ cm $\approx \bnum{\num{0.9}}$ cm \tr{(quase 1 cm!)}{(almost 1 cm!)}.
        \medskip

        (c) $v/d = \dfrac{10^{6}}{10^{10}\times \num{9.46}\times10^{15}} \times \num{3.16}\times10^{7}$
        $\approx \num{3.3}\times10^{-13}$ rad/\tr{ano}{yr} $\approx \bnum{\num{0.07}}\uas$/\tr{ano}{yr}.
      \end{ideia}
      }
    \end{column}
  \end{columns}
\end{frame}

% ==================================================================
\section{\tr{Hipparcos, Gaia e a aceleração da Galáxia}{Hipparcos, Gaia and the Galaxy's acceleration}}

\begin{frame}{\tr{O céu em luz visível: Hipparcos e Gaia}{The sky in visible light: Hipparcos and Gaia}}
  \begin{columns}[T]
    \begin{column}{0.50\textwidth}
      \small
      \tr{Os quasares são fracos em luz visível. Para o céu que se vê, a agência espacial
      europeia (ESA) lançou dois satélites:}{Quasars are faint in visible light. For the sky we
      see, the European Space Agency (ESA) launched two satellites:}
      \smallskip
      \begin{itemize}
        \setlength{\itemsep}{1pt}
        \item \textbf{Hipparcos} (1989--1993): \tr{catálogo de \num{118218} estrelas, com
              $\approx 1\mas$ de precisão (1997);}{a catalog of \num{118218} stars, with
              $\approx 1\mas$ precision (1997);}
        \item \textbf{Gaia} \tr{(lançado em 2013; observou até 15/1/2025)}{(launched 2013; observed
              until 15 Jan 2025)}: \tr{cerca de \bnum{1,8 bilhão} de fontes no DR3 (2022), entre
              elas $\approx$ \bnum{1,6 milhão de quasares}.}{about \bnum{1.8 billion} sources in DR3
              (2022), among them $\approx$ \bnum{1.6 million quasars}.}
      \end{itemize}
      \smallskip
      \tr{Os quasares do Gaia são mais fracos e menos precisos um a um --- mas são
      \bnum{milhões}.}{Gaia's quasars are fainter and less precise one by one --- but there are
      \bnum{millions}.}
    \end{column}
    \begin{column}{0.48\textwidth}
      \centering
      \begin{tikzpicture}[x=0.14cm, y=1.0cm]
        \draw[mStoneMid, -{Stealth[length=2mm]}] (0,0) -- (42.6,0);
        \foreach \aa in {1990,2000,2010,2020} { \draw[mStoneMid] ({\aa-1985},-0.08) -- ({\aa-1985},0.08); \node[rotulo, below] at ({\aa-1985},-0.1) {\aa}; }
        \fill[corSist] (4,0) circle (0.12); \node[rotulo, above, text=corSist, align=center] at (4,0.2) {Hipparcos\\\tr{lançado}{launched}};
        \fill[corSist] (12,0) circle (0.12); \node[rotulo, below, text=corSist, align=center] at (12,-0.55) {\tr{catálogo}{catalog}\\1997};
        \fill[corFrame] (28,0) circle (0.12); \node[rotulo, above, text=corFrame, align=center] at (28,0.2) {Gaia\\\tr{lançado}{launched}};
        \fill[corFrame!60] (37,0) circle (0.1); \node[rotulo, below, text=corFrame, align=center] at (37,-0.55) {DR3\\2022};
        \draw[corFrame, dashed] (41.9,0) circle (0.1); \node[rotulo, above, text=corFrame, align=center] at (40.6,0.2) {DR4\\\tr{dez./2026}{Dec. 2026}};
      \end{tikzpicture}
      \par\smallskip
      {\scriptsize\textcolor{mStoneMid}{\tr{A próxima liberação de dados, a DR4, está prevista para 2/12/2026.}{The next data release, DR4, is planned for 2 Dec 2026.}}}
    \end{column}
  \end{columns}
\end{frame}

\begin{frame}{\tr{Um sistema, vários frames}{One system, several frames}}
  \centering
  \small
  \renewcommand{\arraystretch}{1.35}
  \begin{tabular}{llp{3.2cm}p{3.0cm}}
    \toprule
    \fram{frame} & \tr{luz}{light} & \tr{fontes}{sources} & \tr{adotado pela IAU}{adopted by the IAU} \\
    \midrule
    \textbf{ICRF3} & \tr{rádio (VLBI)}{radio (VLBI)} & \tr{4\,536 (303 definidoras)}{4,536 (303 defining)} & \tr{2018 (Res.~B2)}{2018 (Res.~B2)} \\
    \textbf{Gaia-CRF3} & \tr{visível}{visible} & \tr{$\approx$ 1,6 milhão de quasares}{$\approx$ 1.6 million quasars} & \tr{2021 (Res.~B3)}{2021 (Res.~B3)} \\
    \textbf{HCRF} \tr{(Hipparcos)}{(Hipparcos)} & \tr{visível}{visible} & \tr{118\,218 estrelas}{118,218 stars} & \tr{1997 (Res.~B2)}{1997 (Res.~B2)} \\
    \bottomrule
  \end{tabular}
  \medskip

  \small
  \tr{O \sist{ICRS} é um só. Os \fram{frames} são vários, cada um numa ``cor'' de luz. Eles se
  alinham pelas fontes em comum: \bnum{3\,181} fontes do ICRF3 também estão no Gaia-CRF3.}{The
  \sist{ICRS} is one. The \fram{frames} are several, each in a ``color'' of light. They are
  aligned through the sources they share: \bnum{3,181} ICRF3 sources are also in Gaia-CRF3.}
  \medskip

  {\footnotesize\tr{Desde 1/1/2022 o ICRS é realizado pelo ICRF3 (rádio) \emph{e} pelo Gaia-CRF3 (visível). O Gaia-CRF3 em grande parte substituiu o Hipparcos.}{Since 1 Jan 2022 the ICRS has been realized by the ICRF3 (radio) \emph{and} the Gaia-CRF3 (visible). Gaia-CRF3 has largely superseded Hipparcos.}}
\end{frame}

\begin{frame}{\tr{Nem os quasares são parados: nós aceleramos}{Not even quasars are still: we accelerate}}
  \begin{columns}[T]
    \begin{column}{0.54\textwidth}
      \small
      \tr{O Sistema Solar gira em torno do centro da Galáxia a uns \bnum{230 km/s}: uma volta a cada
      $\approx$ \bnum{220 milhões de anos}.}{The Solar System goes around the center of the Galaxy
      at about \bnum{230 km/s}: one turn every $\approx$ \bnum{220 million years}.}
      \medskip

      \tr{Uma órbita é uma \bnum{aceleração} para o centro: $\approx 2{,}7\times10^{-10}$ m/s$^2$.
      Efeito: todos os quasares parecem derivar na direção do centro da Galáxia, por
      \bnum{$5{,}8\uas$/ano}.}{An orbit is an \bnum{acceleration} toward the center: $\approx
      2.7\times10^{-10}$ m/s$^2$. Effect: all quasars seem to drift toward the center of the Galaxy,
      by \bnum{$5.8\uas$/yr}.}
      \medskip

      \tr{Em 10 anos são $58\uas$ --- mais que o piso de ruído do ICRF3 ($30\uas$). Por isso o
      ICRF3 \bnum{corrige esse efeito}: as posições valem para a época \bnum{2015,0}.}{In 10 years
      that is $58\uas$ --- more than the ICRF3 noise floor ($30\uas$). That is why the ICRF3
      \bnum{corrects this effect}: the positions are valid for epoch \bnum{2015.0}.}
    \end{column}
    \begin{column}{0.42\textwidth}
      \centering
      \begin{tikzpicture}[scale=0.92]
        \fill[corCeu!10] (0,0) circle (2.3);
        \draw[corCeu, thick] (0,0) circle (2.3);
        \fill[corTempo!60] (0,0) circle (0.32);
        \draw[mStoneMid, dashed] (0,0) circle (1.6);
        \fill[corSist] (1.6,0) circle (0.08);
        \node[rotulo, text=corSist, above right] at (1.62,0.05) {\tr{Sol}{Sun}};
        \draw[corErro, seta, line width=1.4pt] (1.55,0) -- (0.85,0);
        \node[rotulo, text=corErro, below] at (1.15,-0.05) {$a$};
        \draw[corSist, seta] (1.6,0.05) -- (1.62,0.85);
        \node[rotulo, text=corSist, right] at (1.65,0.5) {$v$};
        \node[rotulo, text=corTempo, below] at (0,-0.35) {\tr{centro}{center}};
        \node[rotulo, text=mStoneMid] at (0,-2.65) {\tr{visto de cima, fora de escala}{seen from above, not to scale}};
      \end{tikzpicture}
    \end{column}
  \end{columns}
\end{frame}

\begin{frame}{\tr{Para levar para casa --- aula 2}{Take-home --- class 2}}
  \footnotesize
  \begin{enumerate}
    \setlength{\itemsep}{2pt}
    \item \tr{Para ter um céu ``parado'', usam-se objetos \bnum{tão distantes que não se mexem}:
          os \bnum{quasares} (movimento aparente $\approx 0{,}07\uas$/ano, contra $10\asec$/ano da
          estrela de Barnard).}{To get a ``still'' sky, we use objects \bnum{so distant that they do
          not move}: \bnum{quasars} (apparent motion $\approx 0.07\uas$/yr, against $10\asec$/yr for
          Barnard's Star).}
    \item \tr{\bnum{VLBI}: duas antenas distantes gravam o mesmo sinal; o atraso $\tau = B\cos\theta/c$
          dá o ângulo. Com $\delta\tau \approx 10$ ps, $\approx 60\uas$.}{\bnum{VLBI}: two distant
          antennas record the same signal; the delay $\tau = B\cos\theta/c$ gives the angle. With
          $\delta\tau \approx 10$ ps, $\approx 60\uas$.}
    \item \tr{\sist{ICRS} (sistema): origem no baricentro, eixos presos a objetos extragalácticos,
          sem rotação. \fram{ICRF} (frame): a lista de $\approx$ 4\,500 quasares ($303$
          definidoras) --- ICRF3, em vigor desde 2019.}{\sist{ICRS} (system): origin at the
          barycenter, axes tied to extragalactic objects, no rotation. \fram{ICRF} (frame): the list
          of $\approx$ 4,500 quasars ($303$ defining) --- ICRF3, in force since 2019.}
    \item \tr{$30\uas$ no céu valem $\approx$ \bnum{1 mm} na Terra.}{$30\uas$ in the sky are worth
          $\approx$ \bnum{1 mm} on the ground.}
    \item \tr{Hipparcos (1997) e Gaia (2013--2025) cuidam da luz visível: \bnum{um sistema, vários
          frames}. E até os quasares ``andam'': a aceleração da Galáxia, $5{,}8\uas$/ano, já é
          corrigida no ICRF3.}{Hipparcos (1997) and Gaia (2013--2025) take care of visible light:
          \bnum{one system, several frames}. And even quasars ``move'': the Galaxy's acceleration,
          $5.8\uas$/yr, is already corrected in the ICRF3.}
  \end{enumerate}
  \begin{ideia}
    \scriptsize \tr{\textbf{Na próxima aula:} como a Terra está virada em relação a esse céu ---
    rotação, precessão, nutação, movimento polar --- e os \bnum{EOP}.}{\textbf{Next class:} how
    the Earth is turned relative to that sky --- rotation, precession, nutation, polar motion ---
    and the \bnum{EOP}.}
  \end{ideia}
\end{frame}


% ==================================================================
%  AULA 3 · A PONTE ENTRE O CÉU E O CHÃO
% ==================================================================
\begin{frame}[standout]
  \tr{Parte 2}{Part 2}\\[6pt]
  {\normalsize \tr{A ponte entre o céu e o chão}{The bridge between the sky and the ground}}\\[4pt]
  {\scriptsize \tr{aula 3 de 4}{class 3 of 4}}
\end{frame}

\section{\tr{Duas visões, uma rotação}{Two views, one rotation}}

\begin{frame}{\tr{Onde paramos: a aula 2}{Where we left off: class 2}}
  \begin{columns}[T]
    \begin{column}{0.56\textwidth}
      \small
      \tr{Na aula 2:}{In class 2:}
      \begin{enumerate}
        \setlength{\itemsep}{4pt}
        \item \tr{os \bnum{quasares} estão tão longe que não se mexem: são a referência parada;}{\bnum{quasars} are so far away that they do not move: they are the still reference;}
        \item \tr{o \bnum{VLBI} mede o atraso $\tau = B\cos\theta/c$ entre duas antenas e dá a direção a $\approx 60\uas$;}{\bnum{VLBI} measures the delay $\tau = B\cos\theta/c$ between two antennas and gives the direction to $\approx 60\uas$;}
        \item \tr{o \sist{ICRS} (sistema) e o \fram{ICRF} (frame): origem no baricentro, eixos presos aos quasares, sem rotação.}{the \sist{ICRS} (system) and the \fram{ICRF} (frame): origin at the barycenter, axes tied to quasars, no rotation.}
      \end{enumerate}
    \end{column}
    \begin{column}{0.40\textwidth}
      \begin{cuidado}[title={\tr{A pergunta de hoje}{Today's question}}]
        \tr{O céu está parado. A Terra está \alert{virada} --- mas virada de quanto, \alert{neste instante}?}{The sky is still. The Earth is \alert{turned} --- but turned by how much, \alert{right now}?}
        \medskip

        \tr{Resposta provisória: de \bnum{cinco números por dia}.}{Provisional answer: by \bnum{five numbers per day}.}
      \end{cuidado}
    \end{column}
  \end{columns}
\end{frame}

\begin{frame}{\tr{As perguntas de hoje}{Today's questions}}
  \begin{columns}[T]
    \begin{column}{0.58\textwidth}
      \begin{enumerate}
        \setlength{\itemsep}{8pt}
        \item \tr{Quais são as ``viradas'' da Terra?}{What are the Earth's ``turns''?}
        \item \tr{Como escrever uma virada com números?}{How do we write a turn with numbers?}
        \item \tr{De que tamanho é cada uma, no céu e no chão?}{How big is each one, in the sky and on the ground?}
        \item \tr{Quem mede, e quem publica? (os \bnum{EOP} e o \bnum{IERS})}{Who measures, and who publishes? (the \bnum{EOP} and the \bnum{IERS})}
      \end{enumerate}
    \end{column}
    \begin{column}{0.38\textwidth}
      \begin{ideia}[title={\tr{Combinado}{Deal}}]
        \tr{No fim da aula, você responde às quatro --- com as suas palavras.}{By the end of the class you will answer all four --- in your own words.}
      \end{ideia}
    \end{column}
  \end{columns}
\end{frame}

\begin{frame}{\tr{O mesmo ponto, dois endereços}{The same point, two addresses}}
  \begin{columns}[T]
    \begin{column}{0.46\textwidth}
      \small
      \tr{Uma antena está no chão: no \fram{ITRF} as coordenadas dela \bnum{não mudam}. Mas, para o céu, ela não para: dá uma volta por dia.}{An antenna sits on the ground: in the \fram{ITRF} its coordinates \bnum{do not change}. But, to the sky, it never stops: it makes one turn per day.}
      \medskip

      \tr{O quasar está parado no céu: a direção $\hat s$ dele vale no \sist{ICRF}.}{The quasar is still in the sky: its direction $\hat s$ holds in the \sist{ICRF}.}
      \medskip

      \tr{Para o atraso $\tau = -\vec b\cdot\hat s/c$, os dois vetores precisam estar \bnum{no mesmo endereço}, isto é, no mesmo referencial. Converter de um para o outro é o trabalho da \bnum{ponte}.}{For the delay $\tau = -\vec b\cdot\hat s/c$, both vectors must be \bnum{at the same address}, that is, in the same reference system. Converting from one to the other is the job of the \bnum{bridge}.}
    \end{column}
    \begin{column}{0.52\textwidth}
      \centering
      \begin{tikzpicture}[scale=1.12]
        \foreach \ox/\ang/\lab in {0/20/{\tr{agora}{now}}, 4.3/110/{\tr{6 h depois}{6 h later}}} {
          \begin{scope}[shift={(\ox,0)}]
            \draw[mStoneMid!70, thick, fill=corFrame!8] (0,0) circle (1.0);
            \draw[corFrame, seta] (0,0) -- ({0.9*cos(\ang)},{0.9*sin(\ang)});
            \fill[corFrame] ({cos(\ang)},{sin(\ang)}) circle (0.07);
            \node[font=\footnotesize, text=corFrame] at ({1.3*cos(\ang)},{1.3*sin(\ang)}) {$\vec b$};
            \draw[corCeu, seta] (0,0) -- ({1.9*cos(55)},{1.9*sin(55)});
            \node[font=\footnotesize, text=corCeu] at ({2.15*cos(55)},{2.15*sin(55)}) {$\hat s$};
            \node[font=\footnotesize, below] at (0,-1.15) {\lab};
          \end{scope}
        }
        \draw[mStoneMid, -{Stealth[length=2mm]}] (1.45,-0.55) .. controls (2.1,-1.0) and (2.9,-1.0) .. (3.55,-0.55);
        \node[font=\footnotesize, text=corTempo] at (2.5,-1.3) {$+90\gr$};
      \end{tikzpicture}
      \par\smallskip
      {\scriptsize\textcolor{mStoneMid}{\tr{Visto de cima. $\vec b$ gira com a Terra; $\hat s$ fica. O ângulo entre os dois muda, e o atraso também.}{Seen from above. $\vec b$ turns with the Earth; $\hat s$ stays. The angle between them changes, and so does the delay.}}}
    \end{column}
  \end{columns}
\end{frame}

\begin{frame}{\tr{Girar os eixos}{Turning the axes}}
  \begin{columns}[T]
    \begin{column}{0.52\textwidth}
      \small
      \tr{Um ponto $P$ tem coordenadas $(x,y)$ nos eixos do céu. Gire os eixos de um ângulo $\varphi$: o ponto não se mexeu, mas as \bnum{coordenadas mudam}:}{A point $P$ has coordinates $(x,y)$ in the sky axes. Turn the axes by an angle $\varphi$: the point did not move, but the \bnum{coordinates change}:}
      \[ \begin{aligned} x' &= x\cos\varphi + y\sin\varphi \\ y' &= -x\sin\varphi + y\cos\varphi \end{aligned} \]
      \tr{Em forma de tabela (uma \bnum{matriz}):}{As a table (a \bnum{matrix}):}
      \[ \begin{pmatrix} x' \\ y' \end{pmatrix} = \begin{pmatrix} \cos\varphi & \sin\varphi \\ -\sin\varphi & \cos\varphi \end{pmatrix} \begin{pmatrix} x \\ y \end{pmatrix} \]
    \end{column}
    \begin{column}{0.44\textwidth}
      \centering
      \begin{tikzpicture}[x=1.9cm, y=1.9cm]
        \draw[corCeu, seta] (-0.15,0) -- (1.3,0) node[right, rotulo, text=corCeu] {$x$};
        \draw[corCeu, seta] (0,-0.2) -- (0,1.05) node[above, rotulo, text=corCeu] {$y$};
        \draw[corFrame, seta] (0,0) -- (30:1.3) node[right, rotulo, text=corFrame] {$x'$};
        \draw[corFrame, seta] (0,0) -- (120:1.05) node[above, rotulo, text=corFrame] {$y'$};
        \draw[corErro] (0.42,0) arc (0:30:0.42);
        \node[rotulo, text=corErro] at (14:0.56) {$\varphi$};
        \draw[mStoneMid, dotted] (1,0) -- (0.75,0.433);
        \draw[mStoneMid, dotted] (1,0) -- (0.25,-0.433);
        \fill[mStone] (1,0) circle (0.03);
        \node[rotulo, below right] at (1.0,0) {$P$};
        \node[rotulo, text=corFrame, anchor=south east] at (0.74,0.44) {$x'\!=\!\tr{0{,}87}{0.87}$};
        \node[rotulo, text=corFrame, below] at (0.25,-0.44) {$y'\!=\!\tr{-0{,}5}{-0.5}$};
      \end{tikzpicture}
      \par
      {\scriptsize\textcolor{mStoneMid}{\tr{céu: azul; chão (girado): verde}{sky: blue; ground (turned): green}}}
      \smallskip
      \begin{ideia}[title={\tr{Exemplo}{Example}}]
        \footnotesize\tr{$P = (1,\,0)$ e $\varphi = 30\gr$ dão $P' = (\cos30\gr,\,-\sin30\gr) = (\num{0.866},\,\num{-0.5})$.}{$P = (1,\,0)$ and $\varphi = 30\gr$ give $P' = (\cos30\gr,\,-\sin30\gr) = (\num{0.866},\,\num{-0.5})$.}
      \end{ideia}
    \end{column}
  \end{columns}
\end{frame}

\begin{frame}{\tr{Sua vez}{Your turn}}
  \begin{columns}[T]
    \begin{column}{0.52\textwidth}
      \begin{regra}[title={\tr{Girando eixos}{Turning axes}}]
        \small
        \begin{enumerate}
          \setlength{\itemsep}{4pt}
          \item[(a)] \tr{$P = (1,\,0)$ e $\varphi = 90\gr$. Quais são $(x',y')$?}{$P = (1,\,0)$ and $\varphi = 90\gr$. What are $(x',y')$?}
          \item[(b)] \tr{E com $\varphi = 180\gr$?}{And with $\varphi = 180\gr$?}
          \item[(c)] \tr{Um satélite está em $(26\,560,\,0)$ km nos eixos do céu. Os eixos do chão giram $90\gr$. Onde ele está nos eixos do chão?}{A satellite is at $(26\,560,\,0)$ km in the sky axes. The ground axes turn $90\gr$. Where is it in the ground axes?}
        \end{enumerate}
      \end{regra}
    \end{column}
    \begin{column}{0.44\textwidth}
      \onslide<2->{
      \begin{ideia}[title={\tr{Respostas}{Answers}}]
        \small\raggedright
        (a) $(\cos90\gr,\,-\sin90\gr) = \bnum{$(0,\,-1)$}$.
        \medskip

        (b) $(\cos180\gr,\,-\sin180\gr) = \bnum{$(-1,\,0)$}$.
        \medskip

        (c) $\bnum{$(0,\,-26\,560)$}$ km: \tr{o satélite ``parado'' no céu aparece \emph{andando para o oeste} no chão.}{the satellite ``still'' in the sky appears \emph{moving west} on the ground.}
      \end{ideia}
      }
    \end{column}
  \end{columns}
\end{frame}

\begin{frame}{\tr{Em três dimensões: a ordem importa}{In three dimensions: the order matters}}
  \begin{columns}[T]
    \begin{column}{0.54\textwidth}
      \small
      \tr{Na Terra o giro principal é em torno do eixo $z$ (o eixo do polo): $x$ e $y$ mudam, $z$ \bnum{fica}.}{On the Earth the main turn is about the $z$ axis (the pole axis): $x$ and $y$ change, $z$ \bnum{stays}.}
      \[ \mathbf{R}_3(\varphi) = \begin{pmatrix} \cos\varphi & \sin\varphi & 0 \\ -\sin\varphi & \cos\varphi & 0 \\ 0 & 0 & 1 \end{pmatrix} \]
      \tr{Há giros análogos em torno de $x$ ($\mathbf{R}_1$) e de $y$ ($\mathbf{R}_2$). Uma virada qualquer é uma \bnum{sequência} de giros --- e a sequência é uma \bnum{multiplicação de matrizes}.}{There are similar turns about $x$ ($\mathbf{R}_1$) and $y$ ($\mathbf{R}_2$). Any turn is a \bnum{sequence} of turns --- and the sequence is a \bnum{multiplication of matrices}.}
    \end{column}
    \begin{column}{0.42\textwidth}
      \begin{cuidado}[title={\tr{Experimente com um livro}{Try it with a book}}]
        \small\tr{Gire um livro $90\gr$ em torno do eixo \emph{vertical} e depois $90\gr$ em torno de um eixo \emph{horizontal}. Agora volte e faça \bnum{na ordem contrária}.}{Turn a book $90\gr$ about the \emph{vertical} axis and then $90\gr$ about a \emph{horizontal} axis. Now go back and do it \bnum{in the opposite order}.}
        \medskip

        \tr{O livro termina virado de \alert{jeitos diferentes}: a ordem dos giros importa. Por isso a ponte tem uma ordem fixa.}{The book ends up turned \alert{in different ways}: the order of the turns matters. That is why the bridge has a fixed order.}
      \end{cuidado}
    \end{column}
  \end{columns}
\end{frame}

\begin{frame}{\tr{A ponte tem três degraus: $W$, $R$ e $Q$}{The bridge has three steps: $W$, $R$ and $Q$}}
  \centering
  \begin{tikzpicture}[node distance=1.0cm, caixa/.append style={font=\small, text width=2.5cm, minimum height=1.1cm}]
    \node[caixa, draw=corFrame, fill=corFrame!10] (A) {\fram{ITRF}\\\tr{chão}{ground}};
    \node[caixa, draw=corFrame!60!corTempo, fill=corTempo!8, right=of A] (B) {TIRS\\\tr{chão, sem o balanço do polo}{ground, no pole wobble}};
    \node[caixa, draw=corTempo, fill=corTempo!10, right=of B] (C) {CIRS\\\tr{céu, com o eixo da Terra}{sky, with the Earth's axis}};
    \node[caixa, draw=corSist, fill=corSist!8, right=of C] (D) {\sist{ICRF}\\\tr{céu}{sky}};
    \draw[seta] (A) -- node[above, font=\large] {$W$} (B);
    \draw[seta] (B) -- node[above, font=\large] {$R$} (C);
    \draw[seta] (C) -- node[above, font=\large] {$Q$} (D);
    \node[font=\footnotesize, text=corFrame, align=center, text width=2.7cm, anchor=north] at ($(A.south)!0.5!(B.south)+(0,-0.12)$) {\tr{movimento polar: o eixo anda na crosta}{polar motion: the axis wanders in the crust}};
    \node[font=\footnotesize, text=corTempo, align=center, text width=2.7cm, anchor=north] at ($(B.south)!0.5!(C.south)+(0,-0.12)$) {\tr{rotação: 1 volta por dia sideral}{rotation: 1 turn per sidereal day}};
    \node[font=\footnotesize, text=corSist, align=center, text width=2.7cm, anchor=north] at ($(C.south)!0.5!(D.south)+(0,-0.12)$) {\tr{precessão e nutação: o eixo muda no céu}{precession and nutation: the axis moves in the sky}};
  \end{tikzpicture}
  \medskip

  \small
  \tr{\textbf{chão $\to$ céu:}}{\textbf{ground $\to$ sky:}} $\vec r_{\text{ICRF}} = Q(t)\;R(t)\;W(t)\;\vec r_{\text{ITRF}}$
  \hspace{1.2em}
  \tr{\textbf{céu $\to$ chão:} o contrário: giros desfeitos (a matriz \emph{transposta}), na ordem inversa.}{\textbf{sky $\to$ ground:} the opposite: turns undone (the \emph{transposed} matrix), in reverse order.}
  \medskip

  \begin{ideia}
    \footnotesize\tr{Cada degrau é uma matriz \bnum{pequena e conhecida}. O que muda no tempo são alguns ângulos --- e esses ângulos são os \bnum{EOP}.}{Each step is a \bnum{small, known} matrix. What changes in time are a few angles --- and those angles are the \bnum{EOP}.}
  \end{ideia}
\end{frame}

% ------------------------------------------------------------------ A rotação
\section{\tr{A rotação da Terra}{The Earth's rotation}}

\begin{frame}{\tr{Quanto dura um dia?}{How long is a day?}}
  \begin{columns}[T]
    \begin{column}{0.46\textwidth}
      \small
      \tr{O \bnum{dia solar}: o Sol volta ao mesmo lugar no céu. Dura \bnum{24 h}.}{The \bnum{solar day}: the Sun returns to the same place in the sky. It lasts \bnum{24 h}.}
      \medskip

      \tr{O \bnum{dia sideral}: uma estrela --- ou um quasar --- volta ao mesmo lugar. Dura \bnum{\mbox{23 h 56 min 4,09 s}}.}{The \bnum{sidereal day}: a star --- or a quasar --- returns to the same place. It lasts \bnum{\mbox{23 h 56 min 4.09 s}}.}
      \medskip

      \tr{Por quê? Enquanto gira, a Terra também \emph{anda} em volta do Sol, cerca de $1\gr$ por dia. Depois de uma volta completa, o Sol ainda está um pouco de lado: é preciso girar mais $\approx 1\gr$ para ele voltar ao mesmo lugar.}{Why? While it spins, the Earth also \emph{travels} around the Sun, about $1\gr$ per day. After one full turn, the Sun is still a bit to the side: it takes about $1\gr$ more to bring it back to the same place.}
    \end{column}
    \begin{column}{0.52\textwidth}
      \centering
      \begin{tikzpicture}[scale=0.8]
        \fill[corTempo!60] (0,0) circle (0.6);
        \draw[corTempo] (0,0) circle (0.6);
        \node[font=\footnotesize, text=corTempo] at (0,-0.95) {\tr{Sol}{Sun}};
        \draw[mStoneMid!60, dashed] (-18:3.6) arc (-18:18:3.6);
        \coordinate (E0) at (-16:3.6);
        \coordinate (E1) at (16:3.6);
        \fill[corCeu] (-3.07,0.89) circle (0.08);
        \node[font=\footnotesize, text=corCeu, align=center, anchor=south] at (-3.07,1.0) {\tr{estrela distante}{distant star}};
        \draw[corCeu, densely dotted, thick] (E0) -- (-3.07,0.89);
        \draw[corCeu, densely dotted, thick] (E1) -- ($(E1)+(164:6.0)$);
        \draw[mStoneMid!70, thick, fill=corFrame!8] (E0) circle (0.45);
        \draw[mStoneMid!70, thick, fill=corFrame!8] (E1) circle (0.45);
        \draw[corCeu, seta] (E0) -- ($(E0)+(164:1.15)$);
        \draw[corCeu, seta] (E1) -- ($(E1)+(164:1.15)$);
        \draw[corTempo, seta] (E1) -- ($(E1)+(196:1.15)$);
        \draw[corErro, thick] ($(E1)+(164:0.82)$) arc (164:196:0.82);
        \node[font=\footnotesize, text=corErro] at ($(E1)+(180:1.07)$) {$\delta$};
        \node[font=\footnotesize, anchor=north, align=center] at ($(E0)+(0,-0.55)$) {\tr{dia 0}{day 0}\\\tr{(meio-dia)}{(noon)}};
        \node[font=\footnotesize, anchor=south, align=center] at ($(E1)+(0,0.55)$) {\tr{uma volta}{one turn}\\\tr{depois}{later}};
      \end{tikzpicture}
      \par
      {\scriptsize\textcolor{mStoneMid}{\tr{Ângulo $\delta$ exagerado: na realidade são $\approx 0{,}9856\gr$.}{Angle $\delta$ exaggerated: in reality it is $\approx 0.9856\gr$.}}}
    \end{column}
  \end{columns}
\end{frame}

\begin{frame}{\tr{Dois dias, duas contas}{Two days, two counts}}
  \begin{columns}[T]
    \begin{column}{0.58\textwidth}
      \small
      \renewcommand{\arraystretch}{1.4}
      \begin{tabular}{@{}lll@{}}
        \toprule
        & \tr{o que volta}{what returns} & \tr{dura}{lasts} \\
        \midrule
        \tmpo{\tr{dia solar}{solar day}} & \tr{o Sol}{the Sun} & $24$ h $= \num{86400}$ s \\
        \tmpo{\tr{dia sideral}{sidereal day}} & \tr{uma estrela, um quasar}{a star, a quasar} & $\num{86164.09}$ s \\
        \tr{diferença}{difference} & & \bnum{3 min \num{55.9} s} \\
        \bottomrule
      \end{tabular}
      \medskip

      \tr{Em um ano há $\num{365.2422}$ dias solares, mas $\num{366.2422}$ voltas em relação às estrelas: a Terra dá \bnum{uma volta a mais} do que o número de dias.}{In one year there are $\num{365.2422}$ solar days, but $\num{366.2422}$ turns relative to the stars: the Earth makes \bnum{one more turn} than the number of days.}
    \end{column}
    \begin{column}{0.38\textwidth}
      \begin{cuidado}[title={\tr{Teste em casa}{Test it at home}}]
        \small\tr{Escolha uma estrela brilhante atrás de um poste e anote a hora. No dia seguinte, ela aparece \bnum{4 minutos mais cedo}; em uma semana, $\approx$ 27 minutos mais cedo.}{Pick a bright star behind a pole and note the time. The next day it shows up \bnum{4 minutes earlier}; a week later, $\approx$ 27 minutes earlier.}
      \end{cuidado}
    \end{column}
  \end{columns}
\end{frame}

\begin{frame}{\tr{O ângulo de rotação da Terra (ERA)}{The Earth Rotation Angle (ERA)}}
  \begin{columns}[T]
    \begin{column}{0.56\textwidth}
      \small
      \tr{O \bnum{ERA} (\emph{Earth Rotation Angle}) é o ``ponteiro'' que diz quanto a Terra girou em relação ao céu parado: o ângulo, em torno do eixo, entre uma origem \emph{no céu} (o CIO) e uma origem \emph{no chão} (o TIO).}{The \bnum{ERA} (\emph{Earth Rotation Angle}) is the ``hand'' that tells how far the Earth has turned relative to the still sky: the angle, about the axis, between an origin \emph{in the sky} (the CIO) and an origin \emph{on the ground} (the TIO).}
      \[ \begin{aligned} \text{ERA} = 360\gr\,\bigl(&\num{0.7790572732640} \\ &+ \num{1.00273781191135448}\,T_u\bigr) \end{aligned} \]
      {\footnotesize $T_u = \text{JD}(\text{UT1}) - \num{2451545.0}$ \tr{(dias desde J2000.0)}{(days since J2000.0)}}
      \medskip

      \begin{itemize}
        \setlength{\itemsep}{2pt}
        \item \tr{em 1/1/2000, 12 h UT1: $\num{280.4606}\gr$;}{on 1 Jan 2000, 12 h UT1: $\num{280.4606}\gr$;}
        \item \tr{cresce $\num{360.9856}\gr$ por dia solar;}{grows $\num{360.9856}\gr$ per solar day;}
        \item \tr{em 18/10/2026, 0 h UT1: $\num{26.155}\gr$.}{on 18 Oct 2026, 0 h UT1: $\num{26.155}\gr$.}
      \end{itemize}
    \end{column}
    \begin{column}{0.40\textwidth}
      \centering
      \begin{tikzpicture}[scale=1.2]
        \draw[mStoneMid!70, thick, fill=corFrame!6] (0,0) circle (1.3);
        \foreach \a in {0,30,...,330} { \draw[mStoneMid!70] (\a:1.18) -- (\a:1.3); }
        \draw[corCeu, seta, line width=1.3pt] (0,0) -- (1.6,0);
        \node[rotulo, text=corCeu, right] at (1.6,0) {CIO};
        \draw[corFrame, seta, line width=1.3pt] (0,0) -- (52:1.6);
        \node[rotulo, text=corFrame, above right] at (52:1.5) {TIO};
        \draw[corErro, thick] (0.6,0) arc (0:52:0.6);
        \node[rotulo, text=corErro] at (26:0.88) {ERA};
        \draw[corTempo, -{Stealth[length=2mm]}] (1.55,0.95) arc (32:75:1.8);
        \node[rotulo, text=corTempo, align=center] at (0,-1.7) {\tr{visto de cima do polo norte;}{seen from above the north pole;}\\\tr{a Terra gira no sentido anti-horário}{the Earth turns counterclockwise}};
      \end{tikzpicture}
    \end{column}
  \end{columns}
\end{frame}

\begin{frame}{\tr{Sua vez}{Your turn}}
  \begin{columns}[T]
    \begin{column}{0.54\textwidth}
      \begin{regra}[title={\tr{Dia solar, dia sideral e o ERA}{Solar day, sidereal day and the ERA}}]
        \small
        \begin{enumerate}
          \setlength{\itemsep}{4pt}
          \item[(a)] \tr{Uma estrela passa atrás de um poste às 21:00 hoje. A que horas passa amanhã? E daqui a 30 dias?}{A star passes behind a pole at 21:00 today. When does it pass tomorrow? And 30 days from now?}
          \item[(b)] \tr{Quantos graus a Terra gira em 6 h?}{How many degrees does the Earth turn in 6 h?}
          \item[(c)] \tr{Um satélite ``fixo'' no céu, em $(26\,560,\,0)$ km. Depois de 6 h, onde está nos eixos do chão? (use $90\gr$)}{A satellite ``fixed'' in the sky, at $(26\,560,\,0)$ km. After 6 h, where is it in the ground axes? (use $90\gr$)}
        \end{enumerate}
      \end{regra}
    \end{column}
    \begin{column}{0.42\textwidth}
      \onslide<2->{
      \begin{ideia}[title={\tr{Respostas}{Answers}}]
        \small\raggedright
        (a) \tr{$\approx$ 20:56 amanhã ($-3{,}9$ min por dia); em 30 dias, $30\times3{,}93\approx118$ min mais cedo: $\approx$ \bnum{19:02}.}{$\approx$ 20:56 tomorrow ($-3.9$ min per day); in 30 days, $30\times3.93\approx118$ min earlier: $\approx$ \bnum{19:02}.}
        \medskip

        (b) $6\times\num{15.041}\gr \approx \bnum{\num{90.2}\gr}$.
        \medskip

        (c) $(0,\,-26\,560)$ km: \tr{o ponto do chão sob ele andou $90\gr$ para o oeste.}{the ground point under it moved $90\gr$ west.}
      \end{ideia}
      }
    \end{column}
  \end{columns}
\end{frame}

\begin{frame}{\tr{A Terra não é um relógio perfeito}{The Earth is not a perfect clock}}
  \begin{columns}[T]
    \begin{column}{0.54\textwidth}
      \small
      \tr{Uma patinadora gira \bnum{mais depressa} ao fechar os braços: o giro total (o \emph{momento angular}) se conserva. Ela mudou a distribuição de massa, e mudou o ritmo.}{A skater spins \bnum{faster} when she pulls her arms in: the total spin (the \emph{angular momentum}) is conserved. She changed the mass distribution, and changed the pace.}
      \medskip

      \tr{A Terra faz o mesmo, em pequeno: ventos, oceanos, gelo, o núcleo líquido --- tudo que \emph{muda de lugar} muda um pouquinho o ritmo do giro.}{The Earth does the same, on a small scale: winds, oceans, ice, the liquid core --- everything that \emph{moves} changes the pace of the spin a little.}
      \medskip

      \tr{Por isso o ângulo $\text{ERA}$ não é uma conta de calendário: o tempo da rotação (\bnum{UT1}) tem de ser \bnum{medido} todo dia.}{That is why the $\text{ERA}$ is not a calendar calculation: the time of the rotation (\bnum{UT1}) has to be \bnum{measured} every day.}
    \end{column}
    \begin{column}{0.42\textwidth}
      \centering
      \begin{tikzpicture}[scale=1.25]
        % braços abertos, com pesos nas mãos
        \fill[mStone] (0,1.6) circle (0.12);
        \draw[line width=1.2pt, mStone] (0,1.46) -- (0,0.72);
        \draw[line width=1.2pt, mStone] (0,0.72) -- (-0.17,0) (0,0.72) -- (0.17,0);
        \draw[line width=1.2pt, mStone] (-0.85,1.22) -- (0.85,1.22);
        \fill[corTempo] (-0.85,1.22) circle (0.07) (0.85,1.22) circle (0.07);
        \draw[corTempo, -{Stealth[length=2mm]}, thick] (-0.7,0.4) arc (200:340:0.7 and 0.16);
        \node[font=\footnotesize, text=corTempo] at (0,-0.3) {\tr{devagar}{slow}};
        % braços junto ao corpo
        \begin{scope}[shift={(2.5,0)}]
          \fill[mStone] (0,1.6) circle (0.12);
          \draw[line width=1.2pt, mStone] (0,1.46) -- (0,0.72);
          \draw[line width=1.2pt, mStone] (0,0.72) -- (-0.17,0) (0,0.72) -- (0.17,0);
          \draw[line width=1.2pt, mStone] (0,1.3) -- (-0.2,1.08) -- (-0.04,0.92) (0,1.3) -- (0.2,1.08) -- (0.04,0.92);
          \fill[corErro] (-0.04,0.92) circle (0.07) (0.04,0.92) circle (0.07);
          \draw[corErro, -{Stealth[length=2mm]}, thick] (-0.5,0.5) arc (200:340:0.5 and 0.12);
          \draw[corErro, -{Stealth[length=2mm]}, thick] (-0.5,0.3) arc (200:340:0.5 and 0.12);
          \draw[corErro, -{Stealth[length=2mm]}, thick] (-0.5,0.1) arc (200:340:0.5 and 0.12);
          \node[font=\footnotesize, text=corErro] at (0,-0.3) {\tr{depressa}{fast}};
        \end{scope}
      \end{tikzpicture}
      \par\smallskip
      \begin{ideia}
        \small $L = I\,\omega$ \tr{constante}{constant}\\[2pt]
        $I \downarrow \;\Rightarrow\; \omega \uparrow$
      \end{ideia}
    \end{column}
  \end{columns}
\end{frame}

\begin{frame}{\tr{O que faz a rotação variar}{What makes the rotation vary}}
  \begin{columns}[T]
    \begin{column}{0.62\textwidth}
      \footnotesize
      \renewcommand{\arraystretch}{1.5}
      \begin{tabular}{@{}>{\raggedright\arraybackslash}p{2.6cm}>{\raggedright\arraybackslash}p{2.0cm}>{\raggedright\arraybackslash}p{3.2cm}@{}}
        \toprule
        \tr{causa}{cause} & \tr{escala de tempo}{time scale} & \tr{efeito no dia}{effect on the day} \\
        \midrule
        \tr{marés (a Lua freia a Terra)}{tides (the Moon brakes the Earth)} & \tr{séculos}{centuries} & \tr{alonga $\approx$ \bnum{1,8 ms por século}}{lengthens $\approx$ \bnum{1.8 ms per century}} \\
        \tr{ventos e oceanos}{winds and oceans} & \tr{dias a anos}{days to years} & \tr{da ordem de \bnum{1 ms} ao longo do ano}{of the order of \bnum{1 ms} through the year} \\
        \tr{núcleo líquido $\leftrightarrow$ manto}{liquid core $\leftrightarrow$ mantle} & \tr{décadas}{decades} & \tr{alguns \bnum{ms}}{a few \bnum{ms}} \\
        \tr{terremotos gigantes}{giant earthquakes} & \tr{instantâneo}{instant} & \tr{alguns \bnum{$\mu$s} (Japão, 2011: $\approx 1{,}8\ \mu$s)}{a few \bnum{$\mu$s} (Japan, 2011: $\approx 1.8\ \mu$s)} \\
        \bottomrule
      \end{tabular}
    \end{column}
    \begin{column}{0.34\textwidth}
      \begin{cuidado}[title={\tr{Ordem de grandeza}{Order of magnitude}}]
        \small\tr{Um dia tem \num{86400} s. Variações de 1 ms são $\approx$ $10^{-8}$ do dia --- pequenas, mas \alert{grandes para a geodésia}, como mostra o próximo slide.}{A day has \num{86400} s. Variations of 1 ms are $\approx$ $10^{-8}$ of the day --- small, but \alert{large for geodesy}, as the next slide shows.}
      \end{cuidado}
    \end{column}
  \end{columns}
\end{frame}

\begin{frame}{\tr{Quanto vale 1 milissegundo?}{What is 1 millisecond worth?}}
  \begin{columns}[T]
    \begin{column}{0.50\textwidth}
      \small
      \tr{A Terra gira $\approx 15{,}04\asec$ por segundo, isto é, \bnum{$15{,}04\mas$ por milissegundo}.}{The Earth turns $\approx 15.04\asec$ per second, that is, \bnum{$15.04\mas$ per millisecond}.}
      \medskip

      \tr{Se o relógio da rotação (UT1) errar $1$ ms, é como se a Terra estivesse virada $15\mas$ a mais ou a menos:}{If the rotation clock (UT1) is off by $1$ ms, it is as if the Earth were turned $15\mas$ too far or too little:}
      \[ \text{erro} = \omega \, r \, \Delta t \]
      \tr{com $\omega = \num{7.292115e-5}$ rad/s, $r$ a distância ao eixo e $\Delta t = 1$ ms.}{with $\omega = \num{7.292115e-5}$ rad/s, $r$ the distance to the axis and $\Delta t = 1$ ms.}
    \end{column}
    \begin{column}{0.46\textwidth}
      \centering
      \begin{tikzpicture}[x=1.0cm, y=1.0cm]
        \foreach \yy/\len/\cor/\lab in {2.4/0.465/corFrame/{\tr{Equador}{Equator}: \num{0.465} m}, 1.2/0.420/corFrame/{Curitiba: \num{0.420} m}, 0/1.937/corErro/{\tr{órbita do GPS}{GPS orbit}: \num{1.94} m}} {
          \fill[\cor!80] (0,\yy) rectangle ({\len*2.0},{\yy+0.4});
          \node[rotulo, anchor=south west] at (0,{\yy+0.4}) {\lab};
        }
        \draw[mStoneMid] (0,-0.2) -- (0,3.1);
      \end{tikzpicture}
      \par\smallskip
      {\scriptsize\textcolor{mStoneMid}{\tr{Um erro de 1 ms em UT1, em metros (a posição ``gira'' errada).}{A 1 ms error in UT1, in meters (the position is ``turned'' wrongly).}}}
    \end{column}
  \end{columns}
\end{frame}

% ------------------------------------------------------------------ O eixo balança no espaço
\section{\tr{O eixo balança no espaço}{The axis wobbles in space}}

\begin{frame}{\tr{O eixo da Terra, no céu: o CIP}{The Earth's axis, in the sky: the CIP}}
  \begin{columns}[T]
    \begin{column}{0.46\textwidth}
      \small
      \tr{A cada instante o eixo da Terra aponta para um ponto do céu: o \bnum{polo celeste intermediário} (\bnum{CIP}).}{At each instant the Earth's axis points to a point in the sky: the \bnum{Celestial Intermediate Pole} (\bnum{CIP}).}
      \medskip

      \tr{Ele anda por duas razões (aula 1):}{It moves for two reasons (class 1):}
      \begin{itemize}
        \setlength{\itemsep}{1pt}
        \item \tr{\bnum{precessão}: $\approx 20\asec$/ano na posição do polo;}{\bnum{precession}: $\approx 20\asec$/yr in the position of the pole;}
        \item \tr{\bnum{nutação}: um tremor de $9{,}2\asec \times 6{,}8\asec$, em ciclos de $18{,}6$ anos.}{\bnum{nutation}: a wobble of $9.2\asec \times 6.8\asec$, in $18.6$-year cycles.}
      \end{itemize}
      \medskip

      \tr{A posição do CIP é dada por duas coordenadas pequenas, $X$ e $Y$, contadas a partir do polo do ICRS (em 2026, $X \approx 537\asec$).}{The CIP position is given by two small coordinates, $X$ and $Y$, counted from the ICRS pole (in 2026, $X \approx 537\asec$).}
    \end{column}
    \begin{column}{0.52\textwidth}
      \centering
      \begin{tikzpicture}[x=0.0165cm, y=0.115cm]
        \draw[corTempo!70, thick, dashed] (-10,0) -- (385,0);
        \draw[corFrame, thick, domain=0:18.6, samples=240, variable=\t]
          plot ({20*\t + 6.84*sin(deg(2*pi*\t/18.6))},{9.2*cos(deg(2*pi*\t/18.6))});
        \foreach \t in {0,2,...,18} {
          \node[circle, fill=corFrame, inner sep=1.1pt] at ({20*\t + 6.84*sin(deg(2*pi*\t/18.6))},{9.2*cos(deg(2*pi*\t/18.6))}) {};
        }
        \node[circle, fill=corFrame, inner sep=1.1pt] at (372,9.2) {};
        \node[rotulo, text=corFrame, anchor=south] at (0,12) {\tr{início}{start}};
        \node[rotulo, text=corFrame, anchor=south east] at (385,12) {\tr{18,6 anos depois}{18.6 years later}};
        \node[rotulo, text=corTempo, anchor=north] at (190,-12) {\tr{polo médio (só precessão): $20\asec$/ano}{mean pole (precession only): $20\asec$/yr}};
        \draw[mStoneMid] (0,-24) -- (100,-24);
        \draw[mStoneMid] (0,-22.5) -- (0,-25.5) (100,-22.5) -- (100,-25.5);
        \node[rotulo, anchor=west] at (108,-24) {$100\asec$};
        \draw[mStoneMid] (-18,-9.2) -- (-18,0);
        \draw[mStoneMid] (-20,-9.2) -- (-16,-9.2) (-20,0) -- (-16,0);
        \node[rotulo, anchor=east] at (-22,-4.6) {$\num{9.2}\asec$};
      \end{tikzpicture}
      \par\smallskip
      {\scriptsize\textcolor{mStoneMid}{\tr{Pontos a cada 2 anos. Ilustrativo: fase arbitrária; vertical ampliada $\approx$ 7 vezes.}{Points every 2 years. Illustrative: arbitrary phase; vertical exaggerated $\approx$ 7 times.}}}
    \end{column}
  \end{columns}
\end{frame}

\begin{frame}{\tr{Modelo mais medida: $dX$ e $dY$}{Model plus measurement: $dX$ and $dY$}}
  \begin{columns}[T]
    \begin{column}{0.50\textwidth}
      \small
      \tr{O modelo da IAU (precessão IAU 2006 e nutação IAU 2000A, com \bnum{1\,365 termos}) calcula $X$ e $Y$.}{The IAU model (IAU 2006 precession and IAU 2000A nutation, with \bnum{1,365 terms}) computes $X$ and $Y$.}
      \medskip

      \tr{Mas o núcleo líquido da Terra tem uma oscilação própria (\emph{free core nutation}, $\approx$ 430 dias) que o modelo \alert{não prevê}. O polo real difere do modelo por até $\approx 0{,}3\mas$.}{But the Earth's liquid core has a natural oscillation (\emph{free core nutation}, $\approx$ 430 days) that the model \alert{does not predict}. The real pole differs from the model by up to $\approx 0.3\mas$.}
      \medskip

      \tr{O VLBI mede a diferença: são os \bnum{$dX$ e $dY$}, publicados pelo IERS. $0{,}3\mas \approx 9$ mm no chão.}{VLBI measures the difference: these are \bnum{$dX$ and $dY$}, published by the IERS. $0.3\mas \approx 9$ mm on the ground.}
    \end{column}
    \begin{column}{0.46\textwidth}
      \centering
      \begin{tikzpicture}[x=0.52cm, y=3.2cm]
        \draw[mStoneMid] (0,0) -- (9.3,0);
        \draw[mStoneMid] (0,-0.42) -- (0,0.42);
        \foreach \xx/\lab in {1/2020, 3/2022, 5/2024, 7/2026} {
          \draw[mStoneMid] (\xx,-0.015) -- (\xx,0.015);
          \node[rotulo, below] at (\xx,-0.33) {\lab};
        }
        \draw[corFrame, thick, domain=0:9, samples=200, variable=\t]
          plot ({\t},{(0.17+0.09*sin(deg(2*pi*\t/9)))*sin(deg(2*pi*\t/1.18+0.4))});
        \foreach \yy/\lab in {0.3/{$+\num{0.3}\mas$}, -0.3/{$-\num{0.3}\mas$}} {
          \draw[mStoneMid, dotted] (0,\yy) -- (9,\yy);
          \node[rotulo, anchor=east] at (0,\yy) {\lab};
        }
        \node[rotulo, text=corFrame, anchor=south] at (4.5,0.34) {$dX$ \tr{(exemplo ilustrativo)}{(illustrative example)}};
      \end{tikzpicture}
      \par\smallskip
      {\scriptsize\textcolor{mStoneMid}{\tr{Oscilação de $\approx 1{,}2$ ano e amplitude variável --- por isso só dá para \emph{medir}, não calcular.}{Oscillation of $\approx 1.2$ years with variable amplitude --- so it can only be \emph{measured}, not computed.}}}
    \end{column}
  \end{columns}
\end{frame}

\begin{frame}{\tr{Sua vez}{Your turn}}
  \begin{columns}[T]
    \begin{column}{0.54\textwidth}
      \begin{regra}[title={\tr{O tamanho da nutação}{The size of nutation}}]
        \small
        \begin{enumerate}
          \setlength{\itemsep}{4pt}
          \item[(a)] \tr{A nutação principal tem $9{,}2\asec$ de amplitude. Quantos metros no chão? (use $1\asec = 30{,}9$ m)}{The main nutation has an amplitude of $9.2\asec$. How many meters on the ground? (use $1\asec = 30.9$ m)}
          \item[(b)] \tr{O polo médio avança $20\asec$ por ano. Em $18{,}6$ anos, quantos quilômetros no chão?}{The mean pole advances $20\asec$ per year. In $18.6$ years, how many kilometers on the ground?}
          \item[(c)] \tr{Um erro $dX = 0{,}3\mas$ vale quantos milímetros? (use $1\mas = 30{,}9$ mm)}{An error $dX = 0.3\mas$ is worth how many millimeters? (use $1\mas = 30.9$ mm)}
        \end{enumerate}
      \end{regra}
    \end{column}
    \begin{column}{0.42\textwidth}
      \onslide<2->{
      \begin{ideia}[title={\tr{Respostas}{Answers}}]
        \small\raggedright
        (a) $\num{9.2} \times \num{30.9} \approx \bnum{284}$ m.
        \medskip

        (b) $20 \times \num{18.6} = 372\asec \Rightarrow 372 \times \num{30.9} \approx \bnum{\num{11.5}}$ km.
        \medskip

        (c) $\num{0.3} \times \num{30.9} \approx \bnum{\num{9.3}}$ mm.
      \end{ideia}
      }
    \end{column}
  \end{columns}
\end{frame}

% ------------------------------------------------------------------ O polo anda dentro da Terra
\section{\tr{O polo anda dentro da Terra}{The pole moves inside the Earth}}

\begin{frame}{\tr{O eixo também anda dentro da Terra}{The axis also moves inside the Earth}}
  \begin{columns}[T]
    \begin{column}{0.46\textwidth}
      \small
      \tr{O eixo de giro não está preso à crosta: a crosta ``escorrega'' em volta dele. O polo descreve uma curva de até $\approx 0{,}3\asec$, uns \bnum{9 m}, em torno de um ponto de referência (o \bnum{IRP}, o polo de referência do IERS).}{The rotation axis is not fixed to the crust: the crust ``slips'' around it. The pole traces a curve of up to $\approx 0.3\asec$, about \bnum{9 m}, around a reference point (the \bnum{IRP}, the IERS Reference Pole).}
      \medskip

      \tr{É o \bnum{movimento polar}, descrito por duas coordenadas, $x_p$ e $y_p$ (na direção de Greenwich e na de $90\gr$ oeste).}{This is the \bnum{polar motion}, described by two coordinates, $x_p$ and $y_p$ (toward Greenwich and toward $90\gr$ west).}
      \medskip

      \tr{Foi notado primeiro como uma pequena variação da \emph{latitude} das estações, no fim do século XIX.}{It was first noticed as a small variation of the \emph{latitude} of stations, in the late 19th century.}
    \end{column}
    \begin{column}{0.52\textwidth}
      \centering
      \begin{tikzpicture}[scale=1.0]
        \draw[mStoneMid!70, thick, fill=corFrame!8] (0,0) circle (1.25);
        \draw[corTempo, dashed, thick] (0,-1.7) -- (0,1.9);
        \draw[corFrame, thick, rotate=8] (0,-1.5) -- (0,1.5);
        \node[font=\footnotesize, text=corTempo, anchor=south] at (0,1.9) {\tr{eixo de rotação}{rotation axis}};
        \node[font=\footnotesize, text=corFrame, anchor=north] at (0.1,-1.72) {\tr{eixo da crosta}{crust axis}};
        \draw[mStoneMid, thin] (-0.19,1.485) -- (2.5,1.15);
        \draw[mStoneMid, thin] (0,1.9) -- (2.5,2.4);
        \draw[mStoneMid!70, fill=white] (3.7,1.7) circle (1.2);
        \draw[mStoneMid!50, dotted] (3.7,1.7) circle (0.62);
        \draw[mStoneMid, -{Stealth[length=1.6mm]}] (2.6,1.7) -- (5.0,1.7);
        \draw[mStoneMid, -{Stealth[length=1.6mm]}] (3.7,0.62) -- (3.7,2.95);
        \node[font=\footnotesize, anchor=west] at (5.0,1.7) {$x_p$};
        \node[font=\footnotesize, anchor=south] at (3.7,2.95) {$y_p$};
        \draw[corFrame, thick, domain=0:15, samples=300, variable=\u]
          plot ({3.7+0.34*cos(deg(\u))+0.23*cos(deg(1.1855*\u))},{1.7+0.34*sin(deg(\u))+0.23*sin(deg(1.1855*\u))});
        \fill[corErro] (3.7,1.7) circle (0.05);
        \node[font=\footnotesize, text=corErro, anchor=north west] at (3.76,1.64) {IRP};
        \node[font=\footnotesize, text=mStoneMid] at (3.7,0.2) {\tr{ampliação enorme}{huge magnification}};
      \end{tikzpicture}
      \par\smallskip
      {\scriptsize\textcolor{mStoneMid}{\tr{Fora de escala: o ângulo real entre os eixos é de $<0{,}3\asec$, invisível em desenho.}{Not to scale: the real angle between the axes is $<0.3\asec$, invisible in a drawing.}}}
    \end{column}
  \end{columns}
\end{frame}

\begin{frame}{\tr{A espiral do polo}{The pole's spiral}}
  \begin{columns}[T]
    \begin{column}{0.46\textwidth}
      \small
      \tr{Ao longo dos anos o polo desenha uma \bnum{espiral}: abre e fecha a cada $\approx 6{,}4$ anos. Entre as voltas mais abertas e as mais fechadas o raio varia de $\approx 60$ a $\approx 240\mas$.}{Over the years the pole draws a \bnum{spiral}: it opens and closes every $\approx 6.4$ years. Between the widest and the tightest turns the radius varies from $\approx 60$ to $\approx 240\mas$.}
      \medskip

      \tr{Além disso, o polo médio deriva devagar, alguns $\mas$ por ano (centímetros por ano), por causa do rebote pós-glacial e do derretimento do gelo.}{In addition, the mean pole drifts slowly, a few $\mas$ per year (centimeters per year), because of post-glacial rebound and melting ice.}
      \medskip

      \begin{ideia}[title={\tr{Quanto é isso?}{How much is that?}}]
        \footnotesize\tr{$100\mas \approx 3{,}1$ m no chão. Os pontos mais afastados estão a $\approx 7$ m do IRP.}{$100\mas \approx 3.1$ m on the ground. The farthest points are $\approx 7$ m from the IRP.}
      \end{ideia}
    \end{column}
    \begin{column}{0.52\textwidth}
      \centering
      \begin{tikzpicture}[x=0.0085cm, y=0.0085cm]
        \draw[mStoneMid!40, dotted] (0,0) circle (100);
        \draw[mStoneMid!40, dotted] (0,0) circle (200);
        \draw[mStoneMid, -{Stealth[length=1.6mm]}] (-260,0) -- (285,0);
        \draw[mStoneMid, -{Stealth[length=1.6mm]}] (0,-260) -- (0,275);
        \node[rotulo, anchor=west] at (285,0) {$x_p$};
        \node[rotulo, anchor=south] at (0,275) {$y_p$};
        \draw[corFrame, thick, domain=0:12.8, samples=900, variable=\t]
          plot ({150*cos(deg(2*pi*\t/1.1855)) + 90*cos(deg(2*pi*\t))},{150*sin(deg(2*pi*\t/1.1855)) + 90*sin(deg(2*pi*\t))});
        \fill[corErro] (240,0) circle (6);
        \node[rotulo, text=corErro, anchor=south west] at (248,8) {\tr{início}{start}};
        \fill[mStone] (0,0) circle (5);
        \node[rotulo, anchor=north east] at (-8,-8) {IRP};
        \node[rotulo, text=mStoneMid, anchor=north west] at (-262,-266) {$100\mas$, $200\mas$: \tr{círculos pontilhados}{dotted circles}};
      \end{tikzpicture}
      \par\smallskip
      {\scriptsize\textcolor{mStoneMid}{\tr{12,8 anos (duas aberturas). Curva ilustrativa, gerada por um modelo simples --- não são dados reais.}{12.8 years (two openings). Illustrative curve from a simple model --- not real data.}}}
    \end{column}
  \end{columns}
\end{frame}

\begin{frame}{\tr{Duas ondas que se somam}{Two waves that add up}}
  \begin{columns}[T]
    \begin{column}{0.50\textwidth}
      \footnotesize
      \tr{\bnum{Chandler} ($\approx$ 433 dias): oscilação \emph{livre}. Uma Terra rígida oscilaria em 305 dias (Euler); como a Terra é elástica, o período se alonga. Descoberta por Seth Chandler, em 1891.}{\bnum{Chandler} ($\approx$ 433 days): a \emph{free} oscillation. A rigid Earth would oscillate in 305 days (Euler); since the Earth is elastic, the period lengthens. Discovered by Seth Chandler in 1891.}
      \medskip

      \tr{\bnum{Anual} (365 dias): oscilação \emph{forçada} pelas estações (ar, água, neve movendo massa).}{\bnum{Annual} (365 days): an oscillation \emph{forced} by the seasons (air, water, snow moving mass).}
      \medskip

      \tr{Duas ondas de períodos parecidos \emph{batem}, como duas notas quase iguais num violão:}{Two waves with similar periods \emph{beat}, like two nearly equal notes on a guitar:}
      \[ \begin{aligned} T_{\text{bat}} &= \Bigl(\tfrac{1}{\num{365.25}} - \tfrac{1}{433}\Bigr)^{-1} \\ &\approx \bnum{2\,334}\ \text{\tr{dias}{days}} \approx \bnum{\num{6.4}}\ \text{\tr{anos}{years}} \end{aligned} \]
    \end{column}
    \begin{column}{0.46\textwidth}
      \centering
      \begin{tikzpicture}[x=0.40cm, y=0.06cm]
        \draw[mStoneMid] (0,0) -- (14.3,0);
        \foreach \t in {0,2,...,14} {
          \draw[mStoneMid] (\t,-0.9) -- (\t,0.9);
          \node[rotulo, below] at (\t,-27) {\t};
        }
        \draw[corErro, dashed, domain=0:14, samples=200, variable=\t]
          plot ({\t},{sqrt(15^2+9^2+2*15*9*cos(deg(2*pi*0.1565*\t)))});
        \draw[corErro, dashed, domain=0:14, samples=200, variable=\t]
          plot ({\t},{-sqrt(15^2+9^2+2*15*9*cos(deg(2*pi*0.1565*\t)))});
        \draw[corFrame, thick, domain=0:14, samples=700, variable=\t]
          plot ({\t},{15*cos(deg(2*pi*\t/1.1855)) + 9*cos(deg(2*pi*\t))});
        \draw[seta2, mStone] (0,29) -- (6.39,29);
        \node[rotulo, above] at (3.2,29) {$\approx \num{6.4}$ \tr{anos}{years}};
        \node[rotulo, anchor=north] at (7,-36) {\tr{tempo (anos)}{time (years)}};
        \draw[mStoneMid] (-0.2,-24) -- (-0.2,24);
        \node[rotulo, anchor=east] at (-0.3,0) {$x_p$};
      \end{tikzpicture}
      \par\smallskip
      {\scriptsize\textcolor{mStoneMid}{\tr{Soma de uma onda de 433 dias (150 mas) e uma de 365 dias (90 mas): ilustrativo. Em tracejado, o envelope.}{Sum of a 433-day wave (150 mas) and a 365-day wave (90 mas): illustrative. Dashed: the envelope.}}}
    \end{column}
  \end{columns}
\end{frame}

\begin{frame}{\tr{Sua vez}{Your turn}}
  \begin{columns}[T]
    \begin{column}{0.54\textwidth}
      \begin{regra}[title={\tr{O tamanho do polo}{The size of the pole}}]
        \small
        \begin{enumerate}
          \setlength{\itemsep}{4pt}
          \item[(a)] \tr{O polo anda $0{,}3\asec$. Quantos metros? (use $1\asec = 30{,}9$ m)}{The pole moves $0.3\asec$. How many meters? (use $1\asec = 30.9$ m)}
          \item[(b)] \tr{Quantos centímetros valem $1\mas$ no chão?}{How many centimeters is $1\mas$ worth on the ground?}
          \item[(c)] \tr{Duas ondas de 365 e de 433 dias: de quanto em quanto tempo elas voltam a ``entrar em fase''? (dica: $1/365 - 1/433$)}{Two waves of 365 and 433 days: how often do they come back ``in phase''? (hint: $1/365 - 1/433$)}
        \end{enumerate}
      \end{regra}
    \end{column}
    \begin{column}{0.42\textwidth}
      \onslide<2->{
      \begin{ideia}[title={\tr{Respostas}{Answers}}]
        \small\raggedright
        (a) $\num{0.3} \times \num{30.9} \approx \bnum{\num{9.3}}$ m.
        \medskip

        (b) $\num{0.001} \times \num{30.9}\ \text{m} \approx \bnum{\num{3.1}}$ cm.
        \medskip

        (c) $\dfrac{1}{1/365 - 1/433} \approx \bnum{2\,300}$ \tr{dias}{days} $\approx \num{6.3}$ \tr{anos}{years}.
      \end{ideia}
      }
    \end{column}
  \end{columns}
\end{frame}

% ------------------------------------------------------------------ Os EOP
\section{\tr{Os EOP: quem mede e quem publica}{The EOP: who measures and who publishes}}

\begin{frame}{\tr{Cinco números por dia}{Five numbers per day}}
  \centering
  \small
  \renewcommand{\arraystretch}{1.3}
  \begin{tabular}{@{}>{\raggedright\arraybackslash}p{2.0cm}>{\raggedright\arraybackslash}p{5.6cm}>{\raggedright\arraybackslash}p{3.3cm}c@{}}
    \toprule
    \textbf{EOP} & \tr{o que diz}{what it says} & \tr{tamanho típico}{typical size} & \tr{degrau}{step} \\
    \midrule
    $x_p,\ y_p$ & \tr{onde o polo está na crosta}{where the pole is in the crust} & $\lesssim \num{0.3}\asec$ ($\approx$ 9 m) & $W$ \\
    UT1$-$UTC & \tr{quanto a Terra girou, comparado ao relógio atômico}{how far the Earth has turned, compared with the atomic clock} & \tr{menos de $\num{0.9}$ s (regra atual)}{under $\num{0.9}$ s (current rule)} & $R$ \\
    $dX,\ dY$ & \tr{correção da posição do eixo no céu}{correction to the axis position in the sky} & $\approx \num{0.1}$--$\num{0.3}\mas$ & $Q$ \\
    \bottomrule
  \end{tabular}
  \medskip

  \tr{Cinco números independentes. Um sexto, o \bnum{LOD} (a duração do dia menos $\num{86400}$ s), é só o \emph{ritmo} de UT1: $\text{LOD} = -\,\mathrm{d}(\text{UT1}-\text{UTC})/\mathrm{d}t$.}{Five independent numbers. A sixth, the \bnum{LOD} (the length of day minus $\num{86400}$ s), is just the \emph{rate} of UT1: $\text{LOD} = -\,\mathrm{d}(\text{UT1}-\text{UTC})/\mathrm{d}t$.}
  \medskip

  \begin{ideia}
    \small\tr{São os \bnum{EOP}: \emph{Earth Orientation Parameters}. Com eles, a ponte $Q\,R\,W$ fica completa.}{These are the \bnum{EOP}: \emph{Earth Orientation Parameters}. With them, the $Q\,R\,W$ bridge is complete.}
  \end{ideia}
\end{frame}

\begin{frame}{\tr{Quem mede o quê?}{Who measures what?}}
  \begin{columns}[T]
    \begin{column}{0.62\textwidth}
      \footnotesize
      \renewcommand{\arraystretch}{1.5}
      \setlength{\tabcolsep}{4pt}
      \begin{tabular}{@{}l>{\centering\arraybackslash}p{1.15cm}>{\centering\arraybackslash}p{1.45cm}>{\centering\arraybackslash}p{1.15cm}>{\centering\arraybackslash}p{1.15cm}@{}}
        \toprule
        & VLBI & GNSS & SLR & DORIS \\
        \midrule
        $x_p,\ y_p$ & $\checkmark$ & $\checkmark$ & $\checkmark$ & $\checkmark$ \\
        LOD & $\checkmark$ & $\checkmark$ & $\checkmark$ & $\checkmark$ \\
        UT1$-$UTC & \textbf{\textcolor{mTeal}{$\checkmark$}} & \tr{só a taxa}{rate only} & \tr{só a taxa}{rate only} & \tr{só a taxa}{rate only} \\
        $dX,\ dY$ & \textbf{\textcolor{mTeal}{$\checkmark$}} & --- & --- & --- \\
        \bottomrule
      \end{tabular}
      \par\medskip
      \scriptsize\textcolor{mStoneMid}{\tr{VLBI: antenas $\times$ quasares. GNSS: receptores $\times$ satélites GPS e outros. SLR: laser em satélites com espelhos. DORIS: balizas em terra $\times$ efeito Doppler em satélites.}{VLBI: antennas $\times$ quasars. GNSS: receivers $\times$ GPS and other satellites. SLR: laser on satellites with mirrors. DORIS: ground beacons $\times$ Doppler shift on satellites.}}
    \end{column}
    \begin{column}{0.34\textwidth}
      \begin{cuidado}[title={\tr{Por que só o VLBI?}{Why only VLBI?}}]
        \footnotesize\tr{Satélites giram em volta do \emph{centro de massa} da Terra: sentem a Terra por baixo, mas \alert{não sabem onde está o ``céu parado''}. Só os quasares sabem. Por isso $UT1$ e $dX, dY$ \bnum{absolutos} só vêm do VLBI.}{Satellites orbit the Earth's \emph{center of mass}: they feel the Earth below, but \alert{do not know where the ``still sky'' is}. Only quasars do. That is why \bnum{absolute} $UT1$ and $dX, dY$ come only from VLBI.}
      \end{cuidado}
    \end{column}
  \end{columns}
\end{frame}

\begin{frame}{\tr{Quem publica?}{Who publishes?}}
  \begin{columns}[T]
    \begin{column}{0.46\textwidth}
      \small
      \tr{O \bnum{IERS} (\emph{International Earth Rotation and Reference Systems Service}, criado em 1987) reúne os resultados dos serviços de cada técnica (IVS, IGS, ILRS, IDS) e publica os EOP.}{The \bnum{IERS} (\emph{International Earth Rotation and Reference Systems Service}, created in 1987) gathers the results of each technique's service (IVS, IGS, ILRS, IDS) and publishes the EOP.}
      \medskip

      \tr{Quem precisa do valor \emph{agora} (um receptor GNSS, uma antena de VLBI) usa valores \bnum{rápidos} ou \bnum{previstos}; quem precisa da melhor resposta espera os valores \bnum{finais}.}{Whoever needs the value \emph{now} (a GNSS receiver, a VLBI antenna) uses \bnum{rapid} or \bnum{predicted} values; whoever needs the best answer waits for the \bnum{final} values.}
    \end{column}
    \begin{column}{0.50\textwidth}
      \footnotesize
      \renewcommand{\arraystretch}{1.4}
      \begin{tabular}{@{}>{\raggedright\arraybackslash}p{1.5cm}>{\raggedright\arraybackslash}p{4.8cm}@{}}
        \toprule
        \textbf{EOP C04} & \tr{série diária final, com atraso de semanas}{final daily series, with a delay of weeks} \\
        \textbf{\tr{Boletim A}{Bulletin A}} & \tr{semanal: valores rápidos e previsões até 1 ano}{weekly: rapid values and predictions up to 1 year} \\
        \textbf{\tr{Boletim B}{Bulletin B}} & \tr{mensal: valores finais}{monthly: final values} \\
        \textbf{\tr{Boletim C}{Bulletin C}} & \tr{anuncia segundos intercalares (duas vezes por ano)}{announces leap seconds (twice a year)} \\
        \textbf{\tr{Boletim D}{Bulletin D}} & \tr{informa DUT1 $=$ UT1$-$UTC, arredondado a $0{,}1$ s}{reports DUT1 $=$ UT1$-$UTC, rounded to $0.1$ s} \\
        \bottomrule
      \end{tabular}
    \end{column}
  \end{columns}
\end{frame}

\begin{frame}{\tr{O tamanho de cada passo}{The size of each step}}
  \centering
  \small
  \renewcommand{\arraystretch}{1.25}
  \begin{tabular}{@{}c>{\raggedright\arraybackslash}p{4.9cm}rr@{}}
    \toprule
    \tr{degrau}{step} & \tr{o quê}{what} & \tr{ângulo}{angle} & \tr{no chão}{on the ground} \\
    \midrule
    $Q$ & \tr{precessão desde J2000 (em 2026,8)}{precession since J2000 (in 2026.8)} & $537\asec$ & \bnum{\tr{16,6 km}{16.6 km}} \\
    $Q$ & \tr{nutação principal}{main nutation} & $\num{9.2}\asec$ & \bnum{284 m} \\
    $W$ & \tr{movimento polar (máximo)}{polar motion (maximum)} & $\num{0.3}\asec$ & \bnum{\tr{9,3 m}{9.3 m}} \\
    $Q$ & \tr{viés do ICRS $\times$ J2000 (o polo)}{ICRS $\times$ J2000 bias (the pole)} & $\num{17.9}\mas$ & \bnum{\tr{0,55 m}{0.55 m}} \\
    $R$ & \tr{UT1 errado em 1 ms (no Equador)}{UT1 off by 1 ms (at the Equator)} & $15\mas$ & \bnum{\tr{0,47 m}{0.47 m}} \\
    $Q$ & \tr{correções $dX$, $dY$ típicas}{typical $dX$, $dY$ corrections} & $\num{0.2}\mas$ & \bnum{6 mm} \\
    \bottomrule
  \end{tabular}
  \medskip

  \tr{Cada degrau tem um tamanho. \alert{Esquecer um degrau dá um erro do tamanho dele}: por isso a ponte tem de ser completa --- e os EOP, atuais.}{Each step has a size. \alert{Forgetting a step gives an error of its size}: that is why the bridge must be complete --- and the EOP, up to date.}
\end{frame}

\begin{frame}{\tr{Experimente: o simulador GNSS}{Try it: the GNSS simulator}}
  \begin{columns}[T]
    \begin{column}{0.54\textwidth}
      \small
      \tr{Na pasta \texttt{GNSS\_simul} há um globo com os satélites GPS, GLONASS, Galileo e BeiDou.}{In the \texttt{GNSS\_simul} folder there is a globe with the GPS, GLONASS, Galileo and BeiDou satellites.}
      \begin{enumerate}
        \setlength{\itemsep}{2pt}
        \item \tr{Ligue ``Mostrar Trajetórias Orbitais'': as linhas são as órbitas \emph{vistas do chão}.}{Turn on ``Show Orbit Trajectories'': the lines are the orbits \emph{seen from the ground}.}
        \item \tr{O GPS dá 2 voltas por dia sideral: o traço quase se repete a cada dia. O Galileo (17 voltas em 10 dias siderais) só repete depois de 10 dias.}{GPS makes 2 turns per sidereal day: the track almost repeats every day. Galileo (17 turns in 10 sidereal days) only repeats after 10 days.}
        \item \tr{Clique no globo: o receptor vê os satélites acima de $10\gr$ e o simulador calcula o DOP.}{Click on the globe: the receiver sees the satellites above $10\gr$ and the simulator computes the DOP.}
      \end{enumerate}
    \end{column}
    \begin{column}{0.42\textwidth}
      \begin{cuidado}[title={\tr{O que ele deixa de fora}{What it leaves out}}]
        \footnotesize\tr{Passa do céu ao chão só com o ângulo sideral ($R$). Não aplica $W$ (polo, até 9 m) e usa o relógio do computador como UT1 ($\le 0{,}9$ s: até $\approx$ 1,7 km na órbita). Serve para a \emph{geometria} (DOP, visibilidade), não para posicionar com precisão.}{It goes from sky to ground with the sidereal angle only ($R$). It does not apply $W$ (pole, up to 9 m) and uses the computer clock as UT1 ($\le 0.9$ s: up to $\approx$ 1.7 km at the orbit). It is good for \emph{geometry} (DOP, visibility), not for precise positioning.}
      \end{cuidado}
    \end{column}
  \end{columns}
\end{frame}

\begin{frame}{\tr{Para levar para casa --- aula 3}{Take-home --- class 3}}
  \footnotesize
  \begin{enumerate}
    \setlength{\itemsep}{2pt}
    \item \tr{Céu parado e chão girando: converter de um para o outro é \bnum{girar eixos}. A ordem dos giros importa: $\vec r_{\text{ICRF}} = Q\,R\,W\,\vec r_{\text{ITRF}}$.}{Still sky, turning ground: converting between them is \bnum{turning axes}. The order of the turns matters: $\vec r_{\text{ICRF}} = Q\,R\,W\,\vec r_{\text{ITRF}}$.}
    \item \tr{\bnum{$R$, a rotação:} o ângulo ERA, $360{,}9856\gr$ por dia solar (o dia sideral dura \mbox{$23$ h $56$ min $4{,}09$ s}). 1 ms de erro em UT1 $=$ $0{,}47$ m no Equador e $1{,}9$ m na órbita do GPS.}{\bnum{$R$, the rotation:} the ERA angle, $360.9856\gr$ per solar day (the sidereal day lasts \mbox{$23$ h $56$ min $4.09$ s}). 1 ms of error in UT1 $=$ $0.47$ m at the Equator and $1.9$ m at the GPS orbit.}
    \item \tr{\bnum{$Q$, precessão e nutação:} o CIP anda no céu ($20\asec$/ano, tremor de $9{,}2\asec$); o modelo da IAU tem 1\,365 termos e o VLBI mede o resto ($dX$, $dY$ $\approx 0{,}2\mas$).}{\bnum{$Q$, precession and nutation:} the CIP moves in the sky ($20\asec$/yr, a $9.2\asec$ wobble); the IAU model has 1,365 terms and VLBI measures the rest ($dX$, $dY$ $\approx 0.2\mas$).}
    \item \tr{\bnum{$W$, movimento polar:} o polo anda até $0{,}3\asec$ ($\approx$ 9 m) na crosta: onda de Chandler (433 d) mais onda anual (365 d), com batimento de $\approx 6{,}4$ anos.}{\bnum{$W$, polar motion:} the pole moves up to $0.3\asec$ ($\approx$ 9 m) in the crust: Chandler wave (433 d) plus annual wave (365 d), beating every $\approx 6.4$ years.}
    \item \tr{Os \bnum{EOP} ($x_p$, $y_p$, UT1$-$UTC, $dX$, $dY$) são medidos por VLBI, GNSS, SLR e DORIS --- só o VLBI dá UT1 e $dX$, $dY$ absolutos --- e publicados pelo \bnum{IERS}.}{The \bnum{EOP} ($x_p$, $y_p$, UT1$-$UTC, $dX$, $dY$) are measured by VLBI, GNSS, SLR and DORIS --- only VLBI gives absolute UT1 and $dX$, $dY$ --- and published by the \bnum{IERS}.}
  \end{enumerate}
  \begin{ideia}
    \scriptsize \tr{\textbf{Na próxima aula:} o \bnum{tempo} --- UT1, UTC, TAI, o tempo do GNSS e o segundo intercalar --- e a prática: da órbita ao receptor, e o Brasil nessa história.}{\textbf{Next class:} \bnum{time} --- UT1, UTC, TAI, GNSS time and the leap second --- and practice: from the orbit to the receiver, and Brazil in this story.}
  \end{ideia}
\end{frame}

% ==================================================================
%  AULA 4 · O TEMPO E A PRÁTICA
% ==================================================================
\begin{frame}[standout]
  \tr{Parte 2}{Part 2}\\[6pt]
  {\normalsize \tr{O tempo e a prática}{Time and practice}}\\[4pt]
  {\scriptsize \tr{aula 4 de 4}{class 4 of 4}}
\end{frame}

\section{\tr{Vários relógios, uma hora}{Several clocks, one hour}}

\begin{frame}{\tr{Onde paramos: a aula 3}{Where we left off: class 3}}
  \begin{columns}[T]
    \begin{column}{0.56\textwidth}
      \small
      \tr{Na aula 3:}{In class 3:}
      \begin{enumerate}
        \setlength{\itemsep}{4pt}
        \item \tr{a ponte entre o céu e o chão tem três degraus: $Q$ (precessão e nutação), $R$ (rotação) e $W$ (movimento polar);}{the bridge between the sky and the ground has three steps: $Q$ (precession and nutation), $R$ (rotation) and $W$ (polar motion);}
        \item \tr{cinco números por dia, os \bnum{EOP}, dão os ângulos: $x_p$, $y_p$, UT1$-$UTC, $dX$, $dY$;}{five numbers per day, the \bnum{EOP}, give the angles: $x_p$, $y_p$, UT1$-$UTC, $dX$, $dY$;}
        \item \tr{1 ms de erro em UT1 $=$ $0{,}47$ m no Equador e $1{,}9$ m na órbita do GPS.}{1 ms of error in UT1 $=$ $0.47$ m at the Equator and $1.9$ m at the GPS orbit.}
      \end{enumerate}
    \end{column}
    \begin{column}{0.40\textwidth}
      \begin{cuidado}[title={\tr{A pergunta de hoje}{Today's question}}]
        \tr{O que é UT1, de onde vem esse ``tempo da rotação''? E que horas são \alert{para um satélite} que anda a $3{,}9$ km/s?}{What is UT1, where does this ``rotation time'' come from? And what time is it \alert{for a satellite} moving at $3.9$ km/s?}
      \end{cuidado}
    \end{column}
  \end{columns}
\end{frame}

\begin{frame}{\tr{As perguntas de hoje}{Today's questions}}
  \begin{columns}[T]
    \begin{column}{0.60\textwidth}
      \begin{enumerate}
        \setlength{\itemsep}{7pt}
        \item \tr{Por que existem tantos ``tempos'': UT1, UTC, TAI, o do GPS?}{Why are there so many ``times'': UT1, UTC, TAI, the GPS one?}
        \item \tr{O que é o segundo intercalar --- e por que o mundo decidiu acabar com ele, \bnum{esta semana}?}{What is the leap second --- and why did the world decide to end it, \bnum{this week}?}
        \item \tr{Como a órbita do satélite e a Terra que gira viram uma posição no chão?}{How do the satellite's orbit and the spinning Earth become a position on the ground?}
        \item \tr{O que isso tem a ver com o Brasil?}{What does this have to do with Brazil?}
      \end{enumerate}
    \end{column}
    \begin{column}{0.36\textwidth}
      \begin{ideia}[title={\tr{Combinado}{Deal}}]
        \tr{No fim da aula, você responde às quatro --- e fecha a Parte 2.}{By the end of the class you will answer all four --- and close Part 2.}
      \end{ideia}
    \end{column}
  \end{columns}
\end{frame}

\begin{frame}{\tr{O segundo: da Terra ao átomo}{The second: from the Earth to the atom}}
  \centering
  \begin{tikzpicture}[caixa/.append style={font=\small, minimum height=1.5cm}]
    \node[caixa, draw=corTempo, fill=corTempo!10, text width=3.7cm] (a) at (0,0) {\textbf{\tr{até 1960}{until 1960}}\\\tr{$1/\num{86400}$ do dia solar médio}{$1/\num{86400}$ of the mean solar day}\\{\footnotesize\tmpo{\tr{a Terra é o relógio}{the Earth is the clock}}}};
    \node[caixa, draw=corSist, fill=corSist!8, text width=3.7cm] (b) at (4.7,0) {\textbf{1960--1967}\\\tr{$1/\num{31556925.9747}$ do ano tropical de 1900}{$1/\num{31556925.9747}$ of the tropical year 1900}};
    \node[caixa, draw=corFrame, fill=corFrame!10, text width=4.0cm] (c) at (9.5,0) {\textbf{\tr{desde 1967}{since 1967}}\\\tr{$\num{9192631770}$ períodos do césio-133}{$\num{9192631770}$ periods of cesium-133}\\{\footnotesize\fram{\tr{o átomo é o relógio}{the atom is the clock}}}};
    \draw[seta] (a) -- (b);
    \draw[seta] (b) -- (c);
  \end{tikzpicture}
  \bigskip

  \begin{columns}[T]
    \begin{column}{0.48\textwidth}
      \small
      \tr{A Terra gira de modo irregular (aula 3); o átomo \bnum{não}. Por isso o segundo passou a ser definido pelo átomo.}{The Earth spins irregularly (class 3); the atom does \bnum{not}. That is why the second came to be defined by the atom.}
    \end{column}
    \begin{column}{0.48\textwidth}
      \begin{ideia}[title={\tr{Quão bom é um relógio atômico?}{How good is an atomic clock?}}]
        \footnotesize\tr{Césio: $\approx$ 1 s em \bnum{300 milhões de anos}. Relógios ópticos: $\approx$ 1 s em \bnum{30 bilhões de anos} --- mais que a idade do Universo (13,8 bilhões).}{Cesium: $\approx$ 1 s in \bnum{300 million years}. Optical clocks: $\approx$ 1 s in \bnum{30 billion years} --- longer than the age of the Universe (13.8 billion).}
      \end{ideia}
    \end{column}
  \end{columns}
\end{frame}

\begin{frame}{\tr{Dois relógios: a Terra e o átomo}{Two clocks: the Earth and the atom}}
  \centering
  \small
  \renewcommand{\arraystretch}{1.5}
  \begin{tabular}{@{}l>{\raggedright\arraybackslash}p{5.4cm}>{\raggedright\arraybackslash}p{4.8cm}@{}}
    \toprule
    & \tr{o que é}{what it is} & \tr{para que serve}{what it is for} \\
    \midrule
    \tmpo{UT1} & \tr{a hora da \emph{rotação}: vem do ERA; irregular}{the time of the \emph{rotation}: comes from the ERA; irregular} & \tr{orientar a Terra (o degrau $R$)}{orienting the Earth (the $R$ step)} \\
    \tmpo{TAI} & \tr{átomos: segundos do SI contados sem parar, desde 1958}{atoms: SI seconds counted without stopping, since 1958} & \tr{tempo científico; base de todos os outros}{scientific time; the basis of all the others} \\
    \tmpo{UTC} & \tr{TAI menos um número inteiro de segundos; é a hora civil}{TAI minus a whole number of seconds; it is civil time} & \tr{relógios do dia a dia; $|\text{UT1}-\text{UTC}| < \num{0.9}$ s (regra atual)}{everyday clocks; $|\text{UT1}-\text{UTC}| < \num{0.9}$ s (current rule)} \\
    \bottomrule
  \end{tabular}
  \bigskip

  \tr{Hoje: $\text{UTC} = \text{TAI} - \bnum{37}$ s \hspace{0.6em} (desde 1/1/2017).}{Today: $\text{UTC} = \text{TAI} - \bnum{37}$ s \hspace{0.6em} (since 1 Jan 2017).}
\end{frame}

\begin{frame}{\tr{A escada dos segundos intercalares}{The ladder of leap seconds}}
  \begin{columns}[T]
    \begin{column}{0.52\textwidth}
      \centering
      \begin{tikzpicture}[x=0.118cm, y=0.165cm]
        \draw[mStoneMid] (0,0) -- (56,0);
        \draw[mStoneMid] (0,0) -- (0,28.5);
        \foreach \a in {1980,1990,2000,2010,2020} {
          \draw[mStoneMid] ({\a-1972},-0.4) -- ({\a-1972},0.4);
          \node[rotulo, below] at ({\a-1972},-0.6) {\a};
        }
        \foreach \yy/\lab in {0/10, 10/20, 20/30, 27/37} {
          \draw[mStoneMid] (-0.6,\yy) -- (0.6,\yy);
          \node[rotulo, left] at (-0.8,\yy) {\lab};
        }
        \draw[corTempo, thick] (0.0,0) -- (0.5,0) -- (0.5,1) -- (1.0,1) -- (1.0,2) -- (2.0,2) -- (2.0,3) -- (3.0,3) -- (3.0,4) -- (4.0,4) -- (4.0,5) -- (5.0,5) -- (5.0,6) -- (6.0,6) -- (6.0,7) -- (7.0,7) -- (7.0,8) -- (8.0,8) -- (8.0,9) -- (9.5,9) -- (9.5,10) -- (10.5,10) -- (10.5,11) -- (11.5,11) -- (11.5,12) -- (13.5,12) -- (13.5,13) -- (16.0,13) -- (16.0,14) -- (18.0,14) -- (18.0,15) -- (19.0,15) -- (19.0,16) -- (20.5,16) -- (20.5,17) -- (21.5,17) -- (21.5,18) -- (22.5,18) -- (22.5,19) -- (24.0,19) -- (24.0,20) -- (25.5,20) -- (25.5,21) -- (27.0,21) -- (27.0,22) -- (34.0,22) -- (34.0,23) -- (37.0,23) -- (37.0,24) -- (40.5,24) -- (40.5,25) -- (43.5,25) -- (43.5,26) -- (45.0,26) -- (45.0,27) -- (54.8,27);
        \node[rotulo, text=corErro, anchor=south east] at (45,27.6) {\tr{último: 31/12/2016}{last: 31 Dec 2016}};
        \node[rotulo, text=mStoneMid, anchor=north west] at (1,27.5) {TAI$-$UTC (s)};
      \end{tikzpicture}
    \end{column}
    \begin{column}{0.44\textwidth}
      \small
      \tr{Em 1972 o UTC passou a ser TAI menos um número inteiro de segundos. Desde então foram inseridos \bnum{27 segundos intercalares}: 22 até 1998 e 5 depois.}{In 1972 UTC became TAI minus a whole number of seconds. Since then \bnum{27 leap seconds} have been inserted: 22 up to 1998 and 5 afterwards.}
      \medskip

      \tr{Um dia com intercalar tem $\num{86401}$ s: o relógio marca \texttt{23:59:60}.}{A day with a leap second has $\num{86401}$ s: the clock shows \texttt{23:59:60}.}
      \medskip

      \tr{Em 2017 chegamos a $\text{TAI}-\text{UTC} = 37$ s. Dez anos depois, \alert{nenhum} novo.}{In 2017 we reached $\text{TAI}-\text{UTC} = 37$ s. Ten years later, \alert{none} more.}
    \end{column}
  \end{columns}
\end{frame}

\begin{frame}{\tr{Por que aparecem segundos extras?}{Why do extra seconds appear?}}
  \begin{columns}[T]
    \begin{column}{0.52\textwidth}
      \small
      \tr{Em média a Terra gira \bnum{mais devagar} que o relógio atômico: o dia solar dura um pouco mais que $\num{86400}$ s (marés, aula 3). A diferença \emph{se acumula}.}{On average the Earth spins \bnum{more slowly} than the atomic clock: the solar day lasts slightly more than $\num{86400}$ s (tides, class 3). The difference \emph{accumulates}.}
      \medskip

      \tr{Quando $\text{UT1}-\text{UTC}$ chega perto de $0{,}9$ s, o IERS (Boletim C) manda inserir um segundo no UTC, no fim de junho ou de dezembro. O UTC volta para perto do UT1.}{When $\text{UT1}-\text{UTC}$ gets close to $0.9$ s, the IERS (Bulletin C) orders a second to be inserted into UTC, at the end of June or December. UTC comes back close to UT1.}
    \end{column}
    \begin{column}{0.44\textwidth}
      \begin{regra}[title={\tr{Sua vez}{Your turn}}]
        \footnotesize
        \begin{enumerate}
          \setlength{\itemsep}{2pt}
          \item[(a)] \tr{Um dia dura $\num{86400.0015}$ s (1,5 ms a mais). Em quantos dias o atraso soma $0{,}9$ s?}{A day lasts $\num{86400.0015}$ s (1.5 ms more). In how many days does the lag add up to $0.9$ s?}
          \item[(b)] \tr{De 1972 a 2016: 27 intercalares em $44{,}5$ anos. Um a cada quanto tempo?}{From 1972 to 2016: 27 leap seconds in $44.5$ years. One every how long?}
        \end{enumerate}
      \end{regra}
      \onslide<2->{
      \begin{ideia}[title={\tr{Respostas}{Answers}}]
        \footnotesize\raggedright
        (a) $\num{0.9}/\num{0.0015} = \bnum{600}$ \tr{dias}{days} $\approx \num{1.6}$ \tr{ano}{yr}.\quad
        (b) $\num{44.5}/27 \approx \bnum{\num{1.6}}$ \tr{ano}{yr}: \tr{confere!}{matches!}
      \end{ideia}
      }
    \end{column}
  \end{columns}
\end{frame}

\begin{frame}{\tr{Desde 2016 a conta mudou}{Since 2016 the count has changed}}
  \begin{columns}[T]
    \begin{column}{0.56\textwidth}
      \small
      \tr{Nos últimos anos a Terra deu uma \bnum{acelerada}: os dias ficaram um pouco mais curtos (núcleo e manto, aula 3). E o último segundo intercalar foi o de 31/12/2016.}{In recent years the Earth has \bnum{sped up}: days became slightly shorter (core and mantle, class 3). And the last leap second was the one of 31 Dec 2016.}
      \medskip

      \tr{Com a Terra mais rápida, o UTC passaria a ``ficar atrás'': seria preciso tirar um segundo --- um \bnum{intercalar negativo}, que \alert{nunca} foi usado.}{With a faster Earth, UTC would start to ``fall behind'': a second would have to be removed --- a \bnum{negative leap second}, which has \alert{never} been used.}
      \medskip

      \tr{Agnew (\emph{Nature}, 2024): o derretimento do gelo, que desacelera a rotação, adiou esse problema em cerca de 3 anos --- para $\approx$ 2029.}{Agnew (\emph{Nature}, 2024): melting ice, which slows the rotation, postponed this problem by about 3 years --- to $\approx$ 2029.}
    \end{column}
    \begin{column}{0.40\textwidth}
      \begin{cuidado}[title={\tr{O clima mexe no relógio}{Climate moves the clock}}]
        \small\tr{A água do gelo que derrete vai dos polos para os oceanos, longe do eixo: o momento de inércia \emph{aumenta} e a Terra gira mais devagar --- como a patinadora abrindo os braços.}{Water from melting ice moves from the poles to the oceans, away from the axis: the moment of inertia \emph{increases} and the Earth spins more slowly --- like the skater opening her arms.}
      \end{cuidado}
    \end{column}
  \end{columns}
\end{frame}

\begin{frame}{\tr{Outubro de 2026: o fim do segundo intercalar}{October 2026: the end of the leap second}}
  \begin{columns}[T]
    \begin{column}{0.54\textwidth}
      \small
      \tr{Em \bnum{15/10/2026} a 28ª Conferência Geral de Pesos e Medidas (CGPM, Versalhes) adotou a \bnum{Resolução 5}:}{On \bnum{15 Oct 2026} the 28th General Conference on Weights and Measures (CGPM, Versailles) adopted \bnum{Resolution 5}:}
      \begin{itemize}
        \setlength{\itemsep}{3pt}
        \item \tr{o UTC passa a ser \bnum{contínuo}, sem segundos intercalares, a partir de \bnum{20/5/2027};}{UTC becomes \bnum{continuous}, with no leap seconds, from \bnum{20 May 2027};}
        \item \tr{o limite de $|\text{UT1}-\text{UTC}|$ sobe para \bnum{1 hora} ($\num{3600}$ s), a valer por séculos.}{the limit of $|\text{UT1}-\text{UTC}|$ rises to \bnum{1 hour} ($\num{3600}$ s), to hold for centuries.}
      \end{itemize}
      \tr{Em 2022 a 27ª CGPM (Res.~4) já tinha decidido ampliar esse limite até 2035.}{In 2022 the 27th CGPM (Res.~4) had already decided to enlarge that limit by 2035.}
    \end{column}
    \begin{column}{0.42\textwidth}
      \begin{armadilha}[title={\tr{Nota de rigor}{Note on rigor}}]
        \footnotesize\tr{Os detalhes acima vêm da declaração do presidente do CCTF e da imprensa técnica, de poucos dias atrás. O \alert{texto oficial} da Resolução 5 deve ser conferido no BIPM antes de citar.}{The details above come from the statement of the CCTF president and the technical press, from a few days ago. The \alert{official text} of Resolution 5 should be checked at the BIPM before quoting.}
      \end{armadilha}
    \end{column}
  \end{columns}
\end{frame}

\begin{frame}{\tr{O que isso muda para a geodésia}{What this changes for geodesy}}
  \begin{columns}[T]
    \begin{column}{0.54\textwidth}
      \small
      \tr{Hoje $|\text{UT1}-\text{UTC}| < 0{,}9$ s: ignorar a diferença custa até \bnum{419 m} no chão (Equador) e $\approx$ \bnum{1,7 km} na órbita do GPS.}{Today $|\text{UT1}-\text{UTC}| < 0.9$ s: ignoring the difference costs up to \bnum{419 m} on the ground (Equator) and $\approx$ \bnum{1.7 km} at the GPS orbit.}
      \medskip

      \tr{Sem os saltos, $\text{UT1}-\text{UTC}$ cresce livremente: de segundos para \bnum{minutos} e, só no futuro distante, até \bnum{1 hora}. Cada minuto são $15'$ de giro:}{Without the jumps, $\text{UT1}-\text{UTC}$ grows freely: from seconds to \bnum{minutes} and, only in the distant future, up to \bnum{1 hour}. Each minute is $15'$ of rotation:}
      \[ 60\ \text{s} \times 465\ \text{m/s} \approx \bnum{28}\ \text{km}\quad \text{\tr{(no Equador)}{(at the Equator)}} \]
      \tr{No limite de 1 hora seriam $\approx$ \bnum{1\,674 km}. O UTC continua sendo $\text{TAI} - 37$ s, sem saltos --- e $\text{UT1}-\text{UTC}$ vira um EOP que \bnum{todo programa} precisa ler do IERS.}{At the 1-hour limit it would be $\approx$ \bnum{1,674 km}. UTC stays $\text{TAI} - 37$ s, with no jumps --- and $\text{UT1}-\text{UTC}$ becomes an EOP that \bnum{every program} must read from the IERS.}
    \end{column}
    \begin{column}{0.42\textwidth}
      \begin{cuidado}[title={\tr{Quem se importa?}{Who cares?}}]
        \footnotesize\tr{Telescópios e antenas que apontam usando $\text{DUT1}$ ($\le 0{,}9$ s), planetários, softwares de satélites e simuladores como os desta pasta, que usam o relógio do computador como UT1.}{Telescopes and antennas that point using $\text{DUT1}$ ($\le 0.9$ s), planetariums, satellite software and simulators like the ones in this folder, which use the computer clock as UT1.}
      \end{cuidado}
    \end{column}
  \end{columns}
\end{frame}

% ------------------------------------------------------------------ O tempo dos satélites
\section{\tr{O tempo dos satélites}{Satellite time}}

\begin{frame}{\tr{Cada sistema tem a sua escala}{Each system has its own scale}}
  \centering
  \small
  \renewcommand{\arraystretch}{1.45}
  \begin{tabular}{@{}llll@{}}
    \toprule
    \tr{sistema}{system} & \tr{escala}{scale} & \tr{em relação ao UTC (2026)}{relative to UTC (2026)} & \tr{intercalares?}{leap seconds?} \\
    \midrule
    GPS & GPST & $\text{UTC} + 18$ s $\;(= \text{TAI} - 19$ s) & \tr{não}{no} \\
    Galileo & GST & \tr{$\approx$ GPST (dezenas de ns)}{$\approx$ GPST (tens of ns)} & \tr{não}{no} \\
    BeiDou & BDT & $\text{UTC} + 4$ s $\;(= \text{TAI} - 33$ s) & \tr{não}{no} \\
    GLONASS & \tr{UTC(SU)$+3$ h}{UTC(SU)$+3$ h} & \tr{segue o UTC (mais 3 h)}{follows UTC (plus 3 h)} & \tr{\textbf{sim}}{\textbf{yes}} \\
    \bottomrule
  \end{tabular}
  \bigskip

  \tr{Os receptores usam um parâmetro de segundos intercalares, enviado na mensagem de navegação, para passar de GPST para UTC.}{Receivers use a leap-second parameter, sent in the navigation message, to go from GPST to UTC.}
\end{frame}

\begin{frame}{\tr{A escada dos tempos}{The ladder of times}}
  \begin{columns}[T]
    \begin{column}{0.54\textwidth}
      \centering
      \begin{tikzpicture}[x=0.062cm, y=0.8cm]
        \foreach \yy/\nome/\val/\cor in {4/{TT}/69.184/corSist, 3/{TAI}/37/corTempo, 2/{GPST}/18/corFrame, 1/{BDT}/4/corCeu, 0/{UTC}/0/mStoneMid} {
          \fill[\cor!75] (0,\yy) rectangle (\val,{\yy+0.6});
          \node[rotulo, anchor=east] at (-1,{\yy+0.3}) {\nome};
          \node[rotulo, anchor=west] at ({\val+1},{\yy+0.3}) {$+\num{\val}$ s};
        }
        \draw[mStoneMid] (0,-0.3) -- (0,5);
      \end{tikzpicture}
      \par\smallskip
      {\scriptsize\textcolor{mStoneMid}{\tr{Adiantamento de cada escala em relação ao UTC, em 2026.}{Offset of each scale relative to UTC, in 2026.}}}
    \end{column}
    \begin{column}{0.42\textwidth}
      \small
      \tr{O \bnum{TT} (\emph{Terrestrial Time}) é o tempo das efemérides: $\text{TT} = \text{TAI} + \num{32.184}$ s. Entra nas fórmulas de precessão e nutação (degrau $Q$).}{\bnum{TT} (\emph{Terrestrial Time}) is the time of the ephemerides: $\text{TT} = \text{TAI} + \num{32.184}$ s. It enters the precession and nutation formulas (the $Q$ step).}
      \medskip

      \tr{Regra prática: \bnum{UT1} para o ângulo de rotação; \bnum{TT} para precessão e nutação.}{Rule of thumb: \bnum{UT1} for the rotation angle; \bnum{TT} for precession and nutation.}
    \end{column}
  \end{columns}
\end{frame}

\begin{frame}{\tr{Sua vez}{Your turn}}
  \begin{columns}[T]
    \begin{column}{0.54\textwidth}
      \begin{regra}[title={\tr{Contas de tempo}{Time sums}}]
        \small
        \tr{São \texttt{12:00:00} em UTC. Que horas são em:}{It is \texttt{12:00:00} UTC. What time is it in:}
        \begin{enumerate}
          \setlength{\itemsep}{3pt}
          \item[(a)] GPST?
          \item[(b)] TAI?
          \item[(c)] TT?
          \item[(d)] \tr{GLONASS (UTC(SU)$+3$ h)?}{GLONASS (UTC(SU)$+3$ h)?}
        \end{enumerate}
      \end{regra}
    \end{column}
    \begin{column}{0.42\textwidth}
      \onslide<2->{
      \begin{ideia}[title={\tr{Respostas}{Answers}}]
        \small\raggedright
        (a) \texttt{12:00:\bnum{18}}.

        (b) \texttt{12:00:\bnum{37}}.

        (c) \texttt{12:01:\bnum{\tr{09,184}{09.184}}} \tr{(37 + 32,184)}{(37 + 32.184)}.

        (d) \texttt{\bnum{15}:00:00}.
      \end{ideia}
      }
    \end{column}
  \end{columns}
\end{frame}

\begin{frame}{\tr{Relógios no espaço andam diferente}{Clocks in space run differently}}
  \begin{columns}[T]
    \begin{column}{0.54\textwidth}
      \small
      \tr{Dois efeitos de Einstein, por dia, no relógio do GPS:}{Two Einstein effects, per day, on the GPS clock:}
      \begin{itemize}
        \setlength{\itemsep}{3pt}
        \item \tr{lá em cima a gravidade é mais fraca: o relógio anda \bnum{mais depressa}, $+\num{45.7}\ \mu$s;}{up there gravity is weaker: the clock runs \bnum{faster}, $+\num{45.7}\ \mu$s;}
        \item \tr{o satélite se move a $3{,}9$ km/s: o relógio anda \bnum{mais devagar}, $-\num{7.2}\ \mu$s.}{the satellite moves at $3.9$ km/s: the clock runs \bnum{slower}, $-\num{7.2}\ \mu$s.}
      \end{itemize}
      \tr{Saldo: $+\num{38.5}\ \mu$s por dia. Sem correção, o erro seria $c \times 38{,}5\ \mu\text{s} \approx$ \bnum{11,6 km por dia}. Por isso a frequência de \num{10.23} MHz é ajustada para \num{10.22999999543} MHz ($\Delta f/f = \num{4.4645e-10}$) antes do lançamento.}{Balance: $+\num{38.5}\ \mu$s per day. Uncorrected, the error would be $c \times 38.5\ \mu\text{s} \approx$ \bnum{11.6 km per day}. That is why the \num{10.23} MHz frequency is adjusted to \num{10.22999999543} MHz ($\Delta f/f = \num{4.4645e-10}$) before launch.}
    \end{column}
    \begin{column}{0.42\textwidth}
      \centering
      \begin{tikzpicture}[x=1.0cm, y=0.062cm]
        \draw[mStoneMid] (-0.2,0) -- (5.2,0);
        \fill[corFrame!80] (0.2,0) rectangle (1.0,45.7);
        \fill[corErro!80] (2.0,45.7) rectangle (2.8,38.5);
        \fill[corSist!80] (3.8,0) rectangle (4.6,38.5);
        \node[rotulo, above, align=center] at (0.6,45.7) {$+\num{45.7}$\\\tr{gravidade}{gravity}};
        \node[rotulo, above, align=center, text=corErro] at (2.4,45.7) {$-\num{7.2}$\\\tr{velocidade}{velocity}};
        \node[rotulo, above, align=center, text=corSist] at (4.2,38.5) {$+\num{38.5}$\\\tr{saldo}{net}};
        \draw[mStoneMid, dotted] (1.0,45.7) -- (2.0,45.7);
        \draw[mStoneMid, dotted] (2.8,38.5) -- (3.8,38.5);
      \end{tikzpicture}
      \par\smallskip
      {\scriptsize\textcolor{mStoneMid}{\tr{$\mu$s por dia, relógio do GPS em relação ao relógio no chão.}{$\mu$s per day, GPS clock relative to a clock on the ground.}}}
    \end{column}
  \end{columns}
\end{frame}

\begin{frame}{\tr{Relógios que medem altura}{Clocks that measure height}}
  \begin{columns}[T]
    \begin{column}{0.54\textwidth}
      \small
      \tr{A mesma física vale na escada de um prédio: um relógio \emph{mais alto} anda mais depressa. A diferença é pequena:}{The same physics holds on the stairs of a building: a \emph{higher} clock runs faster. The difference is tiny:}
      \[ \frac{\Delta f}{f} = \frac{g\,h}{c^{2}} \approx \num{1.09e-16}\ \text{\tr{por metro}{per meter}} \]
      \tr{Relógios ópticos de $10^{-18}$ \bnum{sentem 1 cm} de diferença de altura. Em 2020, dois deles, no pé e no alto da torre Tokyo Skytree (\bnum{450 m}), mediram o efeito esperado ($\approx \num{4.9e-14}$).}{Optical clocks of $10^{-18}$ \bnum{feel 1 cm} of height difference. In 2020, two of them, at the foot and at the top of the Tokyo Skytree tower (\bnum{450 m}), measured the expected effect ($\approx \num{4.9e-14}$).}
    \end{column}
    \begin{column}{0.42\textwidth}
      \begin{ideia}[title={\tr{Geodésia cronométrica}{Chronometric geodesy}}]
        \small\tr{Comparar relógios dá o \bnum{potencial gravitacional} $W$ --- o mesmo $W$ que define o \emph{geoide} nas aulas de gravimetria. Um dia, relógios podem substituir parte do nivelamento.}{Comparing clocks gives the \bnum{gravitational potential} $W$ --- the same $W$ that defines the \emph{geoid} in the gravimetry classes. One day, clocks may replace part of leveling.}
      \end{ideia}
    \end{column}
  \end{columns}
\end{frame}

% ------------------------------------------------------------------ Da órbita ao receptor
\section{\tr{Da órbita ao receptor}{From the orbit to the receiver}}

\begin{frame}{\tr{A órbita do GPS em duas contas}{The GPS orbit in two sums}}
  \begin{columns}[T]
    \begin{column}{0.46\textwidth}
      \small
      \tr{Em órbita circular a gravidade faz o papel da força centrípeta: $v^{2} = GM/r$, e o período é}{In a circular orbit gravity plays the role of the centripetal force: $v^{2} = GM/r$, and the period is}
      \[ T = 2\pi\sqrt{\frac{r^{3}}{GM}} \]
      \tr{O GPS dá \bnum{2 voltas por dia sideral}: \mbox{$T = 11$ h $58$ min $2$ s}. Resultado: $r \approx \num{26562}$ km (altitude $\approx \num{20184}$ km), $v \approx 3{,}87$ km/s.}{GPS makes \bnum{2 turns per sidereal day}: \mbox{$T = 11$ h $58$ min $2$ s}. Result: $r \approx \num{26562}$ km (altitude $\approx \num{20184}$ km), $v \approx 3.87$ km/s.}
      \medskip

      \tr{Por isso o traço do satélite no chão \bnum{se repete} a cada dia sideral.}{That is why the satellite's track on the ground \bnum{repeats} every sidereal day.}
    \end{column}
    \begin{column}{0.50\textwidth}
      \scriptsize
      \renewcommand{\arraystretch}{1.4}
      \setlength{\tabcolsep}{3pt}
      \begin{tabular}{@{}lccr@{}}
        \toprule
        \tr{sistema}{system} & \tr{voltas/dias}{turns/days} & \tr{período}{period} & \tr{altitude}{altitude} \\
        \midrule
        GPS & 2/1 & 11 h 58 min & $\num{20184}$ km \\
        GLONASS & 17/8 & 11 h 16 min & $\num{19132}$ km \\
        Galileo & 17/10 & 14 h 05 min & $\num{23223}$ km \\
        BeiDou (MEO) & 13/7 & 12 h 53 min & $\num{21529}$ km \\
        \bottomrule
      \end{tabular}
      \par\smallskip
      \textcolor{mStoneMid}{\tr{``Voltas/dias'': voltas do satélite por dias siderais até o traço se repetir. Valores nominais, calculados a partir da relação de repetição.}{``Turns/days'': satellite turns per sidereal days until the track repeats. Nominal values, computed from the repeat ratio.}}
    \end{column}
  \end{columns}
\end{frame}

\begin{frame}{\tr{Do TLE ao mapa: a cadeia do simulador}{From the TLE to the map: the simulator's chain}}
  \centering
  \begin{tikzpicture}[caixa/.append style={font=\small, text width=3.7cm, minimum height=1.2cm}]
    \node[caixa, draw=corSist, fill=corSist!8] (a) at (0,2.1) {TLE\\{\footnotesize\tr{elementos da órbita}{orbital elements}}};
    \node[caixa, draw=corSist, fill=corSist!8] (b) at (5.0,2.1) {SGP4\\{\footnotesize\tr{propaga a órbita no tempo}{propagates the orbit in time}}};
    \node[caixa, draw=corSist, fill=corSist!8] (c) at (10.0,2.1) {TEME $(x,y,z)$\\{\footnotesize\tr{no espaço, ``parado''}{in space, ``still''}}};
    \node[caixa, draw=corTempo, fill=corTempo!10] (d) at (10.0,0) {$R_3(\text{GMST})$\\{\footnotesize\tr{gira a Terra: o degrau $R$}{turns the Earth: the $R$ step}}};
    \node[caixa, draw=corFrame, fill=corFrame!10] (e) at (5.0,0) {ECEF $(X,Y,Z)$\\{\footnotesize\tr{no chão, quase ITRF}{on the ground, almost ITRF}}};
    \node[caixa, draw=corFrame, fill=corFrame!10] (f) at (0,0) {$(\varphi,\lambda,h)$ \tr{e}{and} $(Az, El)$\\{\footnotesize\tr{mapa, visibilidade, DOP}{map, visibility, DOP}}};
    \draw[seta] (a) -- (b);
    \draw[seta] (b) -- (c);
    \draw[seta] (c) -- (d);
    \draw[seta] (d) -- (e);
    \draw[seta] (e) -- (f);
  \end{tikzpicture}
  \bigskip

  \small
  \tr{\bnum{TEME} (\emph{True Equator, Mean Equinox}): o referencial em que o SGP4 entrega a posição. No \texttt{GNSS\_simul} (\texttt{satellite.js}) a passagem para o chão usa \bnum{só o ângulo sideral}: sem $W$ e com $\text{UT1} = \text{UTC}$.}{\bnum{TEME} (\emph{True Equator, Mean Equinox}): the frame in which SGP4 delivers the position. In the \texttt{GNSS\_simul} (\texttt{satellite.js}) the step to the ground uses \bnum{only the sidereal angle}: no $W$ and with $\text{UT1} = \text{UTC}$.}
\end{frame}

\begin{frame}{\tr{Para precisão: o que muda}{For precision: what changes}}
  \centering
  \small
  \renewcommand{\arraystretch}{1.3}
  \begin{tabular}{@{}>{\raggedright\arraybackslash}p{2.6cm}>{\raggedright\arraybackslash}p{4.6cm}>{\raggedright\arraybackslash}p{6.0cm}@{}}
    \toprule
    & \tr{simulador (geometria)}{simulator (geometry)} & \tr{posicionamento preciso}{precise positioning} \\
    \midrule
    \textbf{\tr{órbita}{orbit}} & \tr{TLE: da ordem de 1 km}{TLE: of the order of 1 km} & \tr{órbitas precisas do IGS (cm), no ITRF}{IGS precise orbits (cm), in the ITRF} \\
    \textbf{\tr{céu $\to$ chão}{sky $\to$ ground}} & $R_3(\text{GMST})$ & \tr{$Q\,R\,W$ completa, com EOP do IERS}{full $Q\,R\,W$, with EOP from the IERS} \\
    \textbf{\tr{tempo}{time}} & \tr{UTC no lugar de UT1}{UTC in place of UT1} & \tr{$\text{UT1}-\text{UTC}$ do IERS; TT onde couber}{$\text{UT1}-\text{UTC}$ from the IERS; TT where due} \\
    \textbf{\tr{relógios}{clocks}} & \tr{ignorados}{ignored} & \tr{produtos de relógio do IGS + relatividade}{IGS clock products + relativity} \\
    \textbf{\tr{quadro do usuário}{user's frame}} & \tr{latitude e longitude do clique}{latitude and longitude of the click} & \tr{SIRGAS2000 ($\approx$ ITRF2000, época 2000,4)}{SIRGAS2000 ($\approx$ ITRF2000, epoch 2000.4)} \\
    \bottomrule
  \end{tabular}
  \medskip

  \tr{O simulador \alert{não é errado}: é \emph{barato}. Serve para ver geometria; para centímetros, a ponte tem de estar completa.}{The simulator is \alert{not wrong}: it is \emph{cheap}. It is good for seeing geometry; for centimeters, the bridge must be complete.}
\end{frame}

\begin{frame}{\tr{Sua vez}{Your turn}}
  \begin{columns}[T]
    \begin{column}{0.54\textwidth}
      \begin{regra}[title={\tr{O erro do simulador}{The simulator's error}}]
        \small
        \tr{Um simulador usa UTC no lugar de UT1, mas hoje $\text{UT1}-\text{UTC} = -0{,}4$ s. Use $\omega = \num{7.292115e-5}$ rad/s.}{A simulator uses UTC in place of UT1, but today $\text{UT1}-\text{UTC} = -0.4$ s. Use $\omega = \num{7.292115e-5}$ rad/s.}
        \begin{enumerate}
          \setlength{\itemsep}{3pt}
          \item[(a)] \tr{Quanto erra a posição de um satélite GPS ($r = \num{26560}$ km)?}{How far off is the position of a GPS satellite ($r = \num{26560}$ km)?}
          \item[(b)] \tr{E a de uma estação em Curitiba (distância ao eixo $\num{5763}$ km)?}{And that of a station in Curitiba (distance to the axis $\num{5763}$ km)?}
          \item[(c)] \tr{Em que direção é o erro?}{In which direction is the error?}
        \end{enumerate}
      \end{regra}
    \end{column}
    \begin{column}{0.42\textwidth}
      \onslide<2->{
      \begin{ideia}[title={\tr{Respostas}{Answers}}]
        \small\raggedright
        (a) $\omega\,r\,\Delta t = \num{7.292e-5} \times \num{2.656e7} \times \num{0.4} \approx \bnum{775}$ m.
        \smallskip

        (b) $\num{7.292e-5} \times \num{5.763e6} \times \num{0.4} \approx \bnum{168}$ m.
        \smallskip

        (c) \tr{Ao longo do paralelo (\emph{leste--oeste}): é um giro errado em torno do eixo.}{Along the parallel (\emph{east--west}): it is a wrong turn about the axis.}
      \end{ideia}
      }
    \end{column}
  \end{columns}
\end{frame}

% ------------------------------------------------------------------ O Brasil e o fecho
\section{\tr{O Brasil e o fecho}{Brazil and the wrap-up}}

\begin{frame}{\tr{O Brasil nessa história}{Brazil in this story}}
  \small
  \begin{columns}[T]
    \begin{column}{0.31\textwidth}
      \begin{cuidado}[title={\tr{A hora}{The time}}]
        \footnotesize\raggedright\tr{O \bnum{Observatório Nacional} (Rio de Janeiro) gera a \bnum{Hora Legal Brasileira} com um conjunto de relógios atômicos, ligada ao UTC do BIPM (Decreto 4.264/2002).}{The \bnum{National Observatory} (Rio de Janeiro) generates the \bnum{Brazilian Legal Time} with a set of atomic clocks, linked to the BIPM's UTC (Decree 4,264/2002).}
      \end{cuidado}
    \end{column}
    \begin{column}{0.31\textwidth}
      \begin{regra}[title={\tr{O céu}{The sky}}]
        \footnotesize\raggedright\tr{\bnum{Fortaleza} (CE): estação de VLBI do IVS desde 1993 (Rádio Observatório Espacial do Nordeste). Em nov./2025, nova antena com terminal VGOS. O Hemisfério Sul tem poucas: ajuda os EOP.}{\bnum{Fortaleza} (CE): IVS VLBI station since 1993 (Northeast Space Radio Observatory). In Nov 2025, a new antenna with a VGOS terminal. The Southern Hemisphere has few: it helps the EOP.}
      \end{regra}
    \end{column}
    \begin{column}{0.31\textwidth}
      \begin{ideia}[title={\tr{O chão}{The ground}}]
        \footnotesize\raggedright\tr{A \bnum{RBMC} (IBGE) tem mais de 150 estações GNSS contínuas: a base do \bnum{SIRGAS2000} ($\approx$ ITRF2000, época 2000,4) e parte da rede SIRGAS-CON.}{The \bnum{RBMC} (IBGE) has more than 150 continuous GNSS stations: the basis of \bnum{SIRGAS2000} ($\approx$ ITRF2000, epoch 2000.4) and part of the SIRGAS-CON network.}
      \end{ideia}
    \end{column}
  \end{columns}
  \bigskip

  \centering
  \tr{Hora, céu e chão: \bnum{os três lados do mesmo triângulo} que esta Parte 2 e a Parte 1 desenharam.}{Time, sky and ground: \bnum{the three sides of the same triangle} that this Part 2 and Part 1 drew.}
\end{frame}

\begin{frame}{\tr{Experimente: o relógio dos simuladores}{Try it: the simulators' clock}}
  \begin{columns}[T]
    \begin{column}{0.54\textwidth}
      \small
      \tr{No \texttt{astrometria} a hora de Greenwich na tela é UTC; o céu, porém, gira com UT1. Com $|\text{DUT1}| \le 0{,}9$ s o céu sai até $13{,}5\asec$ do lugar ($\approx$ 419 m). Faltam também a precessão ($\le 0{,}37\gr$ desde J2000) e a nutação ($9{,}2\asec$).}{In \texttt{astrometria} the Greenwich time on screen is UTC; the sky, however, turns with UT1. With $|\text{DUT1}| \le 0.9$ s the sky ends up to $13.5\asec$ off ($\approx$ 419 m). Precession ($\le 0.37\gr$ since J2000) and nutation ($9.2\asec$) are also missing.}
      \medskip

      \bnum{\tr{Desafios}{Challenges}}
      \begin{enumerate}
        \setlength{\itemsep}{2pt}
        \item \tr{Leve Sirius de J2000 a 2026,8 com as fórmulas do apêndice A4: as coordenadas dele mudam $17{,}3\amin$ ($\approx$ 32 km num ponto astronômico).}{Take Sirius from J2000 to 2026.8 with the formulas of appendix A4: its coordinates change by $17.3\amin$ ($\approx$ 32 km in an astronomical fix).}
        \item \tr{Troque UTC por UT1 usando um $\text{UT1}-\text{UTC}$ do IERS e veja a diferença.}{Replace UTC with UT1 using a $\text{UT1}-\text{UTC}$ from the IERS and see the difference.}
      \end{enumerate}
    \end{column}
    \begin{column}{0.42\textwidth}
      \begin{ideia}[title={\tr{Ordem de grandeza}{Order of magnitude}}]
        \small\tr{Um minuto de arco no céu vale uma milha náutica ($1{,}852$ km) no chão. Por isso $17{,}3\amin$ de erro na estrela são $\approx$ 32 km de erro na posição.}{One arcminute in the sky is one nautical mile ($1.852$ km) on the ground. That is why $17.3\amin$ of error on the star is $\approx$ 32 km of error in the position.}
      \end{ideia}
    \end{column}
  \end{columns}
\end{frame}

\begin{frame}{\tr{Orçamento de erros: tudo de uma vez}{Error budget: everything at once}}
  \centering
  \begin{tikzpicture}[x=0.86cm, y=0.47cm]
    \foreach \yy/\val/\cor/\lab in {
      9/7.22/corCeu/{\tr{precessão não aplicada (2026,8)}{precession not applied (2026.8)}},
      8/7.06/corTempo/{\tr{relatividade do GPS, 1 dia}{GPS relativity, 1 day}},
      7/6.24/corTempo/{\tr{UT1$-$UTC ignorado (0,9 s), órbita do GPS}{UT1$-$UTC ignored (0.9 s), GPS orbit}},
      6/5.62/corTempo/{\tr{UT1$-$UTC ignorado (0,9 s), Equador}{UT1$-$UTC ignored (0.9 s), Equator}},
      5/5.48/corTempo/{\tr{relógio com 1 $\mu$s de erro}{clock with 1 $\mu$s error}},
      4/5.45/corCeu/{\tr{nutação principal ($9{,}2\asec$)}{main nutation ($9.2\asec$)}},
      3/3.97/corFrame/{\tr{movimento polar ($0{,}3\asec$)}{polar motion ($0.3\asec$)}},
      2/2.74/corCeu/{\tr{viés do ICRS ($17{,}9\mas$)}{ICRS bias ($17.9\mas$)}},
      1/2.48/corTempo/{\tr{relógio com 1 ns de erro}{clock with 1 ns error}},
      0/0.78/corCeu/{$dX$, $dY$ ($\num{0.2}\mas$)}} {
      \fill[\cor!80] (0,\yy) rectangle (\val,{\yy+0.62});
      \node[rotulo, anchor=east] at (-0.15,{\yy+0.31}) {\lab};
    }
    \draw[mStoneMid] (0,-0.5) -- (8.2,-0.5);
    \foreach \xx/\lab in {0/{1 mm}, 1/{1 cm}, 2/{10 cm}, 3/{1 m}, 4/{10 m}, 5/{100 m}, 6/{1 km}, 7/{10 km}} {
      \draw[mStoneMid] (\xx,-0.62) -- (\xx,-0.38);
      \draw[mStoneMid!25] (\xx,-0.38) -- (\xx,9.8);
      \node[rotulo, below] at (\xx,-0.65) {\lab};
    }
  \end{tikzpicture}
  \par\smallskip
  {\scriptsize\textcolor{mStoneMid}{\tr{Efeito de cada passo no chão ou na órbita, em escala \emph{logarítmica} (cada marca vale 10 vezes a anterior). Cores: azul $=$ degrau $Q$ (céu); âmbar $=$ $R$ e tempo; verde $=$ $W$ (polo).}{Effect of each step on the ground or at the orbit, on a \emph{logarithmic} scale (each mark is 10 times the previous). Colors: blue $=$ $Q$ step (sky); amber $=$ $R$ and time; green $=$ $W$ (pole).}}}
\end{frame}

\begin{frame}{\tr{O mapa das duas partes}{The map of the two parts}}
  \centering
  \begin{tikzpicture}[caixa/.append style={font=\small, text width=3.4cm, minimum height=2.0cm}]
    \node[caixa, draw=corSist, fill=corSist!8] (ceu) at (0,2.4) {\sist{\tr{O CÉU}{THE SKY}}\\ICRS $\to$ ICRF\\{\footnotesize\tr{quasares, VLBI, Gaia}{quasars, VLBI, Gaia}}\\{\footnotesize\tr{não gira}{does not spin}}};
    \node[caixa, draw=corTempo, fill=corTempo!10, text width=3.2cm] (ponte) at (4.9,2.4) {\tmpo{\tr{A PONTE}{THE BRIDGE}}\\$Q\,R\,W$\\{\footnotesize EOP (IERS)}\\{\footnotesize UT1, UTC, TAI}};
    \node[caixa, draw=corFrame, fill=corFrame!10] (chao) at (9.8,2.4) {\fram{\tr{O CHÃO}{THE GROUND}}\\ITRS $\to$ ITRF\\{\footnotesize SIRGAS2000, WGS 84}\\{\footnotesize\tr{gira com a Terra}{turns with the Earth}}};
    \draw[seta2] (ceu) -- (ponte);
    \draw[seta2] (ponte) -- (chao);
    \node[rotulo, text=mStoneMid] at (0,3.9) {\tr{aulas 1 e 2}{classes 1 and 2}};
    \node[rotulo, text=mStoneMid] at (4.9,3.9) {\tr{aulas 3 e 4}{classes 3 and 4}};
    \node[rotulo, text=mStoneMid] at (9.8,3.9) {\tr{Parte 1}{Part 1}};
    \node[caixa, draw=mStoneMid, fill=mStoneLight, text width=9.0cm, minimum height=0.9cm] (gnss) at (4.9,-0.3) {\tr{\textbf{GNSS}: satélite no céu $\to$ receptor no chão (órbita $\to$ TEME $\to$ ECEF)}{\textbf{GNSS}: satellite in the sky $\to$ receiver on the ground (orbit $\to$ TEME $\to$ ECEF)}};
    \draw[mStoneMid, dashed] (ceu.south) -- ($(ceu.south)+(0,-1.2)$) -- (gnss.west);
    \draw[mStoneMid, dashed] (chao.south) -- ($(chao.south)+(0,-1.2)$) -- (gnss.east);
  \end{tikzpicture}
  \medskip

  \small
  \tr{\bnum{VLBI} liga as duas pontas: antenas no chão $\times$ quasares no céu. Os \bnum{EOP} dizem como a ponte está virada em cada dia.}{\bnum{VLBI} links the two ends: antennas on the ground $\times$ quasars in the sky. The \bnum{EOP} say how the bridge is turned on each day.}
\end{frame}

\begin{frame}{\tr{Para levar para casa --- aula 4 e a Parte 2}{Take-home --- class 4 and Part 2}}
  \footnotesize
  \begin{enumerate}
    \setlength{\itemsep}{2pt}
    \item \tr{\bnum{Vários tempos:} UT1 (rotação, irregular), TAI (átomos, contínuo), $\text{UTC} = \text{TAI} - 37$ s (hora civil), $\text{GPST} = \text{TAI} - 19$ s, $\text{TT} = \text{TAI} + \num{32.184}$ s.}{\bnum{Several times:} UT1 (rotation, irregular), TAI (atoms, continuous), $\text{UTC} = \text{TAI} - 37$ s (civil time), $\text{GPST} = \text{TAI} - 19$ s, $\text{TT} = \text{TAI} + \num{32.184}$ s.}
    \item \tr{\bnum{Segundo intercalar:} 27 desde 1972, o último em 2016. A 28ª CGPM (15/10/2026) adotou o UTC contínuo a partir de 20/5/2027, com limite de 1 h para $|\text{UT1}-\text{UTC}|$ (conferir o texto oficial).}{\bnum{Leap second:} 27 since 1972, the last in 2016. The 28th CGPM (15 Oct 2026) adopted continuous UTC from 20 May 2027, with a 1 h limit for $|\text{UT1}-\text{UTC}|$ (check the official text).}
    \item \tr{\bnum{Relógios do GPS:} $+\num{38.5}\ \mu$s/dia pela relatividade, $\approx$ 11,6 km/dia se não corrigido.}{\bnum{GPS clocks:} $+\num{38.5}\ \mu$s/day from relativity, $\approx$ 11.6 km/day if uncorrected.}
    \item \tr{\bnum{Do satélite ao chão:} órbita (GPS: 2 voltas por dia sideral) $\to$ TEME $\to$ $R(\text{GMST})$ $\to$ ECEF. Com precisão: $Q\,R\,W$, EOP e $\text{UT1}-\text{UTC}$ do IERS, órbitas e relógios do IGS.}{\bnum{From satellite to ground:} orbit (GPS: 2 turns per sidereal day) $\to$ TEME $\to$ $R(\text{GMST})$ $\to$ ECEF. With precision: $Q\,R\,W$, EOP and $\text{UT1}-\text{UTC}$ from the IERS, IGS orbits and clocks.}
    \item \tr{\bnum{Brasil:} ON (hora legal), Fortaleza (VLBI), RBMC/SIRGAS2000 (chão).}{\bnum{Brazil:} ON (legal time), Fortaleza (VLBI), RBMC/SIRGAS2000 (ground).}
  \end{enumerate}
  \begin{ideia}
    \scriptsize \tr{\textbf{A Parte 2 em uma frase:} o céu parado (\sist{ICRS}/\fram{ICRF}, quasares) e o chão que gira (ITRS/ITRF) se ligam pela ponte $Q\,R\,W$, alimentada pelos \bnum{EOP} --- medidos por VLBI e satélites, publicados pelo IERS --- e pelo \bnum{tempo}.}{\textbf{Part 2 in one sentence:} the still sky (\sist{ICRS}/\fram{ICRF}, quasars) and the spinning ground (ITRS/ITRF) are linked by the $Q\,R\,W$ bridge, fed by the \bnum{EOP} --- measured by VLBI and satellites, published by the IERS --- and by \bnum{time}.}
  \end{ideia}
\end{frame}

% ==================================================================
\appendix

\begin{frame}[standout]
  \tr{Apêndice}{Appendix}\\[6pt]
  {\normalsize \tr{Para quem quer ir além}{For those who want to go further}}
\end{frame}

\begin{frame}{\tr{A1 · Marcos normativos}{A1 · Normative milestones}}
  \centering
  \footnotesize
  \renewcommand{\arraystretch}{1.2}
  \setlength{\tabcolsep}{4pt}
  \begin{tabular}{@{}r>{\raggedright\arraybackslash}p{4.2cm}>{\raggedright\arraybackslash}p{7.8cm}@{}}
    \toprule
    & \tr{instância}{body} & \tr{decisão}{decision} \\
    \midrule
    1997 & IAU XXIII (Kyoto), Res.~B2 & \tr{ICRS e ICRF (desde 1/1/1998); Hipparcos como realização óptica}{ICRS and ICRF (from 1 Jan 1998); Hipparcos as the optical realization} \\
    2000 & IAU XXIV (Manchester), Res.~B1.3, B1.6--B1.9 & \tr{BCRS e GCRS; precessão-nutação IAU 2000 (desde 2003); CIP; CIO, TIO e ERA; redefinição do TT}{BCRS and GCRS; IAU 2000 precession-nutation (from 2003); CIP; CIO, TIO and ERA; redefinition of TT} \\
    2006 & \tr{IAU XXVI (Praga), Res.~B1}{IAU XXVI (Prague), Res.~B1} & \tr{precessão P03 (IAU 2006), desde 1/1/2009}{P03 precession (IAU 2006), from 1 Jan 2009} \\
    2009 & IAU XXVII (Rio de Janeiro), Res.~B3 & ICRF2 \tr{(desde 1/1/2010)}{(from 1 Jan 2010)} \\
    2018 & \tr{IAU XXX (Viena), Res.~B2}{IAU XXX (Vienna), Res.~B2} & ICRF3 \tr{(desde 1/1/2019)}{(from 1 Jan 2019)} \\
    2021 & IAU XXXI (Busan), Res.~B3 & Gaia-CRF3 \tr{(desde 1/1/2022)}{(from 1 Jan 2022)} \\
    2022 & CGPM 27, Res.~4 & \tr{ampliar o limite de $|\text{UT1}-\text{UTC}|$ até 2035}{enlarge the $|\text{UT1}-\text{UTC}|$ limit by 2035} \\
    2026 & CGPM 28, Res.~5 & \tr{UTC contínuo desde 20/5/2027; limite de 1 h (conferir o texto oficial)}{continuous UTC from 20 May 2027; 1 h limit (check the official text)} \\
    \bottomrule
  \end{tabular}
\end{frame}

\begin{frame}{\tr{A2 · A cadeia $Q\,R\,W$ em fórmulas}{A2 · The $Q\,R\,W$ chain in formulas}}
  \footnotesize
  \tr{Convenções do IERS (2010, cap.~5). Do ITRS ao GCRS (chão $\to$ céu):}{IERS conventions (2010, ch.~5). From the ITRS to the GCRS (ground $\to$ sky):}
  \[ [\text{GCRS}] = Q(t)\,R(t)\,W(t)\,[\text{ITRS}] \]
  \[ Q(t) = \begin{pmatrix} 1-aX^{2} & -aXY & X \\ -aXY & 1-aY^{2} & Y \\ -X & -Y & 1-a(X^{2}+Y^{2}) \end{pmatrix} \mathbf{R}_3(s), \qquad a = \frac{1}{1+\sqrt{1-X^{2}-Y^{2}}} \]
  \[ R(t) = \mathbf{R}_3(-\text{ERA}), \qquad W(t) = \mathbf{R}_3(-s')\,\mathbf{R}_2(x_p)\,\mathbf{R}_1(y_p) \]
  \begin{itemize}
    \setlength{\itemsep}{1pt}
    \item \tr{$X,\,Y$: coordenadas do CIP no GCRS (modelo IAU 2006/2000A mais $dX$, $dY$); $s$: localizador do CIO, $s \approx -XY/2 +$ termos de $\mu$as;}{$X,\,Y$: coordinates of the CIP in the GCRS (IAU 2006/2000A model plus $dX$, $dY$); $s$: CIO locator, $s \approx -XY/2 +$ terms of $\mu$as;}
    \item \tr{$s'$: localizador do TIO, $s' \approx -47\ \mu\text{as} \times t$ ($t$ em séculos desde J2000.0);}{$s'$: TIO locator, $s' \approx -47\ \mu\text{as} \times t$ ($t$ in centuries since J2000.0);}
    \item \tr{$\mathbf{R}_i(\alpha)$: rotação \emph{dos eixos} de $\alpha$ em torno do eixo $i$ (a convenção do slide ``Girar os eixos'').}{$\mathbf{R}_i(\alpha)$: rotation \emph{of the axes} by $\alpha$ about axis $i$ (the convention of the slide ``Turning the axes'').}
  \end{itemize}
  \tr{Do céu ao chão: $[\text{ITRS}] = W^{\mathsf T} R^{\mathsf T} Q^{\mathsf T}\,[\text{GCRS}]$.}{From sky to ground: $[\text{ITRS}] = W^{\mathsf T} R^{\mathsf T} Q^{\mathsf T}\,[\text{GCRS}]$.}
\end{frame}

\begin{frame}{\tr{A3 · ERA, GMST e a hora sideral}{A3 · ERA, GMST and sidereal time}}
  \footnotesize
  \[ \text{ERA} = 2\pi\,(\num{0.7790572732640} + \num{1.00273781191135448}\,T_u), \qquad T_u = \text{JD(UT1)} - \num{2451545.0} \]
  \[ \text{GMST} = \text{ERA} + \num{0.014506}\asec + \num{4612.156534}\asec\,t + \num{1.3915817}\asec\,t^{2} - \num{0.00000044}\asec\,t^{3} - \ldots \]
  \tr{com $t$ em séculos de TT desde J2000.0. A diferença $\text{GMST}-\text{ERA}$ é a precessão em ascensão reta (equinócio contra CIO).}{with $t$ in centuries of TT since J2000.0. The difference $\text{GMST}-\text{ERA}$ is the precession in right ascension (equinox versus CIO).}
  \medskip

  \renewcommand{\arraystretch}{1.3}
  \begin{tabular}{@{}lr@{}}
    \toprule
    \tr{em 18/10/2026, 0 h UT1}{on 18 Oct 2026, 0 h UT1} & \\
    \midrule
    ERA & $\num{26.15528}\gr$ \\
    GMST & $\num{26.49858}\gr$ \\
    $\text{GMST}-\text{ERA}$ & $\num{1235.9}\asec$ \\
    \tr{GMST pela fórmula do simulador (Meeus)}{GMST by the simulator's formula (Meeus)} & $\num{26.49860}\gr$ \\
    \tr{dia sideral pelo ERA (estrela fixa)}{sidereal day by the ERA (fixed star)} & $\num{86164.0989}$ s \\
    \tr{dia sideral pelo GMST (equinócio)}{sidereal day by the GMST (equinox)} & $\num{86164.0905}$ s \\
    \bottomrule
  \end{tabular}
\end{frame}

\begin{frame}{\tr{A4 · Precessão: das coordenadas J2000 às de hoje}{A4 · Precession: from J2000 coordinates to today's}}
  \footnotesize
  \tr{Com $T = (\text{JD(TT)} - \num{2451545.0})/\num{36525}$ e os ângulos de Lieske (IAU 1976):}{With $T = (\text{JD(TT)} - \num{2451545.0})/\num{36525}$ and the Lieske angles (IAU 1976):}
  \[ \begin{aligned}
     \zeta &= \num{2306.2181}\asec\,T + \num{0.30188}\asec\,T^{2} + \num{0.017998}\asec\,T^{3} \\
     z &= \num{2306.2181}\asec\,T + \num{1.09468}\asec\,T^{2} + \num{0.018203}\asec\,T^{3} \\
     \theta &= \num{2004.3109}\asec\,T - \num{0.42665}\asec\,T^{2} - \num{0.041833}\asec\,T^{3}
     \end{aligned} \]
  \[ \begin{aligned}
     A &= \cos\delta_0 \sin(\alpha_0+\zeta), \quad B = \cos\theta \cos\delta_0 \cos(\alpha_0+\zeta) - \sin\theta \sin\delta_0, \\
     C &= \sin\theta \cos\delta_0 \cos(\alpha_0+\zeta) + \cos\theta \sin\delta_0, \quad \alpha = \operatorname{atan2}(A,B) + z, \quad \delta = \arcsin C
     \end{aligned} \]
  \begin{cuidado}[title={\tr{Exemplo: Sirius, de J2000.0 a 18/10/2026}{Example: Sirius, from J2000.0 to 18 Oct 2026}}]
    \tr{$(\alpha,\delta)$: $(\num{101.2872}\gr,\ \num{-16.7161}\gr) \to (\num{101.5865}\gr,\ \num{-16.7457}\gr)$ só pela precessão: $\Delta\alpha = +1078\asec$, $\Delta\delta = -106\asec$, separação $17{,}3\amin$. O movimento próprio ($\num{1.339}\asec$/ano) soma ainda $-32{,}8\asec$ em $\delta$.}{$(\alpha,\delta)$: $(\num{101.2872}\gr,\ \num{-16.7161}\gr) \to (\num{101.5865}\gr,\ \num{-16.7457}\gr)$ by precession alone: $\Delta\alpha = +1078\asec$, $\Delta\delta = -106\asec$, separation $17.3\amin$. The proper motion ($\num{1.339}\asec$/yr) adds a further $-32.8\asec$ in $\delta$.}
  \end{cuidado}
\end{frame}

\begin{frame}{\tr{A5 · Relatividade nos relógios do GNSS: as contas}{A5 · Relativity in GNSS clocks: the sums}}
  \footnotesize
  \tr{Para órbita circular ($v^{2} = GM/a$), relógio de bordo contra relógio no geoide ($R$):}{For a circular orbit ($v^{2} = GM/a$), on-board clock against a clock on the geoid ($R$):}
  \[ \frac{\Delta f}{f} = \underbrace{\frac{GM}{c^{2}}\Bigl(\frac{1}{R}-\frac{1}{a}\Bigr)}_{\text{\tr{gravidade}{gravity}}\ =\ +\num{5.28e-10}} \;-\; \underbrace{\frac{v^{2}}{2c^{2}}}_{\text{\tr{velocidade}{velocity}}\ =\ -\num{8.35e-11}} \;=\; \frac{GM}{c^{2}}\Bigl(\frac{1}{R}-\frac{3}{2a}\Bigr) \approx \num{4.45e-10} \]
  \tr{Em um dia ($\num{86400}$ s): $+\num{45.7}\ \mu$s $-$ $\num{7.2}\ \mu$s $=$ $+\num{38.5}\ \mu$s. Com $W_0$ e a rotação da Terra, $\Delta f/f = \num{4.4645e-10}$ ($\num{38.6}\ \mu$s/dia).}{In one day ($\num{86400}$ s): $+\num{45.7}\ \mu$s $-$ $\num{7.2}\ \mu$s $=$ $+\num{38.5}\ \mu$s. With $W_0$ and the Earth's rotation, $\Delta f/f = \num{4.4645e-10}$ ($\num{38.6}\ \mu$s/day).}
  \medskip

  \begin{itemize}
    \setlength{\itemsep}{1pt}
    \item \tr{Frequência de projeto: $\num{10.23}\ \text{MHz}\,(1-\num{4.4645e-10}) = \num{10.22999999543}$ MHz.}{Design frequency: $\num{10.23}\ \text{MHz}\,(1-\num{4.4645e-10}) = \num{10.22999999543}$ MHz.}
    \item \tr{Sem a correção: $c \times \num{38.6}\ \mu\text{s} \approx 11{,}6$ km por dia.}{Without the correction: $c \times \num{38.6}\ \mu\text{s} \approx 11.6$ km per day.}
    \item \tr{Órbita excêntrica: termo periódico $\Delta t_r = -\dfrac{2\sqrt{GM\,a}}{c^{2}}\,e\sin E$, de uns 7 ns para $e = 0{,}01$ (corrigido pelo receptor).}{Eccentric orbit: periodic term $\Delta t_r = -\dfrac{2\sqrt{GM\,a}}{c^{2}}\,e\sin E$, about 7 ns for $e = 0.01$ (corrected by the receiver).}
    \item \tr{No chão: $g/c^{2} = \num{1.0911e-16}$ por metro, ou $\num{1.09e-18}$ por centímetro.}{On the ground: $g/c^{2} = \num{1.0911e-16}$ per meter, or $\num{1.09e-18}$ per centimeter.}
  \end{itemize}
\end{frame}

\begin{frame}{\tr{A6 · Os 27 segundos intercalares}{A6 · The 27 leap seconds}}
  \footnotesize
  \begin{columns}[T]
    \begin{column}{0.48\textwidth}
      \renewcommand{\arraystretch}{1.3}
      \begin{tabular}{@{}>{\raggedright\arraybackslash}p{2.6cm}>{\raggedright\arraybackslash}p{2.1cm}r@{}}
        \toprule
        \tr{ano}{year} & \tr{data}{date} & \tr{quantos}{how many} \\
        \midrule
        1972 & \tr{30/6 e 31/12}{30 Jun, 31 Dec} & 2 \\
        1973--1979 & \tr{31/12}{31 Dec} & 7 \\
        1981--1983, 1985 & \tr{30/6}{30 Jun} & 4 \\
        1987, 1989, 1990 & \tr{31/12}{31 Dec} & 3 \\
        1992--1994 & \tr{30/6}{30 Jun} & 3 \\
        1995 & \tr{31/12}{31 Dec} & 1 \\
        \bottomrule
      \end{tabular}
    \end{column}
    \begin{column}{0.48\textwidth}
      \renewcommand{\arraystretch}{1.3}
      \begin{tabular}{@{}>{\raggedright\arraybackslash}p{2.6cm}>{\raggedright\arraybackslash}p{2.1cm}r@{}}
        \toprule
        \tr{ano}{year} & \tr{data}{date} & \tr{quantos}{how many} \\
        \midrule
        1997 & \tr{30/6}{30 Jun} & 1 \\
        1998 & \tr{31/12}{31 Dec} & 1 \\
        2005, 2008 & \tr{31/12}{31 Dec} & 2 \\
        2012, 2015 & \tr{30/6}{30 Jun} & 2 \\
        2016 & \tr{31/12}{31 Dec} & 1 \\
        \midrule
        \textbf{\tr{total}{total}} & & \textbf{27} \\
        \bottomrule
      \end{tabular}
    \end{column}
  \end{columns}
  \medskip

  \tr{Cada inserção acontece às \texttt{23:59:60} UTC da data indicada. $\text{TAI}-\text{UTC}$ passou de 10 s (1972) a 37 s (a partir de 1/1/2017).}{Each insertion happens at \texttt{23:59:60} UTC on the date shown. $\text{TAI}-\text{UTC}$ went from 10 s (1972) to 37 s (from 1 Jan 2017).}
\end{frame}

\begin{frame}{\tr{A7 · Glossário (1/3): sistemas e referenciais}{A7 · Glossary (1/3): systems and frames}}
  \centering
  \scriptsize
  \renewcommand{\arraystretch}{1.3}
  \setlength{\tabcolsep}{4pt}
  \begin{tabular}{@{}>{\raggedright\arraybackslash}p{1.8cm}>{\raggedright\arraybackslash}p{5.4cm}>{\raggedright\arraybackslash}p{5.4cm}@{}}
    \toprule
    & \tr{português}{Portuguese} & English \\
    \midrule
    ICRS / ICRF & Sistema / Referencial Celeste Internacional & International Celestial Reference System / Frame \\
    GCRS & Sistema de Referência Celeste Geocêntrico & Geocentric Celestial Reference System \\
    ITRS / ITRF & Sistema / Referencial Terrestre Internacional & International Terrestrial Reference System / Frame \\
    CIP & Polo Celeste Intermediário & Celestial Intermediate Pole \\
    CIO / TIO & Origem Intermediária Celeste / Terrestre & Celestial / Terrestrial Intermediate Origin \\
    ERA & Ângulo de Rotação da Terra & Earth Rotation Angle \\
    GMST & Tempo Sideral Médio de Greenwich & Greenwich Mean Sidereal Time \\
    TEME & Equador Verdadeiro, Equinócio Médio (SGP4) & True Equator, Mean Equinox (SGP4's frame) \\
    \bottomrule
  \end{tabular}
\end{frame}

\begin{frame}{\tr{A8 · Glossário (2/3): movimentos e parâmetros}{A8 · Glossary (2/3): motions and parameters}}
  \centering
  \scriptsize
  \renewcommand{\arraystretch}{1.3}
  \setlength{\tabcolsep}{4pt}
  \begin{tabular}{@{}>{\raggedright\arraybackslash}p{1.8cm}>{\raggedright\arraybackslash}p{5.4cm}>{\raggedright\arraybackslash}p{5.4cm}@{}}
    \toprule
    & \tr{português}{Portuguese} & English \\
    \midrule
    EOP & Parâmetros de Orientação da Terra & Earth Orientation Parameters \\
    LOD & excesso da duração do dia sobre 86\,400 s & length of day (excess over 86\,400 s) \\
    \tr{mov.\ polar}{polar motion} & movimento polar ($x_p$, $y_p$) & polar motion ($x_p$, $y_p$) \\
    \tr{precessão, nutação}{precession, nutation} & precessão, nutação & precession, nutation \\
    \tr{baricentro}{barycenter} & centro de massa do Sistema Solar & center of mass of the Solar System \\
    \tr{linha de base}{baseline} & distância entre duas antenas & distance between two antennas \\
    quasar & ``quase uma estrela'': núcleo ativo de galáxia & quasi-stellar object: active galactic nucleus \\
    mas, $\mu$as & milissegundo e microssegundo de arco & milliarcsecond and microarcsecond \\
    \bottomrule
  \end{tabular}
\end{frame}

\begin{frame}{\tr{A9 · Glossário (3/3): tempo e serviços}{A9 · Glossary (3/3): time and services}}
  \centering
  \scriptsize
  \renewcommand{\arraystretch}{1.3}
  \setlength{\tabcolsep}{4pt}
  \begin{tabular}{@{}>{\raggedright\arraybackslash}p{1.8cm}>{\raggedright\arraybackslash}p{5.4cm}>{\raggedright\arraybackslash}p{5.4cm}@{}}
    \toprule
    & \tr{português}{Portuguese} & English \\
    \midrule
    UT1, UTC, TAI & Tempo Universal (rotação); Coordenado; Atômico Internacional & Universal Time (rotation); Coordinated; International Atomic Time \\
    TT & Tempo Terrestre & Terrestrial Time \\
    GPST, GST, BDT & escalas de tempo do GPS, Galileo e BeiDou & GPS, Galileo and BeiDou system times \\
    \tr{seg.\ intercalar}{leap second} & segundo inserido no UTC & second inserted into UTC \\
    VLBI & Interferometria de Linha de Base Muito Longa & Very Long Baseline Interferometry \\
    SLR, DORIS & telemetria laser a satélites; Doppler de balizas em terra & Satellite Laser Ranging; Doppler beacons (DORIS) \\
    IERS, IVS, IGS & serviços: rotação da Terra; VLBI; GNSS & services: Earth rotation; VLBI; GNSS \\
    BIPM, CGPM & Bureau Internacional de Pesos e Medidas; Conferência Geral & Intl.\ Bureau of Weights and Measures; General Conference \\
    \bottomrule
  \end{tabular}
\end{frame}

\begin{frame}{\tr{Fontes e leituras (1/2)}{Sources and further reading (1/2)}}
  \scriptsize
  \begin{itemize}
    \setlength{\itemsep}{1.5pt}
    \item Petit, G., Luzum, B. (eds.), \emph{IERS Conventions (2010)}, IERS Technical Note 36, Frankfurt am Main: Verlag des BKG, 2010.
    \item IAU, \tr{Resoluções}{Resolutions}: 1997 B2; 2000 B1.3, B1.6--B1.9; 2006 B1; 2009 B3; 2018 B2; 2021 B3.
    \item Ma, C. et al., \emph{AJ} 116, 516 (1998) [ICRF]. Fey, A. L. et al., \emph{AJ} 150, 58 (2015) [ICRF2].
    \item Charlot, P. et al., \emph{A\&A} 644, A159 (2020) [ICRF3].
    \item Gaia Collaboration, Klioner, S. A. et al., \emph{A\&A} 667, A148 (2022) [Gaia-CRF3]. ESA, \emph{The Hipparcos and Tycho Catalogues}, SP-1200 (1997).
    \item Capitaine, N., Wallace, P. T., \emph{A\&A} 450, 855 (2006) [\tr{cadeia baseada no CIO}{CIO-based chain}]. Meeus, J., \emph{Astronomical Algorithms}, 2nd ed., Willmann-Bell (1998).
    \item Stephenson, F. R., Morrison, L. V., Hohenkerk, C. Y., \emph{Proc. R. Soc. A} 472, 20160404 (2016) [\tr{rotação da Terra, 720 a.C.--2015}{Earth's rotation, 720 BC--AD 2015}].
    \item Agnew, D. C., \emph{Nature} 628, 333--336 (2024) [\tr{segundo intercalar e o clima}{leap second and climate}].
    \item BIPM, \tr{\emph{Le Système international d'unités (Brochura do SI)}, 9ª ed. (2019); CGPM 2022, Res.~4; CGPM 2026, Res.~5 (rascunho de resolução e declaração do presidente do CCTF; o texto oficial deve ser conferido)}{\emph{The International System of Units (SI Brochure)}, 9th ed. (2019); CGPM 2022, Res.~4; CGPM 2026, Res.~5 (draft resolution and statement of the CCTF president; the official text must be checked)}.
  \end{itemize}
\end{frame}

\begin{frame}{\tr{Fontes e leituras (2/2)}{Sources and further reading (2/2)}}
  \scriptsize
  \begin{itemize}
    \setlength{\itemsep}{1.5pt}
    \item Ashby, N., \emph{Living Rev. Relativ.} 6, 1 (2003) [\tr{relatividade no GPS}{relativity in GPS}]. Takamoto, M. et al., \emph{Nature Photonics} 14, 411--415 (2020) [\tr{relógios ópticos na Tokyo Skytree}{optical clocks at Tokyo Skytree}].
    \item Vallado, D. A., Crawford, P., Hujsak, R., Kelso, T. S., \emph{Revisiting Spacetrack Report \#3}, AIAA 2006-6753 [SGP4, TEME].
    \item ESA, \emph{Navipedia}: ``Time References in GNSS'' \tr{e}{and} ``Atomic Time''.
    \item IVS, \emph{Biennial Report 2017+2018}: \tr{relatório da estação de Fortaleza}{Fortaleza station report}. \tr{Universidade Presbiteriana Mackenzie, comunicado sobre a nova antena do Rádio Observatório Espacial do Nordeste (nov./2025).}{Mackenzie Presbyterian University, press release on the new antenna of the Northeast Space Radio Observatory (Nov 2025).}
    \item \tr{Observatório Nacional, Hora Legal Brasileira (Decreto 4.264/2002). IBGE, RBMC (153 estações ativas em dez./2025).}{National Observatory, Brazilian Legal Time (Decree 4,264/2002). IBGE, RBMC (153 active stations in Dec 2025).}
    \item \tr{Neste repositório: \texttt{sgr\_terrestres} (Parte 1), \texttt{quase\_geoide}, \texttt{gravimetria}; simuladores \texttt{astrometria} e \texttt{GNSS\_simul}; números e figuras: \texttt{sgr\_celestes\_numeros.py}.}{In this repository: \texttt{sgr\_terrestres} (Part 1), \texttt{quase\_geoide}, \texttt{gravimetria}; simulators \texttt{astrometria} and \texttt{GNSS\_simul}; numbers and figures: \texttt{sgr\_celestes\_numeros.py}.}
  \end{itemize}
  \medskip

  \begin{ideia}
    \footnotesize\tr{Valores e datas conferidos em 18/10/2026. Os dados do Brasil e as notícias da 28ª CGPM vieram de resumos de busca, porque alguns sites oficiais ficaram inacessíveis; confira nas fontes primárias antes de citar.}{Values and dates checked on 18 Oct 2026. The Brazilian data and the news of the 28th CGPM came from search summaries, because some official sites were unreachable; check the primary sources before quoting.}
  \end{ideia}
\end{frame}

\end{document}
